% 翻译检查点：源 /tmp/sga3-audit/units/XXIV.tex；保留原编号与引用键。
\markboth{第 XXIV 讲. 约化群的自同构}{第 XXIV 讲. 约化群的自同构}
\setcounter{footnote}{0}
\SGARefTarget{sga3:xxiv:chapter:xxiv}\chapter*{第 XXIV 讲\\约化群的自同构}
\addcontentsline{toc}{chapter}{第 XXIV 讲. 约化群的自同构}
\markboth{第 XXIV 讲. 约化群的自同构}{第 XXIV 讲. 约化群的自同构}
\begin{center}
  \textit{Michel Demazure 著}
\end{center}
\begingroup
\renewcommand{\thefootnote}{0}
\footnote{编者注：2024 年 10 月 13 日版本。}
\addtocounter{footnote}{-1}
\endgroup
本讲的第一部分（第 1 至 5 节）直接由约化群（reductive group）有“足够多的外自同构”这一事实推出；这一事实本身又是钉扎群（pinned group）同构定理的最弱形式的推论。第二部分（第 6、7 节）阐述上一讲中更精确结果的两个应用；特别地，第 7 节使用了生成元与关系定理的显式形式。最后，我们在附录（第 8 节）中收入正文用到的“Galois”上同调结果。

先说明上同调记号。若 \(S\) 是概形，\(G\) 是 \(S\)-群概形，则以 \(H^1(S,G)\) 表示按 \textup{(fpqc)} 拓扑计算的 \(S\) 以 \(G\) 为系数的第一上同调集；它也就是 \(G\) 下主齐性 \textup{(fpqc)} 层的同构类集合。以 \(H^1_{\mathrm{\acute et}}(S,G)\) 表示平展拓扑下的相应集合；因此它是 \(H^1(S,G)\) 中由 \(G\) 下拟等平凡的齐性层（即平展拓扑下局部平凡的齐性层）的类组成的部分。以 \(\Fib(S,G)\) 表示 \(H^1(S,G)\) 中由可表层（主齐性丛）的类组成的部分。于是有包含关系
\[
  H^1_{\mathrm{\acute et}}(S,G)\subset H^1(S,G),
  \qquad
  \Fib(S,G)\subset H^1(S,G).
\]

若 \(G\) 下的每个主齐性层都可表（例如 \(G\) 在 \(S\) 上拟仿射时；参见 SGA~1，VIII，\SGARefLink{sga3:viii:proposition:7:9}{7.9}），则
\[
  \Fib(S,G)=H^1(S,G).
\]

若 \(S'\to S\) 是 \textup{(fpqc)} 拓扑的覆盖态射，则以 \(H^1(S'/S,G)\) 表示典范映射
\[
  H^1(S,G)\longrightarrow H^1(S',G_{S'}).
\]
的核。已知 \(H^1(S'/S,G)\) 可以用单纯方法计算（TDTE~I，\S A.4）；这意味着，当 \(S'\to S\) 是平展拓扑的覆盖时，\(H^1(S'/S,G)\) 也是映射
\[
  H^1_{\mathrm{\acute et}}(S,G)
  \longrightarrow
  H^1_{\mathrm{\acute et}}(S',G_{S'}).
\]
的核。

最后，按照\SGARefLink{sga3:viii:corollary:4:5}{第 VIII 讲 4.5}，把以下断言称为“\SGARefLink{sga3:viii:corollary:4:5}{定理 90}”：“\(G_{m,S}\) 下的每个主齐性层都可表且局部平凡。”这一断言等价于
\[
  H^1(S,G_{m,S})=\operatorname{Pic}(S),
\]
也等价于：当 \(S\) 局部（或更一般地，半局部）时，
\[
  H^1(S,G_{m,S})=0
\]

\SGARefTarget{sga3:xxiv:section:1}
\section*{1. 约化群的自同构概形}
\addcontentsline{toc}{section}{1. 约化群的自同构概形}
\SGARefTarget{sga3:xxiv:1.0}
\subsection*{1.0.}
首先须使上一讲中的某些定义精确化。设 \(S\) 是非空概形，\(G\) 是约化 \(S\)-群，并设
\[
  \mathscr R=(M,M^*,R,R^*,\Delta)
\]
是钉扎约化根数据（\SGARefLink{sga3:XXIII:1.5}{第 XXIII 讲 1.5}）。\emph{\(G\) 的类型为 \(\mathscr R\) 的钉扎}，或称 \emph{\(G\) 的 \(\mathscr R\)-钉扎}，由以下资料组成：
\begin{enumerate}
\item[(i)] 从 \(D_S(M)\) 到 \(G\) 的一个极大环面 \(T\) 的同构（等价地，一个单态射 \(D_S(M)\to G\)，其像为 \(G\) 的极大环面 \(T\)），它将 \(R\) 等同于 \(G\) 关于 \(T\) 的根系（\SGARefLink{sga3:xix:definition:3-6}{第 XIX 讲 3.6}），将 \(R^*\) 等同于相应的余根集；
\item[(ii)] 对每个 \(\alpha\in\Delta\)，一个元素
  \[
    X_\alpha\in\Gamma(S,\mathfrak g^\alpha)^\times.
  \]
\end{enumerate}
群 \(G\) 有一个 \(\mathscr R\)-钉扎，当且仅当它分裂且类型为 \(\mathscr R\)（\SGARefLink{sga3:xxii:def:2.7}{第 XXII 讲 2.7}）。

若 \(u:G\to G'\) 是约化 \(S\)-群的同构，则 \(G\) 的每个 \(\mathscr R\)-钉扎 \(\mathcal E\) 都通过“结构转移”确定 \(G'\) 的一个 \(\mathscr R\)-钉扎 \(u(\mathcal E)\)。若 \(v:\mathscr R'\to\mathscr R\) 是钉扎根数据的同构，则 \(G'\) 的每个 \(\mathscr R'\)-钉扎 \(\mathcal E'\) 都通过结构转移确定 \(G\) 的一个 \(\mathscr R\)-钉扎 \(v(\mathcal E')\)。

称三元组 \((G,\mathscr R,\mathcal E)\) 为\emph{钉扎群}，其中 \(G\) 是约化 \(S\)-群，\(\mathscr R\) 是钉扎约化根数据，\(\mathcal E\) 是 \(G\) 的类型为 \(\mathscr R\) 的钉扎。从钉扎群 \((G,\mathscr R,\mathcal E)\) 到钉扎群 \((G',\mathscr R',\mathcal E')\) 的同构是一个有序对 \((u,v)\)，其中 \(u:G\to G'\) 是同构，\(v:\mathscr R'\to\mathscr R\) 是钉扎根数据的同构，并且
\[
  u(\mathcal E)=v(\mathcal E').
\]
\footnote{编者注：因此，对每个 \(\alpha\in R'\)，有 \(\Lie(u)(X_{v(\alpha)})=X'_\alpha\)。}

\textit{注意。} 若 \(S\) 非空，则 \(v\) 由 \(u\) 唯一确定；因而也把 \(u\) 称为钉扎群的同构，这是术语上的滥用。特别地，若 \((G,\mathscr R,\mathcal E)\) 是钉扎群，则 \((G,\mathscr R,\mathcal E)\) 的自同构就是 \(G\) 的自同构 \(u\)，使得存在 \(\mathscr R\) 的自同构 \(v\)，满足 \(u(\mathcal E)=v(\mathcal E)\)。因此它是 \(G\) 的如下自同构：正规化 \(T\)，在 \(T\) 上诱导形如 \(D_S(h)\) 的自同构，其中 \(h\) 是 \(M\) 的自同构，并置换元素 \(X_\alpha\)，\(\alpha\in\Delta\)。（容易看出，上述条件实际上刻画了 \((G,\mathscr R,\mathcal E)\) 的自同构。）

显然有一个反变函子
\[
  \Phi:\quad
  (G,\mathscr R,\mathcal E)\longmapsto\mathscr R,
  \qquad
  (u,v)\longmapsto v,
\]
而上一讲的主要结果（\SGARefLink{sga3:xxiii:theorem:4-1}{第 XXIII 讲 4.1}）表明，这个函子完全忠实（而且我们将在下一讲看到，它是范畴等价）。特别地，钉扎群 \((G,\mathscr R,\mathcal E)\) 的自同构群典范同构于钉扎根数据 \(\mathscr R\) 的自同构群（参见\SGARefLink{sga3:XXIII:5.5}{第 XXIII 讲 5.5}）。

\SGARefTarget{sga3:xxiv:1.1}
\subsection*{1.1.}
设 \(S\) 是概形；赋予 \((\mathbf{Sch})/S\) 以 \textup{(fpqc)} 拓扑，并考虑 \(S\)-群层 \(\underline{\Aut}_{S\text{-}\mathrm{gr.}}(G)\)，其中 \(G\) 是 \(S\)-群概形。有一个 \(S\)-群层的正合序列
\[
  1\longrightarrow \underline{\Centr}(G)
  \longrightarrow G
  \xrightarrow{\mathrm{int}}
  \underline{\Aut}_{S\text{-}\mathrm{gr.}}(G),
\]
它定义一个单态射
\[
  j:G/\underline{\Centr}(G)
  \longrightarrow \underline{\Aut}_{S\text{-}\mathrm{gr.}}(G).
\]
\(j\) 的像层就是 \(G\) 的内自同构层。\(G\) 的自同构 \(u\) 是内自同构，当且仅当存在覆盖族 \(\{S_i\to S\}\)，并对每个 \(i\) 存在元素 \(g_i\in G(S_i)\)，使得
\[
  \operatorname{int}(g_i)=u_{S_i}.
\]
此时，若 \(v\) 是 \(G\) 的另一个自同构，则立即可见
\[
  \operatorname{int}(v)u=vuv^{-1}
\]
是由族 \(g'_i=v(g_i)\) 定义的内自同构。因此 \(j\) 的像在 \(\underline{\Aut}_{S\text{-}\mathrm{gr.}}(G)\) 中正规。记商群层为 \(\underline{\Autext}(G)\)，称为 \(G\) 的外自同构层。于是有正合序列
\[
  1\longrightarrow G/\underline{\Centr}(G)
  \longrightarrow \underline{\Aut}_{S\text{-}\mathrm{gr.}}(G)
  \longrightarrow \underline{\Autext}(G)
  \longrightarrow 1.
\]
上述定义均与基变换相容。当然，它们在任意景中都成立。

\SGARefTarget{sga3:xxiv:1.2}
\subsection*{1.2.}
设 \(S\) 是概形，\((G,\mathscr R,\mathcal E)\) 是钉扎约化 \(S\)-群。设 \(E\) 为钉扎根数据 \(\mathscr R\) 的（抽象）自同构群，即 \(\mathscr R\) 的正规化 \(\Delta\) 的自同构组成的群。由\SGARefLink{sga3:XXIII:5.5}{第 XXIII 讲 5.5}，有典范单态射
\[
  E\longrightarrow \underline{\Aut}_{S\text{-}\mathrm{gr.}}(G)
\]
它把 \(h\in E\) 对应到钉扎群 \(G\) 的唯一自同构 \(u\)，使得 \(\Phi(u)=h\)。这个单态射典范地定义层的单态射
\[
  \tag{*}\label{sga3:xxiv:subsection:1-2:equation:star}
  a:E_S\longrightarrow \underline{\Aut}_{S\text{-}\mathrm{gr.}}(G).
\]
\(G\) 的自同构 \(u\) 是 \(a\) 的像层的截面，当且仅当下列条件成立：
\begin{enumerate}
\item[(i)] \(u\) 正规化 \(T\)。此时已知 \(u\) 置换 \(G\) 关于 \(T\) 的根。若 \(\alpha\in R\)，则
  \[
    u(\alpha):t\longmapsto\alpha\bigl(u^{-1}(t)\bigr)
  \]
  在 \(S\) 上局部地具有形式 \(t\longmapsto\beta(t)\)，其中 \(\beta\in R\)。第二个条件于是可写成：
\item[(ii)] 若 \(\alpha\in\Delta\)，\(U\) 是 \(S\) 的开子集，并且 \(u(\alpha)_U\in R\)，则 \(u(\alpha)_U\in\Delta\)，且
  \[
    \Lie(u_U)(X_\alpha)_U=(X_{u(\alpha)})_U.
  \]
\end{enumerate}
由定义立即得知，\(a(E_S)\) 的截面正规化由钉扎确定的 \(G\) 的子群
\[
  T,\quad B,\quad B^{-},\quad U,\quad U^{-}.
\]

作好这些定义后，有如下结果。

% 已译源行 1--261。
\SGARefTarget{sga3:xxiv:1.3}
\paragraph{\SGARefLink{sga3:xxiv:1.3}{定理 1.3}.}
\emph{设 \(S\) 是概形，\(G\) 是约化 \(S\)-群。考虑 \(S\)-群层的典范正合序列}
\footnote{编者注：回忆 \(\operatorname{ad}(G)=G/\underline{\Centr}(G)\) 表示 \(G\) 的伴随群；参见\SGARefLink{sga3:xxii:def:4.3.6}{第 XXII 讲 4.3.6}。}
\[
  1\longrightarrow \operatorname{ad}(G)
  \longrightarrow \underline{\Aut}_{S\text{-}\mathrm{gr.}}(G)
  \xrightarrow{p}\underline{\Autext}(G)
  \longrightarrow 1.
\]
\emph{则：}
\begin{enumerate}
\item[\SGAInlineRefTarget{sga3:xxiv:1.3:item:i}(i)] \(\underline{\Aut}_{S\text{-}\mathrm{gr.}}(G)\) 可由一个光滑分离的 \(S\)-概形表示；
\item[\SGAInlineRefTarget{sga3:xxiv:1.3:item:ii}(ii)] \(\underline{\Autext}(G)\) 可由一个有限生成的扭常 \(S\)-概形表示（\SGARefLink{sga3:x:definition:5-1}{第 X 讲 5.1}）；
\item[\SGAInlineRefTarget{sga3:xxiv:1.3:item:iii}(iii)] 若 \(G\) 分裂，则上述正合序列分裂。更精确地，对 \(G\) 的每个钉扎，态射（参见 \SGARefLink{sga3:xxiv:1.2}{1.2}~\((*)\)）
  \[
    p\circ a:E_S\longrightarrow\underline{\Autext}(G)
  \]
  是同构。
\end{enumerate}

先说明如何由以下引理推出定理。
\SGARefTarget{sga3:xxiv:1.4}
\paragraph{\SGARefLink{sga3:xxiv:1.4}{引理 1.4}.}
\emph{在 \textup{\SGARefLink{sga3:xxiv:1.3:item:iii}{(iii)}} 的假设下，\(\underline{\Aut}_{S\text{-}\mathrm{gr.}}(G)\) 是半直积}
\[
  a(E_S)\cdot\operatorname{ad}(G).
\]
\footnote{编者注：这里及下文中，用记号 \(\operatorname{ad}(G)\) 代替 \(\operatorname{int}(G)\)。}

当 \(G\) 分裂时，引理立即推出定理。由于 \(G\) 在平展拓扑下局部分裂（\SGARefLink{sga3:xxii:cor:2.3}{第 XXII 讲 2.3}），因此在 \textup{(fppf)} 拓扑下也局部分裂，而后者对于扭常态射的纤维范畴是“有效下降的”（\SGARefLink{sga3:x:corollary:5-5}{第 X 讲 5.5}），故一般情形的断言 \textup{\SGARefLink{sga3:xxiv:1.3:item:ii}{(ii)}} 得证（参见\SGARefLink{sga3:iv:proposition:4:6:8}{第 IV 讲 4.6.8}）。为推出 \textup{\SGARefLink{sga3:xxiv:1.3:item:i}{(i)}}，注意到 \(\operatorname{ad}(G)\) 在 \(S\) 上仿射；因此只要 \(\underline{\Aut}_{S\text{-}\mathrm{gr.}}(G)\) 可表，\(p\) 就是仿射的，结果遂由仿射概形的下降得到。\footnote{编者注：参见 SGA~1，VIII，\SGARefLink{sga3:viii:proposition:2:1}{2.1}。}

只须证明 \SGARefLink{sga3:xxiv:1.4}{1.4}。为此，只须证明：
\SGARefTarget{sga3:xxiv:1.5}
\paragraph{\SGARefLink{sga3:xxiv:1.5}{引理 1.5}.}
\emph{若 \((\mathscr R,\mathcal E)\) 和 \((\mathscr R',\mathcal E')\) 是约化 \(S\)-群 \(G\) 的两个钉扎，则存在唯一的 \(G\) 在 \(S\) 上的内自同构 \(u\)，将一个钉扎变为另一个；即存在同构
\[
  v:\mathscr R'\xrightarrow{\sim}\mathscr R
\]
使得 \(u(\mathcal E)=v(\mathcal E')\)；参见 \SGARefLink{sga3:xxiv:1.0}{1.0}。}

\SGARefTarget{sga3:xxiv:1.5.1}
\subsubsection*{1.5.1. 唯一性。}
只须证明，若 \(G\) 是钉扎 \(S\)-群，\(\operatorname{int}(g)\) 是钉扎群的自同构（\(g\in G(S)\)），则 \(\operatorname{int}(g)=\operatorname{id}\)。首先，
\[
  \operatorname{int}(g)T=T,\qquad
  \operatorname{int}(g)B=B,
\]
因而
\[
  g\in
  \underline{\operatorname{Norm}}_G(T)(S)
  \cap\underline{\operatorname{Norm}}_G(B)(S)
  =T(S)
\]
（例如参见\SGARefLink{sga3:xxii:prop:5.6.1}{第 XXII 讲 5.6.1}）。因此 \(\operatorname{int}(g)\) 正规化每个 \(U_\alpha\)，并且对每个 \(\alpha\in\Delta\)，有
\[
  \Lie(\operatorname{int}(g))X_\alpha
  =\operatorname{Ad}(g)X_\alpha
  =\alpha(g)X_\alpha
\]
于是对 \(\alpha\in\Delta\) 有 \(\alpha(g)=1\)，所以
\[
  g\in
  \bigcap_{\alpha\in\Delta}(\ker\alpha)(S)
  =\underline{\Centr}(G)(S)
\]
（\SGARefLink{sga3:xxii:rem:4.1.8}{第 XXII 讲 4.1.8}）。\qed
% 已译源行 263--368。

\SGARefTarget{sga3:xxiv:1.5.2}
\subsubsection*{1.5.2. 存在性。}
只须在 \textup{(fpqc)} 拓扑下局部地证明存在性。设两个钉扎为
\[
  (T,M,R,\Delta,(X_\alpha)_{\alpha\in\Delta})
  \quad\text{和}\quad
  (T',M',R',\Delta',(X'_{\alpha'})_{\alpha'\in\Delta'})
\]
由极大环面的共轭性，可设 \(T=T'\)。将 \(S\) 限制到适当部分后，可设同构
\[
  D_S(M)\simeq D_S(M')
\]
来自一个将 \(R\) 映为 \(R'\) 的同构 \(M\simeq M'\)；从而归结到如下情形：
\[
  T=T',\qquad M=M',\qquad R=R'.
\]
由于单根系在 Weyl 群的作用下共轭（\SGARefLink{sga3:XXI:3.3.7}{第 XXI 讲 3.3.7}），也可设 \(\Delta=\Delta'\)。于是，对每个 \(\alpha\in\Delta\)，有标量 \(z_\alpha\in\mathbb G_m(S)\)，使得
\[
  X'_\alpha=z_\alpha X_\alpha,
\]
并且只须在 \textup{(fpqc)} 拓扑下局部地构造 \(T\) 的截面 \(t\)，使得对每个 \(\alpha\in\Delta\) 都有 \(\alpha(t)=z_\alpha\)。但以 \(\alpha\)，\(\alpha\in\Delta\)，为分量的态射
\[
  T\longrightarrow(\mathbb G_{m,S})^\Delta
\]
对偶于单射 \(\mathbb Z^\Delta\to M\)，因而忠实平坦。这就证明了 \SGARefLink{sga3:xxiv:1.5.2}{1.5.2}，从而也证明了 \SGARefLink{sga3:xxiv:1.3}{1.3}。

\SGARefTarget{sga3:xxiv:1.6}
\paragraph{\SGARefLink{sga3:xxiv:1.6}{推论 1.6}.}
\emph{设 \(S\) 是概形，\(G\) 是约化 \(S\)-群。下列条件等价：}
\begin{enumerate}
\item[\SGAInlineRefTarget{sga3:xxiv:1.6:item:i}(i)] \(\underline{\Aut}_{S\text{-}\mathrm{gr.}}(G)\) 在 \(S\) 上仿射\emph{（分别为有限型、有限呈示、拟紧）}；
\item[\SGAInlineRefTarget{sga3:xxiv:1.6:item:ii}(ii)] \(\underline{\Autext}(G)\) 在 \(S\) 上有限；
\item[\SGAInlineRefTarget{sga3:xxiv:1.6:item:iii}(iii)] 对每个 \(s\in S\)，
  \[
    \operatorname{rg}_{\mathrm{red}}(G_s)
    -\operatorname{rg}_{\mathrm{ss}}(G_s)\leq 1.
  \]
\end{enumerate}

事实上，由于 \(\operatorname{ad}(G)\) 在 \(S\) 上仿射、平坦且有限呈示，态射
\[
  p:\underline{\Aut}_{S\text{-}\mathrm{gr.}}(G)
  \longrightarrow\underline{\Autext}(G)
\]
仿射、忠实平坦且有限呈示。

若 \(\underline{\Autext}(G)\) 在 \(S\) 上有限，则它在 \(S\) 上仿射，因而 \(\underline{\Aut}_{S\text{-}\mathrm{gr.}}(G)\) 也是如此；这证明了 \textup{\SGARefLink{sga3:xxiv:1.6:item:ii}{(ii)}}\(\Rightarrow\)\textup{\SGARefLink{sga3:xxiv:1.6:item:i}{(i)}}。若 \(\underline{\Aut}_{S\text{-}\mathrm{gr.}}(G)\) 在 \(S\) 上拟紧，则它在 \(S\) 上有限呈示，因为它总是在 \(S\) 上局部有限呈示且分离。由\SGARefLink{sga3:v:proposition:9:1}{第 V 讲 9.1}，\(\underline{\Autext}(G)\) 也在 \(S\) 上有限呈示，因而有限；这证明了 \textup{\SGARefLink{sga3:xxiv:1.6:item:i}{(i)}}\(\Rightarrow\)\textup{\SGARefLink{sga3:xxiv:1.6:item:ii}{(ii)}}。最后，为证明 \textup{\SGARefLink{sga3:xxiv:1.6:item:ii}{(ii)}} 与 \textup{\SGARefLink{sga3:xxiv:1.6:item:iii}{(iii)}} 等价，可以假设 \(G\) 分裂，此时断言归结为\SGARefLink{sga3:XXI:6.7.8}{第 XXI 讲 6.7.8}。

\SGARefTarget{sga3:xxiv:1.7}
\paragraph{\SGARefLink{sga3:xxiv:1.7}{推论 1.7}.}
\emph{设 \(S\) 是概形，\(G\) 是约化 \(S\)-群。则}
\[
  \underline{\Aut}_{S\text{-}\mathrm{gr.}}(G)^0
  \simeq\operatorname{ad}(G).
\]

\SGARefTarget{sga3:xxiv:1.8}
\paragraph{\SGARefLink{sga3:xxiv:1.8}{推论 1.8}.}
\emph{设 \(S\) 是概形，\(G\) 是约化 \(S\)-群，\(H\) 是在 \(S\) 上光滑、仿射且纤维连通的 \(S\)-群概形。则 \(S\)-函子}
\[
  \underline{\operatorname{Isom}}_{S\text{-}\mathrm{gr.}}(G,H)
\]
\emph{可由一个光滑分离的 \(S\)-概形表示；当 \(G\) 半单时，该概形在 \(S\) 上仿射。}

事实上，设 \(U\) 为使 \(H_s\) 约化的点 \(s\in S\) 的集合；这是开子集（\SGARefLink{sga3:xix:corollary:2:6}{第 XIX 讲 2.6}）。若 \(S'\) 是 \(S\)-概形，则 \(H_{S'}\) 约化，当且仅当 \(S'\to S\) 经 \(U\) 分解。因此典范态射
\[
  \underline{\operatorname{Isom}}_{S\text{-}\mathrm{gr.}}(G,H)
  \longrightarrow S
\]
经 \(U\) 分解。于是可设 \(S=U\)，从而归结到下面的推论。
% 已译源行 370--479。

\SGARefTarget{sga3:xxiv:1.9}
\paragraph{\SGARefLink{sga3:xxiv:1.9}{推论 1.9}.}
\emph{设 \(S\) 是概形，\(G\) 与 \(G'\) 是两个约化 \(S\)-群。则}
\[
  F=\underline{\operatorname{Isom}}_{S\text{-}\mathrm{gr.}}(G,G')
\]
\emph{可由一个光滑分离的 \(S\)-概形表示；当 \(G\) 或 \(G'\) 半单时，它是仿射的。此外，\(S\) 是两个开子概形 \(S_1\)、\(S_2\) 的不交并，使得 \(F_{S_1}=\varnothing\)，而 \(F_{S_2}\) 是 \(\underline{\Aut}_{S_2\text{-}\mathrm{gr.}}(G'_{S_2})\)（分别为 \(\underline{\Aut}_{S_2\text{-}\mathrm{gr.}}(G_{S_2})\)）下的左（分别为右）主齐性丛。}

事实上，设 \(S_2\) 是 \(S\) 中使 \(G\) 与 \(G'\) 类型相同的点的集合，\(S_1\) 是它的补集。由于约化群的类型是局部常值函数，\(S_1\) 与 \(S_2\) 都开。显然 \(F_{S_1}=\varnothing\)，并且可设 \(S=S_2\)。由\SGARefLink{sga3:XXIII:5.6}{第 XXIII 讲 5.6}，\(F\) 是 \(\underline{\Aut}_{S\text{-}\mathrm{gr.}}(G)\) 下的主齐性层，在平展拓扑下局部平凡。因此
\[
  F_0=F/\operatorname{ad}(G)
\]
是 \(\underline{\Autext}(G)\) 下的主齐性层，在平展拓扑下局部平凡，因而可表（\SGARefLink{sga3:x:corollary:5-5}{第 X 讲 5.5}）。由于 \(\operatorname{ad}(G)\) 仿射，故 \(F\) 也可表。

注意，这个证明还给出以下结果。
\SGARefTarget{sga3:xxiv:1.10}
\paragraph{\SGARefLink{sga3:xxiv:1.10}{推论 1.10}.}
\emph{设 \(S\) 是概形，\(G\) 与 \(G'\) 是两个约化 \(S\)-群，它们在每个 \(s\in S\) 处具有相同的类型。则 \(\operatorname{ad}(G)\) 在 \(\underline{\operatorname{Isom}}_{S\text{-}\mathrm{gr.}}(G,G')\) 上自由右作用。商层}
\[
  \underline{\operatorname{Isomext}}(G,G')
  =
  \underline{\operatorname{Isom}}_{S\text{-}\mathrm{gr.}}(G,G')
  /\operatorname{ad}(G)
\]
\emph{可由一个扭常 \(S\)-概形表示；它是 \(\underline{\Autext}(G)\) 下的主齐性丛，故当 \(G\) 半单时在 \(S\) 上有限。此外，同构}
\[
  \underline{\operatorname{Isom}}_{S\text{-}\mathrm{gr.}}(G,G')
  \simeq
  \underline{\operatorname{Isom}}_{S\text{-}\mathrm{gr.}}(G',G),
  \qquad u\longmapsto u^{-1},
\]
\emph{诱导一个同构}
\[
  \underline{\operatorname{Isomext}}(G,G')
  \simeq
  \underline{\operatorname{Isomext}}(G',G).
\]

\SGARefTarget{sga3:xxiv:1.11}
\paragraph{\SGARefLink{sga3:xxiv:1.11}{注 1.11}.}
若 \(\underline{\operatorname{Isomext}}(G,G')(S)\neq\varnothing\)，则称 \(G\) 为 \(G'\) 的\emph{内扭形式}（inner twisted form）。此时 \(G'\) 也是 \(G\) 的内扭形式，并且 \(\underline{\operatorname{Isom}}_{S\text{-}\mathrm{gr.}}(G,G')\) 的结构群可约化到 \(\operatorname{ad}(G)\)。更精确地，设
\[
  u\in\underline{\operatorname{Isomext}}(G,G')(S)
\]
并将 \(u\) 视为截面 \(u:S\to\underline{\operatorname{Isomext}}(G,G')\)。以
\[
  \underline{\operatorname{Isomint}}_u(G,G')
\]
表示 \(\underline{\operatorname{Isomext}}(G,G')\) 中作为 \(u\) 的像的闭子概形在典范态射
\[
  \underline{\operatorname{Isom}}_{S\text{-}\mathrm{gr.}}(G,G')
  \longrightarrow
  \underline{\operatorname{Isomext}}(G,G'),
\]
下的逆像。\(\operatorname{ad}(G)\) 在 \(\underline{\operatorname{Isomint}}_u(G,G')\) 上的自然作用使这个概形成为主齐性丛。沿
\[
  \operatorname{ad}(G)\longrightarrow\underline{\Aut}(G),
\]
扩张结构群，\footnote{底本的编者注将 \(\operatorname{ad}(G)\to\underline{\Autext}(G)\) 改为 \(\operatorname{ad}(G)\to\underline{\Aut}(G)\)。}就从 \(\underline{\operatorname{Isomint}}_u(G,G')\) 恢复了 \(\underline{\operatorname{Isom}}_{S\text{-}\mathrm{gr.}}(G,G')\)。

由 Hensel 引理（\SGARefLink{sga3:xi:corollary:1:11}{第 XI 讲 1.11}），\SGARefLink{sga3:xxiv:1.8}{推论 1.8} 立即给出：
% 已译源行 481--589。

\SGARefTarget{sga3:xxiv:1.12}
\paragraph{\SGARefLink{sga3:xxiv:1.12}{推论 1.12}.}
\emph{设 \(S\) 是 Hensel 局部概形，\(G\) 是约化 \(S\)-群，\(G'\) 是具有连通纤维的光滑仿射 \(S\)-群，\(s\) 是 \(S\) 的闭点。若 \(G_s\) 与 \(G'_s\) 是 \(\kappa(s)\) 上同构的代数群，则 \(G\) 与 \(G'\) 同构。更精确地，每个 \(\kappa(s)\)-同构}
\[
  G_s\simeq G'_s
\]
\emph{都来自一个 \(S\)-同构 \(G\simeq G'\)。}

将\SGARefLink{sga3:xxiv:1.7}{推论 1.7}（分别为\SGARefLink{sga3:xxiv:1.12}{推论 1.12}）用于域上的对偶数概形，由\SGARefLink{sga3:iii:remark:2-10}{第 III 讲 2.10}（分别为 \SGARefLink{sga3:xxiv:3.10}{3.10}），即得下一推论的 \textup{(i)}（分别为 \textup{(ii)}）。

\SGARefTarget{sga3:xxiv:1.13}
\paragraph{\SGARefLink{sga3:xxiv:1.13}{推论 1.13}.}
\emph{设 \(k\) 是域，\(G\) 是约化 \(k\)-群。}
\begin{enumerate}
\item[\SGAInlineRefTarget{sga3:xxiv:1.13:item:i}(i)] \emph{若 \(G\) 是伴随的，则}
\[
  H^1(G,\Lie(G/k))=0.
\]
\footnote{底本的编者注修正了原文；原文在未假定 \(G\) 为伴随群的情况下陈述了 \textup{\SGARefLink{sga3:xxiv:1.13:item:i}{(i)}}。由 \(H^i(G,-)\) 作为 \(\operatorname{Hom}_G(k,-)\) 的导出函子的描述（\SGARefLink{sga3:i:theorem:5-3-1}{第 I 讲 5.3.1}；另见 [Ja87]，I，4.16），对每个 \(G\)-模 \(V\) 有 \(H^i(G,V)=\operatorname{Ext}^i_G(k,V)\)。若 \(\operatorname{char}(k)=p>0\) 且 \(G=\operatorname{SL}_{p,k}\)，则 \(\Lie(\operatorname{GL}_p/k)\) 是 \(k\) 被 \(\mathfrak g=\Lie(\operatorname{SL}_p/k)\) 的非平凡扩张，故 \(H^1(G,\mathfrak g)\neq0\)。另见下文增补 \SGARefLink{sga3:xxiv:1.15.1}{1.15.1}。}
\item[\SGAInlineRefTarget{sga3:xxiv:1.13:item:ii}(ii)] \emph{有}
\[
  H^2(G,\Lie(G/k))=0.
\]
\end{enumerate}

\SGARefTarget{sga3:xxiv:1.14}
\paragraph{\SGARefLink{sga3:xxiv:1.14}{注 1.14}.}
\begin{enumerate}
\item[(i)] 关于 \(H^1\) 的断言是 Chevalley 已知的；关于 \(H^2\) 的断言，Chevalley 已在分类中的大多数情形下证明。
\item[(ii)] 事实上，\SGARefLink{sga3:xxiv:1.13}{推论 1.13} 与代数闭域上的唯一性定理结合，实质上等价于唯一性定理。因此，\SGARefLink{sga3:xxiv:1.13}{推论 1.13} 的直接证明将给出一种途径，由 Chevalley 的域上唯一性定理推出一般的唯一性定理。
\end{enumerate}

每个概形上都存在任意类型的约化群（第 XXV 讲），这表明，将约化 \(k\)-群 \(G\) 提升到剩余域为 \(k\) 的 Artin 环上的障碍均为零；根据\SGARefLink{sga3:iii:corollary:3-8}{第 III 讲 3.8}，这些障碍属于
\[
  H^3(G,\Lie(G/k)\otimes V)
  \simeq H^3(G,\Lie(G/k))\otimes V,
\]
其中 \(V\) 是赋予 \(G\) 的平凡作用的某个 \(k\)-向量空间。这提示
\[
  H^3(G,\Lie(G/k))=0.
\]
若此断言为真，则其直接证明无疑同样会给出一种途径，由域上的存在性定理（Chevalley 的 Tôhoku 论文）推出一般的存在性定理。

\SGARefTarget{sga3:xxiv:1.15}
\paragraph{\SGARefLink{sga3:xxiv:1.15}{推论 1.15}.}
\emph{设 \(k\) 是域，\footnote{底本的编者注将“概形”改为“域”。}\(G\) 是约化 \(k\)-群。将 \(k\) 视为平凡 \(G\)-模。则}
\[
  H^1(G,k)=H^2(G,k)=0.
\]
\footnote{鉴于 \SGARefLink{sga3:xxiv:1.13}{1.13} 中所作的修正，底本编者另给出 \(H^1\) 的证明。此外，G.~Kempf 的一个定理给出：对每个 \(i>0\)，都有 \(H^i(G,k)=0\)；参见 [Ja87]，II，\SGARefLink{sga3:ii:proposition:4-5}{4.5} 与 4.11。}
由于
\[
  H^i(G_{\bar k},\bar k)=H^i(G,k)\otimes_k\bar k,
\]
可设 \(k\) 代数闭。\(H^1(G,k)\) 的一个元素无非是 \(k\)-群态射
\[
  \varphi:G\longrightarrow\mathbb G_{a,k}.
\]
于是
\[
  \varphi(G)=G/\ker(\varphi)
\]
是 \(\mathbb G_{a,k}\) 的光滑连通约化子群（参见 XIX，\SGARefLink{sga3:xix:paragraph:1-7}{1.7}），因而平凡。所以 \(H^1(G,k)=0\)。
% 已译源行 591--697。

现在考虑约化 \(k\)-群
\[
  H=G\times_k\mathbb G_{m,k}.
\]
有一个 \(H\)-稳定的分解
\[
  \Lie(H/k)=\Lie(G/k)\oplus k.
\]
对每个 \(H\)-模 \(V\)，
\[
  H^i(H,V)=H^i\bigl(G,H^0(\mathbb G_{m,k},V)\bigr).
\]
事实上，\(H^i(H,-)\) 是
\[
  H^0(H,-)
  =H^0\bigl(G,H^0(\mathbb G_{m,k},-)\bigr),
\]
的导出函子，而函子 \(H^0(\mathbb G_{m,k},-)\) 正合（参见\SGARefLink{sga3:i:theorem:5-3-1}{第 I 讲 5.3.1} 和 \SGARefLink{sga3:i:theorem:5-3-3}{5.3.3}）。特别地，对每个 \(i\)，
\[
  H^i(H,\Lie(H/k))
  =
  H^i(G,\Lie(G/k))\oplus H^i(G,k).
\]
将 \SGARefLink{sga3:xxiv:1.13:item:ii}{1.13(ii)} 用于约化群 \(H\)，即得 \(H^2(G,k)=0\)。

\SGARefTarget{sga3:xxiv:1.15.1}
\medskip\noindent\textbf{推论 1.15.1.} ---
\footnote{编者注：根据 O.~Gabber 的评论增补这一推论，以精确化 \SGARefLink{sga3:xxiv:1.13:item:i}{1.13(i)}。}
{\itshape 设 \(k\) 是域，\(G\) 是约化 \(k\)-群，\(Z\) 是它的中心，\(\mathscr R=(M,M^*,R,R^*)\) 是 \(G\) 的类型。则
\[
  H^1(G,\Lie(G/k))
  \simeq
  \operatorname{Ext}^1_{\mathbb Z}(M/\Gamma_0(R),k).
\]
特别地，\(H^1(G,\Lie(G/k))=0\)，当且仅当 \(Z\) 在 \(k\) 上光滑。}

设 \(\mathfrak g\)、\(\mathfrak z\)、\(\mathfrak g_{\mathrm{ad}}\) 分别是 \(G\)、\(Z\)、\(G_{\mathrm{ad}}=G/Z\) 的 Lie 代数。由 \SGARefLink{sga3:xxiv:1.15}{1.15} 及其证明，得到
\[
  H^1(G,\mathfrak g)
  =
  H^1(G_{\mathrm{ad}},\mathfrak g)
  \simeq
  H^1(G_{\mathrm{ad}},\mathfrak g/\mathfrak z).
\]
置 \(C=\operatorname{Coker}(\mathfrak g\to\mathfrak g_{\mathrm{ad}})\)。有正合序列
\[
  0\longrightarrow\mathfrak g/\mathfrak z
  \longrightarrow\mathfrak g_{\mathrm{ad}}
  \longrightarrow C\longrightarrow0.
\]
由 \SGARefLink{sga3:xxiv:1.13}{1.13}，有 \(H^1(G_{\mathrm{ad}},\mathfrak g_{\mathrm{ad}})=0\)，又有
\[
  H^0(G_{\mathrm{ad}},\mathfrak g_{\mathrm{ad}})=0
\]
（参见\SGARefLink{sga3:ii:theorem:5-2-3}{第 II 讲 5.2.3}），故得到
\[
  H^1(G_{\mathrm{ad}},\mathfrak g/\mathfrak z)
  \simeq H^0(G_{\mathrm{ad}},C).
\]
为计算右端，可设 \(k\) 代数闭。设 \((G,T,M,R)\) 是 \(G\) 的一个分裂结构。则
\[
  Z=D_k(M/\Gamma_0(R))
\]
并且 \(C\) 等同于
\[
  \operatorname{Coker}\!\left(
    \operatorname{Hom}_{\mathbb Z}(M,k)
    \longrightarrow
    \operatorname{Hom}_{\mathbb Z}(\Gamma_0(R),k)
  \right)
  \simeq
  \operatorname{Ext}^1_{\mathbb Z}(M/\Gamma_0(R),k),
\]
其上的 \(G_{\mathrm{ad}}\)-作用平凡。推论随之得到，因为 \(Z\) 在 \(k\) 上不光滑，当且仅当 \(\operatorname{char}(k)=p>0\) 且 \(M/\Gamma_0(R)\) 有 \(p\)-挠元（\SGARefLink{sga3:ix:proposition:2:1}{第 IX 讲 2.1}）。

\SGARefTarget{sga3:xxiv:1.16}
\medskip\noindent\textbf{定义 1.16.} ---
{\itshape 设 \(S\) 是概形，\(G\) 是约化 \(S\)-群。\emph{\(G\) 在 \(S\) 上的形式}是一个 \(S\)-群概形 \(G'\)，它在 \textup{(fpqc)} 拓扑下局部同构于 \(G\)。等价地，也可要求 \(G'\) 在平展拓扑下局部同构于 \(G\)（参见\SGARefLink{sga3:XXIII:5.6}{第 XXIII 讲 5.6}），或要求 \(G'\) 为约化 \(S\)-群，且在 \(S\) 的每一点处与 \(G\) 具有相同类型。}

\SGARefTarget{sga3:xxiv:1.17}
\medskip\noindent\textbf{推论 1.17.} ---
{\itshape 设 \(S\) 是概形，\(G\) 是约化 \(S\)-群。
\begin{enumerate}
\item[(i)] 函子
\[
  G'\longmapsto
  \underline{\operatorname{Isom}}_{S\text{-}\mathrm{gr}}(G,G')
\]
给出 \(G\) 在 \(S\) 上的形式所成范畴与 \(\underline{\Aut}_{S\text{-}\mathrm{gr}}(G)\) 下主齐性丛所成范畴之间的等价。
\item[(ii)] 若 \(S'\to S\) 是覆盖态射，则由 \(S'\) 平凡化的 \(G\) 的形式对应于由 \(S'\) 平凡化的丛。
\item[\SGAInlineRefTarget{sga3:xxiv:1.17:item:iii}(iii)] \(\underline{\Aut}_{S\text{-}\mathrm{gr}}(G)\) 下的每个主齐性层都可表且拟等平凡，即在平展拓扑下局部平凡。
\end{enumerate}}

第一个断言在层范畴中是形式的，例如可取 \textup{(fpqc)} 拓扑。每个在 \textup{(fpqc)} 拓扑下局部同构于 \(G\) 的层都可表，因为 \(G\) 在 \(S\) 上仿射；而且它在平展拓扑下局部同构于 \(G\)。最后，对 \(G\) 的每个形式 \(G'\)，\(S\)-层 \(\underline{\operatorname{Isom}}_{S\text{-}\mathrm{gr}}(G,G')\) 由 \SGARefLink{sga3:xxiv:1.8}{1.8} 可表。推论即得。

\SGARefTarget{sga3:xxiv:1.18}
\medskip\noindent\textbf{推论 1.18.} ---
{\itshape 约化群 \(G\) 在 \(S\) 上的形式的同构类集合，同构于
\[
  H^1\!\left(S,\underline{\Aut}_{S\text{-}\mathrm{gr}}(G)\right)
  =
  H^1_{\mathrm{\acute et}}\!
    \left(S,\underline{\Aut}_{S\text{-}\mathrm{gr}}(G)\right)
  =
  \Fib\!\left(S,\underline{\Aut}_{S\text{-}\mathrm{gr}}(G)\right).
\]
若 \(S'\to S\) 是覆盖态射，则由 \(S'\) 平凡化的形式所成的子集同构于
\[
  H^1\!\left(S'/S,\underline{\Aut}_{S\text{-}\mathrm{gr}}(G)\right).
\]}

\SGARefTarget{sga3:xxiv:1.19}
\medskip\noindent\textbf{推论 1.19.} ---
{\itshape 设 \(S\) 是概形，\(\mathscr R\) 是钉扎约化根数据，使得 \(\operatorname{\acute{E}p}_S(\mathscr R)\) 存在；\footnote{编者注：参见 XXIII，\SGARefLink{sga3:XXIII:5.11}{定义 5.11}。}由第 XXV 讲，这一条件总是满足。置
\[
  A_S(\mathscr R)
  =
  \underline{\Aut}_{S\text{-}\mathrm{gr}}
    (\operatorname{\acute{E}p}_S(\mathscr R))
  =
  \operatorname{ad}(\operatorname{\acute{E}p}_S(\mathscr R))
  \cdot E(\mathscr R)_S.
\]
\begin{enumerate}
\item[\SGAInlineRefTarget{sga3:xxiv:1.19:item:i}(i)] 类型为 \(\mathscr R\) 的约化 \(S\)-群（\SGARefLink{sga3:xxii:def:2.7}{第 XXII 讲 2.7}）的同构类集合，由\SGARefLink{sga3:XXIII:5.12}{第 XXIII 讲 5.12}，同构于
\[
  H^1(S,A_S(\mathscr R))
  =
  H^1_{\mathrm{\acute et}}(S,A_S(\mathscr R))
  =
  \Fib(S,A_S(\mathscr R)).
\]
\item[\SGAInlineRefTarget{sga3:xxiv:1.19:item:ii}(ii)] 若 \(S'\to S\) 是覆盖态射，则由在 \(S'\) 上可分裂的群的类组成的子集，同构于 \(H^1(S'/S,A_S(\mathscr R))\)。
\end{enumerate}}

\SGARefTarget{sga3:xxiv:1.20}
\medskip\noindent\textbf{注 1.20.} ---
沿用上述记号，每个类型为 \(\mathscr R\) 的约化 \(S\)-群 \(G\) 都典范地伴有一个 \(A_S(\mathscr R)\) 下的右主齐性丛：
\[
  \underline{\operatorname{Isom}}_{S\text{-}\mathrm{gr}}
    (\operatorname{\acute{E}p}_S(\mathscr R),G)=P.
\]
概形 \(P\) 可以解释为“\(G\) 的类型为 \(\mathscr R\) 的钉扎概形”（参见 \SGARefLink{sga3:xxiv:1.0}{1.0}）。它也是 \(\underline{\Aut}_{S\text{-}\mathrm{gr}}(G)\) 下的左主齐性丛；从上述描述立即可见这一结构。

\SGARefTarget{sga3:xxiv:1.21}
\medskip\noindent\textbf{命题 1.21.} ---
{\itshape 设 \(S\) 是 Hensel 局部概形，\(s\) 是它的闭点。函子
\[
  G\longmapsto G_s
\]
从约化 \(S\)-群的同构类集合到约化 \(\kappa(s)\)-群的同构类集合诱导双射。特别地，对每个约化 \(S\)-群 \(G\)，存在有限平展满射态射 \(S'\to S\)，使得 \(G_{S'}\) 可分裂。}

利用 \(\operatorname{\acute{E}p}_S(\mathscr R)\) 的存在性（第 XXV 讲），只须取 \(H=A_S(\mathscr R)\)，证明典范映射
\[
  \Fib(S,H)\longrightarrow\Fib(\kappa(s),H_s)
\]
为双射，并且 \(\Fib(S,H)\) 的每个元素都有上述性质。\(H\) 的每个有限子集都包含于一个仿射开子概形：对常群这是显然的，而 \(H\) 在一个常群上仿射。因此可应用附录中证明的结果 \SGARefLink{sga3:xxiv:proposition:8-1}{8.1}。

\SGARefTarget{sga3:xxiv:section:2}
\section*{2. 自同构与子群}

引入以下记号。若
\[
  \mathcal H=\underline{\Aut}_{S\text{-}\mathrm{gr}}(G)
\]
且 \(X\) 是 \(G\) 的子函子，则置
\[
  \underline{\Aut}_{S\text{-}\mathrm{gr}}(G,X)
  =\operatorname{Norm}_{\mathcal H}(X),
  \qquad
  \underline{\Aut}_{S\text{-}\mathrm{gr}}(G,\operatorname{id}_X)
  =\underline{\Centr}_{\mathcal H}(X).
\]
若 \(Y\) 是 \(G\) 的另一个子函子，则定义
\[
  \underline{\Aut}_{S\text{-}\mathrm{gr}}(G,X,Y)
  =
  \underline{\Aut}_{S\text{-}\mathrm{gr}}(G,X)
  \cap
  \underline{\Aut}_{S\text{-}\mathrm{gr}}(G,Y).
\]
若 \(G'\) 是另一个 \(S\)-群，\(X'\) 是 \(G'\) 的子函子，\footnote{英译者注：法文此处印为 \(G\)；由定义中各对象的类型可知，应为 \(G'\)。}则以
\[
  \underline{\operatorname{Isom}}_{S\text{-}\mathrm{gr}}
    (G,X;G',X')
\]
表示 \(\underline{\operatorname{Isom}}_{S\text{-}\mathrm{gr}}(G,G')\) 的如下子函子：对每个 \(S'\to S\)，它由
\[
  \left\{
    u\in
    \underline{\operatorname{Isom}}_{S\text{-}\mathrm{gr}}(G,G')(S')
    \,\middle|\,
    u(X_{S'})=X'_{S'}
  \right\}.
\]
定义。同样定义记号 \(\underline{\operatorname{Isom}}_{S\text{-}\mathrm{gr}}(G,X,Y;G',X',Y')\) 等。\footnote{编者注：上述定义已予写明。原文仅说以类似方式定义 \(\underline{\Aut}_{S\text{-}\mathrm{gr}}(G,X,Y),\ldots\) 与 \(\underline{\operatorname{Isom}}_{S\text{-}\mathrm{gr}}(G,X;G',X'),\ldots\)。}

\SGARefTarget{sga3:xxiv:2.1}
\medskip\noindent\textbf{命题 2.1.} ---
{\itshape 设 \(S\) 是概形，\(G\) 是约化 \(S\)-群，\(T\) 是 \(G\) 的极大环面（分别地，\(B\) 是 \(G\) 的 Borel 子群，或 \(B\supset T\) 是 \(G\) 的 Killing 偶）。以 \(T_{\mathrm{ad}}\)（分别为 \(B_{\mathrm{ad}}\)）表示 \(\operatorname{ad}(G)\) 中相应的极大环面（分别为 Borel 子群）：
\[
  B_{\mathrm{ad}}
  \simeq B/\underline{\Centr}(G)
  =B/\underline{\Centr}(B),
  \qquad
  T_{\mathrm{ad}}\simeq T/\underline{\Centr}(G).
\]
则 \(\underline{\Aut}_{S\text{-}\mathrm{gr}}(G,T)\)（分别为 \(\underline{\Aut}_{S\text{-}\mathrm{gr}}(G,B)\)，或 \(\underline{\Aut}_{S\text{-}\mathrm{gr}}(G,B,T)\)）由 \(\underline{\Aut}_{S\text{-}\mathrm{gr}}(G)\) 的闭子概形表示，它在 \(S\) 上光滑，而且\SGARefLink{sga3:xxiv:1.3}{定理 1.3} 的正合序列诱导以下正合序列：
\[
  1\longrightarrow
  \operatorname{Norm}_{\operatorname{ad}(G)}(T_{\mathrm{ad}})
  \longrightarrow
  \underline{\Aut}_{S\text{-}\mathrm{gr}}(G,T)
  \longrightarrow
  \underline{\Autext}(G)
  \longrightarrow1,
\]
\[
  1\longrightarrow B_{\mathrm{ad}}
  \longrightarrow
  \underline{\Aut}_{S\text{-}\mathrm{gr}}(G,B)
  \longrightarrow
  \underline{\Autext}(G)
  \longrightarrow1,
\]
\[
  1\longrightarrow T_{\mathrm{ad}}
  \longrightarrow
  \underline{\Aut}_{S\text{-}\mathrm{gr}}(G,B,T)
  \longrightarrow
  \underline{\Autext}(G)
  \longrightarrow1.
\]}

由闭子概形的下降，\footnote{编者注：参见 SGA~1，VIII，1.9。}立即归结到如下情形：\(G\) 已钉扎，\(B\supset T\) 是它的典范 Killing 偶（参见\SGARefLink{sga3:xxii:cor:5.5.5:item:iv}{第 XXII 讲 5.5.5(iv)}）。由于 \SGARefLink{sga3:xxiv:1.2}{1.2} 中的群 \(E\) 正规化 \(B\) 和 \(T\)，结果立即由 \(\operatorname{ad}(G)\) 中的正规化定理（\SGARefLink{sga3:xxii:cor:5.3.12}{第 XXII 讲 5.3.12} 和 5.6.1）推出。

现在使用共轭定理（参见\SGARefLink{sga3:XXIII:5.12}{第 XXIII 讲 5.12}），并仿照第 1 节论证，得到：

\SGARefTarget{sga3:xxiv:2.2}
\medskip\noindent\textbf{推论 2.2.} ---
{\itshape 设 \(S\) 是概形，\(G\)、\(G'\) 是在每一点处类型相同的两个约化 \(S\)-群，\(B\supset T\)（分别为 \(B'\supset T'\)）是 \(G\)（分别为 \(G'\)）的 Killing 偶。

\textup{\SGARefLink{sga3:xxiv:1.19:item:i}{(i)}} \(S\)-函子
\[
  \underline{\operatorname{Isom}}_{S\text{-}\mathrm{gr}}
    (G,T;G',T')
\]
由 \(\underline{\operatorname{Isom}}_{S\text{-}\mathrm{gr}}(G,G')\) 的一个光滑闭子概形表示，它是 \(\underline{\Aut}_{S\text{-}\mathrm{gr}}(G,T)\) 下的主齐性空间。此外，\(\operatorname{Norm}_{\operatorname{ad}(G)}(T_{\mathrm{ad}})\) 在这个概形上自由作用，并有典范同构
\[
  \underline{\operatorname{Isom}}_{S\text{-}\mathrm{gr}}
    (G,T;G',T')
  \big/
  \operatorname{Norm}_{\operatorname{ad}(G)}(T_{\mathrm{ad}})
  \simeq
  \underline{\operatorname{Isomext}}(G,G').
\]

\textup{\SGARefLink{sga3:xxiv:1.19:item:ii}{(ii)}} \(S\)-函子
\[
  \underline{\operatorname{Isom}}_{S\text{-}\mathrm{gr}}
    (G,B;G',B')
\]
由 \(\underline{\operatorname{Isom}}_{S\text{-}\mathrm{gr}}(G,G')\) 的一个光滑闭子概形表示，它是 \(\underline{\Aut}_{S\text{-}\mathrm{gr}}(G,B)\) 下的主齐性空间。此外，\(B_{\mathrm{ad}}\) 在这个概形上自由作用，并有典范同构
\[
  \underline{\operatorname{Isom}}_{S\text{-}\mathrm{gr}}
    (G,B;G',B')
  \big/B_{\mathrm{ad}}
  \simeq
  \underline{\operatorname{Isomext}}(G,G').
\]

\textup{\SGARefLink{sga3:xxiv:1.17:item:iii}{(iii)}} \(S\)-函子
\[
  \underline{\operatorname{Isom}}_{S\text{-}\mathrm{gr}}
    (G,B,T;G',B',T')
\]
由 \(\underline{\operatorname{Isom}}_{S\text{-}\mathrm{gr}}(G,G')\) 的一个光滑闭子概形表示，它是 \(\underline{\Aut}_{S\text{-}\mathrm{gr}}(G,B,T)\) 下的主齐性空间。此外，\(T_{\mathrm{ad}}\) 在这个概形上自由作用，并有典范同构
\[
  \underline{\operatorname{Isom}}_{S\text{-}\mathrm{gr}}
    (G,B,T;G',B',T')
  \big/T_{\mathrm{ad}}
  \simeq
  \underline{\operatorname{Isomext}}(G,G').
\]}

再一次仿照第 1 节论证，得到：
\SGARefTarget{sga3:xxiv:2.3}
\medskip\noindent\textbf{推论 2.3.} ---
{\itshape 设 \(S\) 是概形，\(G\) 是约化 \(S\)-群，\(B\supset T\) 是 \(G\) 的 Killing 偶。函子}
\[
  (G',T')
  \longmapsto
  \underline{\operatorname{Isom}}_{S\text{-}\mathrm{gr}}
    (G,T;G',T'),
\]
{\itshape 分别地}
\[
  (G',B')
  \longmapsto
  \underline{\operatorname{Isom}}_{S\text{-}\mathrm{gr}}
    (G,B;G',B'),
\]
{\itshape 分别地}
\[
  (G',B',T')
  \longmapsto
  \underline{\operatorname{Isom}}_{S\text{-}\mathrm{gr}}
    (G,B,T;G',B',T'),
\]
{\itshape 给出以下两个范畴之间的等价：一方是有序对 \((G',T')\)（分别为有序对 \((G',B')\)、三元组 \((G',B',T')\)）所成范畴，其中 \(G'\) 是 \(G\) 的形式，\(T'\) 是 \(G'\) 的极大环面（分别地，\(B'\) 是 \(G'\) 的 Borel 子群；\(B'\supset T'\) 是 \(G'\) 的 Killing 偶）；另一方是 \(S\)-群 \(H\) 下主齐性丛所成范畴，其中 \(H=\underline{\Aut}_{S\text{-}\mathrm{gr}}(G,T)\)（分别为 \(H=\underline{\Aut}_{S\text{-}\mathrm{gr}}(G,B)\)、\(H=\underline{\Aut}_{S\text{-}\mathrm{gr}}(G,B,T)\)）。此外，\(H\) 下每个主齐性层都可表且拟等平凡，因此}
\[
  H^1(S,H)=H^1_{\mathrm{\acute et}}(S,H)=\Fib(S,H).
\]

\SGARefTarget{sga3:xxiv:2.4}
\medskip\noindent\textbf{注 2.4.} ---
在 \SGARefLink{sga3:xxiv:2.2}{2.2} 的假设下，在集合层面记作 \(u\mapsto u(T)\)（分别为 \(u\mapsto u(B)\)、\(u\mapsto(u(B),u(T))\)）的态射诱导同构
\[
  \underline{\operatorname{Isom}}_{S\text{-}\mathrm{gr}}(G,G')
  \big/
  \underline{\Aut}_{S\text{-}\mathrm{gr}}(G,T)
  \simeq \operatorname{Tor}(G'),
\]
分别地，
\[
  \underline{\operatorname{Isom}}_{S\text{-}\mathrm{gr}}(G,G')
  \big/
  \underline{\Aut}_{S\text{-}\mathrm{gr}}(G,B)
  \simeq \operatorname{Bor}(G'),
\]
分别地，
\[
  \underline{\operatorname{Isom}}_{S\text{-}\mathrm{gr}}(G,G')
  \big/
  \underline{\Aut}_{S\text{-}\mathrm{gr}}(G,B,T)
  \simeq \operatorname{Kil}(G').
\]
证明是直接的：只须在 \textup{(fpqc)} 下局部地证明，故可设 \(G\simeq G'\)，此时结果归结为\SGARefLink{sga3:xxii:cor:5.8.3:item:iii}{第 XXII 讲 5.8.3(iii)}。

\SGARefTarget{sga3:xxiv:2.5}
\medskip\noindent\textbf{注 2.5.} ---
上述结果立即可用结构群的约化来解释。若 \(G'\) 是 \(G\) 的形式，对应于主丛 \(\underline{\operatorname{Isom}}_{S\text{-}\mathrm{gr}}(G,G')\)，则给定这个丛的结构群到 \(\underline{\Aut}_{S\text{-}\mathrm{gr}}(G,T)\) 的约化，就等于给定 \(G'\) 的一个极大环面 \(T'\)。这里的双射，一方为 \SGARefLink{sga3:xxiv:2.4}{2.4} 中的双射，另一方为映射
\[
  T'\longmapsto
  \underline{\operatorname{Isom}}_{S\text{-}\mathrm{gr}}
    (G,T;G',T').
\]
Borel 子群与 Killing 偶的情形也相同。

\SGARefTarget{sga3:xxiv:2.6}
\medskip\noindent\textbf{命题 2.6.} ---
{\itshape 设 \(S\) 是概形，\(G\)、\(G'\) 是在每一点处类型相同的两个约化 \(S\)-群，\(T\)（分别为 \(T'\)）是 \(G\)（分别为 \(G'\)）的极大环面。则 \(T_{\mathrm{ad}}\) 在 \(\underline{\operatorname{Isom}}_{S\text{-}\mathrm{gr}}(G,T;G',T')\) 上自由作用，商}
\[
  P=
  \underline{\operatorname{Isom}}_{S\text{-}\mathrm{gr}}
    (G,T;G',T')
  \big/T_{\mathrm{ad}}
\]
{\itshape 可表；它是}
\[
  A=
  \underline{\Aut}_{S\text{-}\mathrm{gr}}(G,T)
  \big/T_{\mathrm{ad}},
\]
{\itshape 下的主齐性丛，其中 \(A\) 由扭常 \(S\)-概形表示，并且是 \(\underline{\Autext}(G)\) 被}
\[
  W_{\operatorname{ad}(G)}(T_{\mathrm{ad}})
  =
  \operatorname{Norm}_{\operatorname{ad}(G)}(T_{\mathrm{ad}})
  \big/T_{\mathrm{ad}}.
\]
{\itshape 的扩张。此外，若让 \(A\) 以显然的方式作用于 \(T\)，则与 \(P\) 相伴的丛恰为 \(T'\)。}\footnote{编者注：因此，若 \(P\) 有截面，则 \(T\) 与 \(T'\) 同构；特别地，\(T_P\) 与 \(T'_P\) 同构。参见增补 \SGARefLink{sga3:xxiv:proposition:4-5-1}{4.5.1}，那里在 \(G=\operatorname{\acute{E}p}_S(\mathscr R)\) 的情形下进一步讨论了这一观察。}

命题的第一部分立即由前述结果推出。为证明第二部分，注意有显然的态射
\[
  P\times_S T\longrightarrow T',
  \qquad (u,t)\longmapsto u(t).
\]
为证明取 \(A\) 的商后它诱导同构，可再设 \((G,T)\simeq(G',T')\)，此时断言直接成立。

完全类似地，有：
\SGARefTarget{sga3:xxiv:2.7}
\medskip\noindent\textbf{命题 2.7.} ---
{\itshape 设 \(S\) 是概形，\(G\)、\(G'\) 是在每一点处类型相同的两个约化 \(S\)-群，\(B\supset T\)（分别为 \(B'\supset T'\)）是 \(G\)（分别为 \(G'\)）的 Killing 偶。若让
\[
  \underline{\Aut}_{S\text{-}\mathrm{gr}}(G,B,T)
  \big/T_{\mathrm{ad}}
  \simeq \underline{\Autext}(G)
\]
以显然的方式作用于 \(T\)，则与 \(\underline{\operatorname{Isomext}}(G,G')\) 相伴的丛恰为 \(T'\)。}

\SGARefTarget{sga3:xxiv:2.8}
\medskip\noindent\textbf{推论 2.8.} ---
{\itshape 设 \(G\)、\(G'\) 是互为内扭形式的两个约化 \(S\)-群，\(B\supset T\)（分别为 \(B'\supset T'\)）是 \(G\)（分别为 \(G'\)）的 Killing 偶。则 \(T\) 与 \(T'\) 同构。}\footnote{编者注：这由 \SGARefLink{sga3:xxiv:2.7}{2.7} 推出，因为 \(\underline{\operatorname{Isomext}}(G,G')\) 有截面（参见前一编者注）。}
% 已译源行 699--1262。

\SGARefTarget{sga3:xxiv:2.9}
\medskip\noindent\textbf{注 2.9.} ---
一般并不能断言 \(B\) 与 \(B'\) 同构；不过，它们互为内扭形式（参见第 5 节）。

可以建立上述结果的“\(\operatorname{Isomint}\)”版本。\footnote{编者注：回忆 \(\underline{\operatorname{Isomint}}_u(G,G')\) 定义于 \SGARefLink{sga3:xxiv:1.11}{1.11}。}这里举出一个：

\SGARefTarget{sga3:xxiv:2.10}
\medskip\noindent\textbf{命题 2.10.} ---
{\itshape 设 \(S\) 是概形，\(G\)、\(G'\) 是在每一点处类型相同的两个约化 \(S\)-群，\(B\supset T\)（分别为 \(B'\supset T'\)）是 \(G\)（分别为 \(G'\)）的 Killing 偶，并设 \(u\in\underline{\operatorname{Isomext}}(G,G')(S)\)。考虑}
\[
\begin{aligned}
  &\underline{\operatorname{Isomint}}_u
     (G,B,T;G',B',T')\\
  &\qquad =
  \underline{\operatorname{Isom}}_{S\text{-}\mathrm{gr}}
     (G,B,T;G',B',T')
  \cap
  \underline{\operatorname{Isomint}}_u(G,G').
\end{aligned}
\]
{\itshape 它是一个光滑仿射 \(S\)-概形，而且是 \(T_{\mathrm{ad}}\) 下的主齐性丛。特别地，}
\[
  \underline{\operatorname{Isomint}}_u
    (G,B,T;G',B',T')(S)\ne\varnothing
\]
{\itshape 当且仅当 \(H^1(S,T_{\mathrm{ad}})\) 中相应的元素为零。}

在本节末尾，我们证明以后要用到的两个结果。
\SGARefTarget{sga3:xxiv:2.11}
\medskip\noindent\textbf{命题 2.11.} ---
{\itshape 设 \(S\) 是概形，\(G\) 是约化 \(S\)-群，\(T\) 是 \(G\) 的极大环面。显然的态射}
\[
  T_{\mathrm{ad}}
  \xrightarrow{\ \sim\ }
  \underline{\Aut}_{S\text{-}\mathrm{gr}}(G,\operatorname{id}_T)
\]
{\itshape 是同构。}

这由前述断言推出。在 \SGARefLink{sga3:xxiv:1.5.2}{1.5.2} 的证明过程中也已给出了一个直接证明。

\SGARefTarget{sga3:xxiv:2.12}
\medskip\noindent\textbf{推论 2.12.} ---
{\itshape 在上述假设下，有序对 \((G',f)\) 所成范畴与 \(T_{\mathrm{ad}}\) 下主齐性丛所成范畴之间有等价，其中 \(G'\) 是 \(G\) 的形式，\(f\) 是从 \(T\) 到 \(G'\) 的一个极大环面的同构。}

\SGARefTarget{sga3:xxiv:2.13}
\medskip\noindent\textbf{推论 2.13.} ---
{\itshape 若 \(H^1(S,T_{\mathrm{ad}})=0\)，而 \(G'\) 是 \(G\) 的一个形式，且有与 \(T\) 同构的极大环面，则 \(G'\) 与 \(G\) 同构。}

\SGARefTarget{sga3:xxiv:2.14}
\medskip\noindent\textbf{推论 2.14.} ---
{\itshape 设 \(S\) 是满足 \(\operatorname{Pic}(S)=0\) 的概形，\(G\) 是类型恒定的约化 \(S\)-群。则 \(G\) 可分裂，当且仅当 \(G\) 有分裂极大环面。}

设 \(G\) 是约化群，\(\operatorname{rad}(G)\) 是它的根基（\SGARefLink{sga3:xxii:rem:4.3.9}{第 XXII 讲 4.3.9}）。由于 \(\operatorname{rad}(G)\) 在 \(G\) 中是中心的特征子群，有典范态射
\[
  q:\underline{\Autext}(G)
  \longrightarrow
  \underline{\Aut}_{S\text{-}\mathrm{gr}}
    \bigl(\operatorname{rad}(G)\bigr).
\]

\SGARefTarget{sga3:xxiv:2.15}
\medskip\noindent\textbf{命题 2.15.} ---
{\itshape 设 \(G\) 是约化 \(S\)-群。以下序列正合：}
\[
  1\longrightarrow
  \operatorname{ad}(G)
  \longrightarrow
  \underline{\Aut}_{S\text{-}\mathrm{gr}}
    \bigl(G,\operatorname{id}_{\operatorname{rad}(G)}\bigr)
  \xrightarrow{p}
  \underline{\Autext}(G)
  \xrightarrow{q}
  \underline{\Aut}_{S\text{-}\mathrm{gr}}
    \bigl(\operatorname{rad}(G)\bigr),
\]
{\itshape 并且 \(\ker(q)=\operatorname{im}(p)\) 是 \(\underline{\Autext}(G)\) 的一个既开又闭的子概形，在 \(S\) 上有限。}

可设 \(G\) 分裂。第一个断言是直接的，第二个断言由\SGARefLink{sga3:XXI:6.7.5}{第 XXI 讲 6.7.5} 与 6.7.7 推出。

记
\[
  H=
  \underline{\Aut}_{S\text{-}\mathrm{gr}}
    \bigl(G,\operatorname{id}_{\operatorname{rad}(G)}\bigr),
\]
得到：
\SGARefTarget{sga3:xxiv:2.16}
\medskip\noindent\textbf{推论 2.16.} ---
{\itshape 有序对 \((G',f)\) 所成范畴与某个 \(S\)-群概形 \(H\) 下主丛所成范畴之间有等价，其中 \(G'\) 是 \(G\) 的形式，\(f\) 是从 \(\operatorname{rad}(G)\) 到 \(G'\) 的根基的同构，而 \(H\) 满足存在正合序列}
\[
  1\longrightarrow \operatorname{ad}(G)
  \longrightarrow H
  \longrightarrow F
  \longrightarrow 1,
\]
{\itshape 其中 \(S\)-群 \(F\) 在 \(S\) 上平展且有限。}

\SGARefTarget{sga3:xxiv:section:3}
\section*{3. 约化群的 Dynkin 概形. 拟分裂群}

\SGARefTarget{sga3:xxiv:3.1}
\medskip\noindent\textbf{3.1.} ---
回忆（\SGARefLink{sga3:XXI:7.4.1}{第 XXI 讲 7.4.1}），Dynkin 图是一个有限集合，其结构由不同元素的若干有序对（键）以及到 \(\{1,2,3\}\) 的映射（长度）定义。每个钉扎约化根数据 \(\mathscr R\) 都伴有一个 Dynkin 图 \(\Delta(\mathscr R)\)，其底集是单根集。

\SGARefTarget{sga3:xxiv:3.2}
\medskip\noindent\textbf{3.2.} ---
设 \(S\) 是概形。Dynkin \(S\)-概形是一个有限扭常 \(S\)-概形 \(\Delta\)，赋予由 \(\Delta\times_S\Delta\) 的一个与对角线交为空的子概形 \(L\)，以及态射
\[
  \Delta\longrightarrow\{1,2,3\}_S.
\]
所定义的结构。若 \(S'\) 是连通 \(S\)-概形，\footnote{编者注：补入了 \(S'\) 连通这一假设。}则 \(\Delta(S')\) 自然带有 Dynkin 图结构。

立刻可定义以下概念：两个 Dynkin 概形的同构、Dynkin 概形的基扩张，以及与一个 Dynkin 图相伴的常 Dynkin 概形。

Dynkin 概形上的每个平展下降资料都是有效的。

\SGARefTarget{sga3:xxiv:3.3}
\medskip\noindent\textbf{3.3.} ---
我们将为每个约化 \(S\)-群 \(G\) 构造一个相伴的 Dynkin \(S\)-概形。先设 \(G\) 在 \(S\) 上可分裂。对 \(G\) 的每个钉扎 \(E\)，设 \(\Delta(E)\) 是与 \(E\) 所定义的钉扎根数据相伴的常 Dynkin 概形。若 \(E\)、\(E'\) 是 \(G\) 的两个钉扎，则由 \SGARefLink{sga3:xxiv:1.5}{1.5}，存在唯一的 \(G\) 在 \(S\) 上的内自同构，将 \(E\) 变为 \(E'\)。\(G\) 的这个自同构定义一个同构
\[
  a_{EE'}:\Delta(E)\xrightarrow{\ \sim\ }\Delta(E').
\]
这些 \(a_{EE'}\) 显然组成传递系统，因而可将各 \(\Delta(E)\) 等同（即取它们的归纳极限）；所得是一个常 Dynkin 概形，记为 \(\operatorname{Dyn}(G)\)。现在若 \(G\) 是任意约化 \(S\)-群，则存在平展覆盖族 \(\{S_i\to S\}\)，使得 \(G_{S_i}\) 可分裂。仿照上述论证，便得到各 \(\operatorname{Dyn}(G_{S_i})\) 上的典范下降资料；由下降便构造出 Dynkin \(S\)-概形 \(\operatorname{Dyn}(G)\)。

\SGARefTarget{sga3:xxiv:3.4}
\medskip\noindent\textbf{3.4.} ---
这个构造有以下性质（而且这些性质实质上刻画了它）：
\begin{enumerate}
\item[\SGAInlineRefTarget{sga3:xxiv:3.4:item:i}(i)] 每个约化 \(S\)-群 \(G\) 都伴有 Dynkin 概形 \(\operatorname{Dyn}(G)\)；每个同构
\[
  u:G\xrightarrow{\ \sim\ }G'
\]
都函子性地伴有同构
\[
  \operatorname{Dyn}(u):
  \operatorname{Dyn}(G)\xrightarrow{\ \sim\ }\operatorname{Dyn}(G').
\]
\item[\SGAInlineRefTarget{sga3:xxiv:3.4:item:ii}(ii)] 若 \(S'\) 是 \(S\)-概形，\(G\) 是约化 \(S\)-群，则
\[
  \operatorname{Dyn}(G\times_S S')
  \simeq
  \operatorname{Dyn}(G)\times_S S'.
\]
\item[(iii)] 若 \(E\) 是 \(G\) 的钉扎，定义了 Dynkin 图为 \(\Delta\) 的钉扎根数据，则
\[
  \operatorname{Dyn}(G)\simeq\Delta_S.
\]
\item[\SGAInlineRefTarget{sga3:xxiv:3.4:item:iv}(iv)] 若 \(u\) 是 \(G\) 的内自同构，则 \(\operatorname{Dyn}(u)\) 是 \(\operatorname{Dyn}(G)\) 的恒等自同构。
\end{enumerate}

\SGARefTarget{sga3:xxiv:3.5}
\medskip\noindent\textbf{3.5.} ---
设 \(S\) 是概形，\(G\) 是约化 \(S\)-群。显然，\(\operatorname{Dyn}(G)\) 作为 Dynkin 概形的自同构函子
\[
  \underline{\Aut}_{\mathrm{Dyn}}
    \bigl(\operatorname{Dyn}(G)\bigr)
\]
由一个有限扭常 \(S\)-概形表示。由 \SGARefLink{sga3:xxiv:3.4:item:i}{3.4(i)} 与 \SGARefLink{sga3:xxiv:3.4:item:ii}{(ii)}，有典范态射
\[
  \underline{\Aut}_{S\text{-}\mathrm{gr}}(G)
  \longrightarrow
  \underline{\Aut}_{\mathrm{Dyn}}
    \bigl(\operatorname{Dyn}(G)\bigr),
\]
而由 \SGARefLink{sga3:xxiv:3.4:item:iv}{(iv)}，它经由态射
\[
  \underline{\Autext}(G)
  \longrightarrow
  \underline{\Aut}_{\mathrm{Dyn}}
    \bigl(\operatorname{Dyn}(G)\bigr).
\]
分解。更一般地，若 \(G\)、\(G'\) 是两个约化 \(S\)-群，则有典范态射
\[
  \underline{\operatorname{Isomext}}(G,G')
  \longrightarrow
  \underline{\operatorname{Isom}}_{\mathrm{Dyn}}
    \bigl(\operatorname{Dyn}(G),\operatorname{Dyn}(G')\bigr).
\]
特别地，若 \(G'\) 是 \(G\) 的内扭形式（\SGARefLink{sga3:xxiv:1.11}{1.11}），则 Dynkin 概形 \(\operatorname{Dyn}(G)\) 与 \(\operatorname{Dyn}(G')\) 同构。

\SGARefTarget{sga3:xxiv:3.6}
\medskip\noindent\textbf{3.6.} ---
若 \(G\) 半单（分别为伴随或单连通），则态射
\[
  \underline{\Autext}(G)
  \longrightarrow
  \underline{\Aut}_{\mathrm{Dyn}}
    \bigl(\operatorname{Dyn}(G)\bigr)
\]
是单态射（分别为同构）。事实上，可设 \(G\) 已钉扎，此时断言归结为钉扎约化根数据的相应结果（参见\SGARefLink{sga3:XXI:7.4.5}{第 XXI 讲 7.4.5}）。

对 \(\operatorname{Isom}\) 函子也有类似结果。特别地，两个半单伴随（分别为单连通）\(S\)-群互为内扭形式，当且仅当它们的 Dynkin 概形同构。

\SGARefTarget{sga3:xxiv:3.7}
\medskip\noindent\textbf{3.7.} ---
还可以用另一种方式构造约化群所伴的 Dynkin 概形。设 \(\mathscr R\) 是钉扎约化根数据，\(G\) 是类型为 \(\mathscr R\) 的约化 \(S\)-群，并以 \(\Delta(\mathscr R)\) 表示由 \(\mathscr R\) 定义的 Dynkin 图。由 \SGARefLink{sga3:xxiv:3.5}{3.5}，有典范态射
\[
  A_S(\mathscr R)
  =
  \underline{\Aut}_{S\text{-}\mathrm{gr}}
    \bigl(\operatorname{\acute{E}p}_S(\mathscr R)\bigr)
  \longrightarrow
  \underline{\Aut}_{\mathrm{Dyn}}
    \bigl(\Delta(\mathscr R)_S\bigr).
\]
约化 \(S\)-群 \(G\) 对应于一个 \(A_S(\mathscr R)\) 下的主齐性丛（\SGARefLink{sga3:xxiv:1.17}{1.17}）
\[
  \underline{\operatorname{Isom}}_{S\text{-}\mathrm{gr}}
    \bigl(\operatorname{\acute{E}p}_S(\mathscr R),G\bigr),
\]
其相伴的 \(\underline{\Aut}_{\mathrm{Dyn}}(\Delta(\mathscr R)_S)\) 下的丛，对应于 \(\Delta(\mathscr R)_S\) 在 \(S\) 上的一个形式；这就是 \(\operatorname{Dyn}(G)\)。换言之，这个相伴丛正是
\[
  \underline{\operatorname{Isom}}_{\mathrm{Dyn}}
    \bigl(\Delta(\mathscr R)_S,\operatorname{Dyn}(G)\bigr).
\]
以这种形式表述后，证明是直接的。

\SGARefTarget{sga3:xxiv:3.8}
\medskip\noindent\textbf{3.8. Dynkin 概形与 Killing 偶。} ---
设 \(S\) 是概形，\(G\) 是约化 \(S\)-群，\(B\supset T\) 是 \(G\) 的 Killing 偶。有典范态射
\[
  i:\operatorname{Dyn}(G)
  \longrightarrow
  \underline{\operatorname{Hom}}_{S\text{-}\mathrm{gr}}
    (T,\mathbb G_{m,S})
\]
它将 \(\operatorname{Dyn}(G)\) 等同于“\(B\) 关于 \(T\) 的单根概形”。这个态射直接由钉扎情形下降而定义。还应注意，给定 \(T\) 与 \(i\)，就能恢复 \(B\)（“单根系与正根系之间的一一对应”）。

由上述 \(D=\operatorname{Dyn}(G)\) 的描述，可知 \(B_D\) 有一个关于 \(T_D\) 的典范根：这个根 \(\alpha_D\) 是 \(D\) 的恒等态射在 \(i(D)\) 下的像。由此得到一个典范的可逆 \(\mathcal O_D\)-模 \(\mathfrak g_D\)：
\[
  \mathfrak g_D
  =
  \bigl(\mathfrak g\otimes_{\mathcal O_S}\mathcal O_D\bigr)^{\alpha_D}.
\]
在钉扎情形，
\[
  D=\coprod_{\alpha\in\Delta}S_\alpha,
\]
其中每个 \(S_\alpha\) 都是 \(S\) 的副本，而 \(\mathfrak g_D\) 就是在每个 \(S_\alpha\) 上诱导 \(\mathfrak g_\alpha\) 的 \(\mathcal O_D\)-模，\(\alpha\in\Delta\)。

\SGARefTarget{sga3:xxiv:3.9}
\medskip\noindent\textbf{3.9. 拟钉扎. 拟钉扎群。} ---
若 \(G\) 是约化 \(S\)-群，\(G\) 的一个\emph{拟钉扎}是以下资料：
\begin{enumerate}
\item[(i)] \(G\) 的一个 Killing 偶 \(B\supset T\)；
\item[(ii)] 一个截面
\[
  X\in\Gamma\bigl(\operatorname{Dyn}(G),\mathfrak g_D\bigr)^\times.
\]
\end{enumerate}
若约化 \(S\)-群有拟钉扎，则称它\emph{可拟分裂}。\emph{拟钉扎群}是赋予拟钉扎的约化群。

设 \(B\supset T\) 是约化 \(S\)-群 \(G\) 的 Killing 偶。则 \(G\) 关于这个 Killing 偶可拟钉扎，当且仅当 \(\mathfrak g_D\) 有一个在每一点都非零的截面，即 \(\mathfrak g_D\) 所定义的 \(\operatorname{Pic}(\operatorname{Dyn}(G))\) 的元素为零。特别地，设 \(S\) 半局部。则 \(\operatorname{Dyn}(G)\) 也半局部，故 \(\operatorname{Pic}(\operatorname{Dyn}(G))=0\)。于是：

\SGARefTarget{sga3:xxiv:3.9.1}
\medskip\noindent\textbf{命题 3.9.1.} ---
{\itshape 设 \(S\) 是半局部概形，\(G\) 是约化 \(S\)-群。则 \(G\) 可拟分裂，当且仅当它有 Borel 子群。}
% 已译源行 1264--1654。

\footnote{编者注：已详细补写对第 XXII 讲 5.9.7 的引用。}
事实上，\(S\) 仿射，所以由\SGARefLink{sga3:xxii:cor:5.9.7}{第 XXII 讲 5.9.7} 的第一个断言，若 \(G\) 有 Borel 子群 \(B\)，则它也有 Killing 偶 \(B\supset T\)。此时，由于 \(\operatorname{Pic}(D)=0\)，模 \(\mathfrak g_D\) 有一个处处非零的截面 \(X\)。

仍设 \(A\) 是半局部环，\(S=\operatorname{Spec}(A)\)。对每个约化 \(S\)-群 \(G\)，态射
\[
  \operatorname{Bor}(G)\longrightarrow S
\]
是满的（因为对每个 \(s\in S\)，\(G_s\) 都有 Borel 子群）、光滑的、射影的（\SGARefLink{sga3:xxii:cor:5.8.3}{第 XXII 讲 5.8.3}）。因此，在作一次有限、满、平展的基扩张后，它便有截面。\footnote{编者注：原文引用了未出版的 EGA~IV 第 24 节。要点在于使用“Bertini 定理”。这里给出 O.~Gabber 向我们指出的论证的细节。设 \(X\to S\) 是满的光滑射影态射。将 \(S\) 换成一个连通分支后，可设 \(X/S\) 在每一点处的相对维数都是 \(d\)（参见 EGA~IV\textsubscript{4}，17.10.2）。还可设 \(X\) 是某个射影空间 \(\mathbb P^n_S=\mathbb P(E)\) 的闭子概形，其中 \(E\) 是秩为 \(n+1\) 的自由 \(A\)-模（参见 EGA~II，5.3.3）。设 \(s_1,\ldots,s_r\) 是 \(S\) 的闭点。由 Bertini 定理（例如参见 \textup{[Jou83]}，I，6.10），
\[
  P=\mathbb P(E^\ast)^d
\]
有一个纤维非空的开子集 \(U\)，使得对 \(U_{s_i}\) 的每个点 \(u=(f_1,\ldots,f_d)\)，\(X_{\kappa(u)}\) 与由 \(u\) 定义的 \(\mathbb P^n_{\kappa(u)}\) 中的 \(d\) 个超平面的交，在 \(\kappa(u)\) 上平展。缩小 \(U\) 后，可设 \(U\) 是仿射空间 \(\mathbb A^{nd}_S\) 中某个次数为 \(m>0\) 的多项式 \(Q\) 的零点集 \(V(Q)\) 的补集。对变量数作简单归纳便知，若 \(\lvert\kappa(s_i)\rvert>m\)，则 \(V(Q)\) 不可能包含 \(\mathbb A^{nd}_{\kappa(s_i)}\) 的全部有理点。由于态射
\[
  A\longrightarrow\prod_i\kappa(s_i)
\]
是满射，可找到一个次数为 \(m\) 的首一多项式 \(R\in A[X]\)，使其在 \(\kappa(s_i)[X]\) 中的像在 \(\kappa(s_i)\) 有限时不可约，而在 \(\kappa(s_i)\) 无限时有 \(m\) 个不同的根。置
\[
  A'=A[X]/(R),
  \qquad
  S'=\operatorname{Spec}(A').
\]
于是 \(S'\to S\) 平展、有限且满，并且每个闭点 \(s'_1,\ldots,s'_t\) 处剩余域的基数至少为 \(2m\)。因此，在 \(S'\) 的每个闭点上方，\(U\) 都有一个有理点 \(u_i\)，而由于
\[
  A'\longrightarrow\prod_i\kappa(s'_i)
\]
是满射\emph{（英译者注：法文印为 \(\kappa(s_i)\)；此处所指的 \(S'\) 的闭点应为 \(s'_i\)）}，这些点可提升为 \(P_{S'}\) 的一个截面 \(u\)。设 \(Z\) 是 \(X_{S'}\) 与由 \(u\) 定义的 \(\mathbb P^n_{S'}\) 中 \(d\) 个超平面的交，\(V\) 是 \(Z\) 的开子集，其点恰为 \(Z\) 在 \(S'\) 上平展的那些点。由 EGA~IV\textsubscript{3}，11.3.8，对每个 \(i\)，\(V\) 包含纤维 \(Z_{s'_i}\)。由于 \(\pi:Z\to S'\) 固有，闭集 \(\pi(Z-V)\) 为空，故 \(V=Z\)。所以 \(Z\to S'\) 满、平展且固有，因而有限；复合态射 \(Z\to S\) 也如此。这就给出 \(X\to S\) 所需的截面。}
由此得到：

\SGARefTarget{sga3:xxiv:3.9.2}
\medskip\noindent\textbf{推论 3.9.2.} ---
{\itshape 设 \(S\) 是半局部概形，\(G\) 是约化 \(S\)-群。存在平展、有限、满的态射 \(S'\to S\)，使得 \(G_{S'}\) 可拟分裂。}

\SGARefTarget{sga3:xxiv:3.9.3}
\medskip\noindent\textbf{注 3.9.3.} ---
在上述假设下，设 \(T\) 是 \(G\) 的极大环面（参见\SGARefLink{sga3:xiv:3.20}{第 XIV 讲 3.20}）。还可要求 \(G_{S'}\) 关于 \(T_{S'}\) 可拟分裂：只须将上述论证用于“包含 \(T\) 的 Borel 子群的概形”，它在 \(S\) 上有限且平展（\SGARefLink{sga3:xxii:corollary:5-5-5:item:ii}{第 XXII 讲 5.5.5(ii)}）。

\SGARefTarget{sga3:xxiv:3.10}
\medskip\noindent\textbf{3.10.} ---
设 \(E\)、\(E'\) 是约化 \(S\)-群 \(G\) 的两个拟钉扎。存在唯一的 \(G\) 的内自同构，将 \(E\) 变为 \(E'\)。事实上，立即可归结到分裂情形，而此时断言已经证明（\SGARefLink{sga3:xxiv:1.5}{1.5}；只须注意，对 \(G\) 的内自同构而言，保持一个钉扎等价于保持其底下的拟钉扎）。仿照第 1 节，得到约化 \(S\)-群 \(G\) 的一个拟钉扎定义了以下正合序列的一个分裂 \(h\)：
\[
\begin{tikzcd}[column sep=small]
1 \arrow[r]
&
\underline{\operatorname{ad}}(G) \arrow[r]
&
\underline{\Aut}_{S\text{-}\mathrm{gr.}}(G)
  \arrow[r,shift right=.7ex,"p"']
&
\underline{\Autext}(G)
  \arrow[l,shift right=.7ex,dotted,"h"']
  \arrow[r]
&
1
\end{tikzcd}
\]
其中 \(h\) 的像是 \(\underline{\Aut}_{S\text{-}\mathrm{gr}}(G)\) 中保持该拟钉扎的子群。

同样，若 \(G\)、\(G'\) 是两个拟钉扎 \(S\)-群，则自然定义了 \(\underline{\operatorname{Isom}}_{S\text{-}\mathrm{gr}}(G,G')\) 的子函子
\[
  \underline{\operatorname{Isom}}_{S\text{-}\mathrm{gr},\,q\text{-}\acute{e}p}
    (G,G')
\]
从 \(\underline{\operatorname{Isom}}_{S\text{-}\mathrm{gr}}(G,G')\) 到 \(\underline{\operatorname{Isomext}}(G,G')\) 的投影诱导同构
\[
  (*)\qquad
  \underline{\operatorname{Isom}}_{S\text{-}\mathrm{gr},\,q\text{-}\acute{e}p}
    (G,G')
  \xrightarrow{\ \sim\ }
  \underline{\operatorname{Isomext}}(G,G').
\]

\SGARefTarget{sga3:xxiv:3.11}
\medskip\noindent\textbf{定理 3.11.} ---
{\itshape 设 \(S\) 是概形，\(\mathscr R\) 是钉扎约化根数据，使得 \(\operatorname{\acute{E}p}_S(\mathscr R)\) 存在（参见第 XXV 讲），\(E\) 是它的自同构群。考虑以下三个范畴：
\begin{enumerate}
\item[\SGAInlineRefTarget{sga3:xxiv:3.11:item:i}(i)] \(S\) 的以 \(E\) 为群的主 Galois 覆盖所成范畴 \(\mathsf{Rev}\)（态射为 \(E\)-丛的同构）；
\item[\SGAInlineRefTarget{sga3:xxiv:3.11:item:ii}(ii)] 范畴 \(\mathsf{Redext}\)，其对象是类型为 \(\mathscr R\) 的约化 \(S\)-群（\SGARefLink{sga3:xxii:def:2.7}{第 XXII 讲 2.7}），从 \(G\) 到 \(G'\) 的态射为 \(\underline{\operatorname{Isomext}}(G,G')(S)\) 中的元素；
\item[(iii)] 类型为 \(\mathscr R\) 的拟钉扎约化 \(S\)-群所成范畴 \(\mathsf{Q\acute{e}p}\)（态射为保持拟钉扎的同构）。
\end{enumerate}
这三个范畴彼此等价。更精确地，有如下函子图，它在相差函子性同构的意义下交换：}
\[
\begin{tikzcd}[column sep=large,row sep=large]
\mathsf{Rev}
  \arrow[rr,"q\acute{e}p"]
&&
\mathsf{Q\acute{e}p}
  \arrow[dl,"i"]
\\
&
\mathsf{Redext}
  \arrow[ul,"\operatorname{rev}"]
&
\end{tikzcd}
\]

下面描述这三个函子，将该图交换性的验证留给读者。\footnote{编者注：在 \SGARefLink{sga3:xxiv:3.11.4}{3.11.4} 中，已简述如何验证函子 \(i\)、\(\operatorname{rev}\)、\(q\acute{e}p\) 完全忠实，以及复合 \(\operatorname{rev}\circ i\circ q\acute{e}p\) 同构于 \(\mathsf{Rev}\) 的恒等函子。}

\SGARefTarget{sga3:xxiv:3.11.1}
\medskip\noindent\textbf{3.11.1. 函子 \(i\)。} ---
这是显然的函子：
\[
  i(G)=G,
  \qquad
  i(f)=\text{\(f\) 在同构 3.10 \((*)\) 下的像}.
\]

\SGARefTarget{sga3:xxiv:3.11.2}
\medskip\noindent\textbf{3.11.2. 函子 \(q\acute{e}p\)。} ---
设 \(S'\) 是 \(S\) 的以 \(E\) 为群的主 Galois 覆盖。设
\[
  \bigl(G,T,M,\mathscr R,\Delta,
    (X_\alpha)_{\alpha\in\Delta}\bigr)
\]
是 \(G=\operatorname{\acute{E}p}_S(\mathscr R)\) 的钉扎，\((X'_\alpha)_{\alpha\in\Delta}\) 是
\[
  G'=G_{S'}=\operatorname{\acute{E}p}_{S'}(\mathscr R).
\]
的相应钉扎。由\SGARefLink{sga3:XXIII:5.5bis}{第 XXIII 讲 5.5 bis}，有唯一的群同态
\[
  \theta:E
  \longrightarrow
  \underline{\Aut}_{S\text{-}\mathrm{gr}}
    \bigl(\operatorname{\acute{E}p}_S(\mathscr R)\bigr)
  =A(\mathscr R)(S)
\]
使得对每个 \(h\in E\)，自同构 \(\theta(h)\) 在 \(T\) 上诱导 \(D_S(h)\)，而且
\[
  \Lie(\theta(h))(X_\alpha)=X_{h(\alpha)}
  \qquad(\alpha\in\Delta).
\]
结合 \(E\) 在 \(S'\) 上的作用，便得到 \(E\) 在 \(G'\) 上的一个作用，它与 \(E\) 在 \(S'\) 上的作用相容，并保持 \(G'\) 的典范拟钉扎。\footnote{编者注：上述讨论扩充了原文，以说明 \(E\) 在 \(G'\) 上的作用是将 \(E\) 在 \(\operatorname{\acute{E}p}_{S'}(\mathscr R)\) 上的自然作用与给定的 \(S'\) 上的作用结合而得到的。这个作用保持 \(G'\) 的典范 Killing 偶，也保持由
\[
  \Gamma\!\left(
    \Delta_{S'},
    \Lie(G'/S')_{\Delta_{S'}}
  \right),
\]
中元素 \(\widetilde X'\) 给出的拟钉扎；该元素在以 \(\alpha\) 为指标的 \(S'\) 副本上等于 \(X'_\alpha\)。事实上，\(E\) 在 \(\Delta_{S'}\) 上通过置换 \(S'\) 的各副本而作用，故对每个 \(h\in E\)，有 \(h(\widetilde X')=\widetilde X'\)。}
由于 \(S'\to S\) 是 \textup{(fpqc)} 拓扑的覆盖，它对于仿射态射的纤维范畴是有效下降态射。以
\[
  q\acute{e}p(S')
  =
  Q\text{-}\operatorname{\acute{E}p}_{S'/S}(\mathscr R)
\]
表示经 Galois 下降得到的拟钉扎 \(S\)-群。\footnote{编者注：例如参见 TDTE~I，第 22 页，\SGARefLink{sga3:xxiv:section:1}{例 1}。还可将 \(q\acute{e}p(S')\) 描述为 \(G=\operatorname{\acute{E}p}_S(\mathscr R)\) 由 \(E\)-主齐性丛 \(S'/S\) 的扭转，即 \(S'\times_SG\) 对 \(E\) 的右作用
\[
  (s',g)\cdot h=(s'h,h^{-1}(g)).
\]
所取的 \textup{(fpqc)} 层商。由于 \(G\) 在 \(S\) 上仿射，且 \(E\) 通过群自同构作用于 \(G\)，这个层由一个 \(S\)-群 \(G^\sharp\) 表示；它是 \(G\) 的“扭”形式，\(D^\sharp=\operatorname{Dyn}(G^\sharp)\) 是常 Dynkin 概形 \(\Delta_S\) 由主齐性丛 \(S'/S\) 的扭转。由于 \(E\) 正规化 \(B\) 与 \(T\)，同样得到 \(G^\sharp\) 的一个 Killing 偶 \((B^\sharp,T^\sharp)\)。另一方面，置 \(\mathfrak g=\Lie(G/S)\) 及 \(\mathfrak g^\sharp=\Lie(G^\sharp/S)\)。对每个 \(U\to S\)，\(\mathfrak g^\sharp\) 在 \(U\) 上的截面就是 \(E\)-等变的 \(S\)-态射
\[
  U\times_SS'\longrightarrow W(\mathfrak g).
\]
由于 \(D^\sharp\times_SS'\) 与赋予作用
\[
  (\alpha,s')\cdot h=(h^{-1}(\alpha),s'h),
\]
的 \(\Delta\times S'\) 是 \(E\)-同构的，态射 \((\alpha,s')\mapsto X_\alpha\) 就给出 \(\mathfrak g^\sharp\) 在 \(D^\sharp\) 上的一个截面，它是一个拟钉扎，即
\[
  \bigl(
    \mathfrak g^\sharp\otimes_{\mathcal O_S}\mathcal O_{D^\sharp}
  \bigr)^{D^\sharp}
\]
的处处非零截面（参见 \SGARefLink{sga3:xxiv:3.8}{3.8}）。}

\SGARefTarget{sga3:xxiv:3.11.3}
\medskip\noindent\textbf{3.11.3. 函子 \(\operatorname{rev}\)。} ---
设 \(G\) 是类型为 \(\mathscr R\) 的约化 \(S\)-群。记
\[
  \operatorname{rev}(G)
  =
  \underline{\operatorname{Isomext}}
    \bigl(\operatorname{\acute{E}p}_S(\mathscr R),G\bigr).
\]
这是 \(E_S\) 下的主齐性丛（参见 \SGARefLink{sga3:xxiv:1.10}{1.10} 及 \SGARefLink{sga3:xxiv:1.19}{1.19}），即 \(\mathsf{Rev}\) 的一个对象。

\SGARefTarget{sga3:xxiv:3.11.4}
\medskip\noindent\textbf{3.11.4.} ---
\footnote{编者注：这一小节号是补入的。}
由同构 \SGARefLink{sga3:xxiv:3.10}{3.10} \((*)\)，显然函子 \(i\) 完全忠实。另一方面，若 \(G_0,G,G'\) 是约化 \(S\)-群，则自然态射
\[
  \underline{\operatorname{Isom}}(G_0,G)
  \times_S
  \underline{\operatorname{Isom}}(G,G')
  \longrightarrow
  \underline{\operatorname{Isom}}(G_0,G')
\]
诱导态射
\[
  (\star)\qquad
  \underline{\operatorname{Isomext}}(G_0,G)
  \times_S
  \underline{\operatorname{Isomext}}(G,G')
  \longrightarrow
  \underline{\operatorname{Isomext}}(G_0,G'),
\]
因为对每个 \(S'\to S\)，若
\[
  g_0\in G_0(S'),\quad g\in G(S'),\quad
  u\in\underline{\operatorname{Isom}}(G_0,G)(S'),\quad
  v\in\underline{\operatorname{Isom}}(G,G')(S'),
\]
则有
\[
  v\circ\operatorname{int}(g)\circ u\circ\operatorname{int}(g_0)
  =
  v\circ u\circ
  \operatorname{int}\bigl(u^{-1}(g)g_0\bigr).
\]
取 \(G_0=\operatorname{\acute{E}p}_S(\mathscr R)\)，得到映射
\[
\begin{aligned}
  \underline{\operatorname{Isomext}}(G,G')(S)
  \longrightarrow
  \operatorname{Hom}_{\mathsf{Rev}}\bigl(
  &\underline{\operatorname{Isomext}}
    (\operatorname{\acute{E}p}_S(\mathscr R),G),\\
  &\underline{\operatorname{Isomext}}
    (\operatorname{\acute{E}p}_S(\mathscr R),G')
  \bigr),
\end{aligned}
\]
为证明这个映射是双射，利用 \textup{(fpqc)} 下降可设 \(G\) 与 \(G'\) 已钉扎；此时结论显然。

同样，若 \(S_1,S_2\) 是 \(\mathsf{Rev}\) 的两个对象，并以
\[
  S_i\times^E\operatorname{\acute{E}p}_S(\mathscr R)
\]
表示通过 \(E\)-主齐性丛 \(S_i\) 扭转 \(\operatorname{\acute{E}p}_S(\mathscr R)\) 而得到的 \(S\)-群 \(q\acute{e}p(S_i)\)（参见编者注 (21)），则有自然映射
\[
\begin{aligned}
  \operatorname{Hom}_{\mathsf{Rev}}(S_1,S_2)
  \longrightarrow
  \operatorname{Hom}_{\mathsf{Q\acute{e}p}}\bigl(
    &S_1\times^E\operatorname{\acute{E}p}_S(\mathscr R),\\
    &S_2\times^E\operatorname{\acute{E}p}_S(\mathscr R)
  \bigr).
\end{aligned}
\]
为证明它是双射，利用 \textup{(fpqc)} 下降可设 \(S_1\)、\(S_2\) 是平凡 \(E\)-丛；此时断言显然。

因此图中的三个函子都完全忠实。此外，对 \(\mathsf{Rev}\) 的每个对象 \(S'\)，有自然态射
\[
  S'
  \longrightarrow
  \underline{\operatorname{Isomext}}\!\left(
    \operatorname{\acute{E}p}_S(\mathscr R),
    S'\times^E\operatorname{\acute{E}p}_S(\mathscr R)
  \right),
\]
为证明它是同构，利用 \textup{(fpqc)} 下降可设 \(S'\) 是平凡 \(E\)-丛；此时断言显然。

于是有函子的同构
\[
  \operatorname{Id}_{\mathsf{Rev}}
  \longrightarrow
  \operatorname{rev}\circ i\circ q\acute{e}p.
\]
由于 \(\operatorname{rev}\) 完全忠实，对每个约化 \(S\)-群 \(H\)（不必拟分裂），得到关于 \(H\) 函子性的双射
\[
  \underline{\operatorname{Isomext}}\!\left(
    Q\text{-}\operatorname{\acute{E}p}_{\operatorname{rev}(H)/S}
      (\mathscr R),
    H
  \right)(S)
  \simeq
  \operatorname{Aut}_{\mathsf{Rev}}\bigl(\operatorname{rev}(H)\bigr),
\]
特别地，\(\operatorname{rev}(H)\) 的恒等态射对应于一个“典范”元素
\[
  u_0\in
  \underline{\operatorname{Isomext}}\!\left(
    Q\text{-}\operatorname{\acute{E}p}_{\operatorname{rev}(H)/S}
      (\mathscr R),
    H
  \right)(S).
\]
此外，这个集合的每个元素 \(u\) 对应于 \(\operatorname{rev}(H)\) 的一个自同构 \(\varphi_u\)，而 \(q\acute{e}p(\varphi_u)\) 是钉扎 \(S\)-群
\[
  Q\text{-}\operatorname{\acute{E}p}_{\operatorname{rev}(H)/S}
    (\mathscr R)
\]
的唯一自同构，使得
\[
  u_0\circ q\acute{e}p(\varphi_u)=u.
\]
这就证明了\SGARefLink{sga3:xxiv:3.11}{定理 3.11}。

我们展开讨论 \SGARefLink{sga3:xxiv:3.11}{3.11} 的一个推论。
\SGARefTarget{sga3:xxiv:3.12}
\medskip\noindent\textbf{推论 3.12.} ---
{\itshape 对每个约化 \(S\)-群 \(G\)，存在一个拟钉扎 \(S\)-群 \(G_{q\text{-}\acute{e}p}\) 和一个“外同构”}
\[
  u\in
  \underline{\operatorname{Isomext}}
    (G_{q\text{-}\acute{e}p},G)(S).
\]
{\itshape 有序对 \((G_{q\text{-}\acute{e}p},u)\) 在相差唯一同构的意义下唯一。}

事实上，可设 \(G\) 的类型恒为 \(\mathscr R\)，并取
\[
  G_{q\text{-}\acute{e}p}
  =
  Q\text{-}\operatorname{\acute{E}p}_{\operatorname{rev}(G)/S}
    (\mathscr R).
\]

\SGARefTarget{sga3:xxiv:3.12.1}
\medskip\noindent\textbf{3.12.1.} ---
注意，有序对 \((G_{q\text{-}\acute{e}p},u)\) 典范地定义了 \(S\)-概形（参见 \SGARefLink{sga3:xxiv:1.11}{1.11}）
\[
  Q=
  \underline{\operatorname{Isomint}}
    (G_{q\text{-}\acute{e}p},G),\footnote{编者注：由于有序对 \((G_{q\text{-}\acute{e}p},u)\) 在相差唯一同构的意义下唯一，我们写 \(\underline{\operatorname{Isomint}}(G_{q\text{-}\acute{e}p},G)\) 而不写 \(\underline{\operatorname{Isomint}}_u(G_{q\text{-}\acute{e}p},G)\)（参见 1.11）。}
\]
它是 \(\operatorname{ad}(G_{q\text{-}\acute{e}p})\) 下的主齐性丛，而给定它与给定 \(G_{q\text{-}\acute{e}p}\) 的内扭形式 \(G\) 是“等价”的。此外，\(Q\) 恰是“\(G\) 的拟钉扎概形”；这个定义解释了它作为 \(\operatorname{ad}(G)\) 下的（左）主齐性丛的结构（\SGARefLink{sga3:xxiv:3.10}{3.10}）；比较 \SGARefLink{sga3:xxiv:1.20}{1.20}。
% 已译源行 1655--2166。

\textbf{\SGAInlineRefTarget{sga3:xxiv:proposition:3-13}命题 3.13.} \emph{设 \(S\) 是概形，\(G\) 是半单伴随（分别为单连通）\(S\)-群，\(B \supset T\) 是 \(G\) 的 Killing 偶，\(\operatorname{Dyn}(G)\) 是 \(G\) 的 Dynkin \(S\)-概形。存在 \(S\)-群概形的典范同构}
\[
  T\xrightarrow{\sim}
  \prod_{\operatorname{Dyn}(G)/S}
  \mathbb G_{m,\operatorname{Dyn}(G)}.
\]
\emph{（回忆，参见第 II 讲第 1 节，右端按定义是将 \(S' \to S\) 对应到 \(G_{m}(\operatorname{Dyn}(G) \times_{S} S')\)，或等价地对应到 \(G_{m}(Dyn(G_{S'}))\) 的 \(S\)-函子。）}

\emph{第一个证明。} 为简便起见，考察伴随情形。考虑复合态射
\[
  T\longrightarrow\prod_{D/S}T_D
  \longrightarrow\prod_{D/S}\mathbb G_{m,D}.
\]
其中第一个态射是典范态射，第二个是 \(\prod_{D/S} \alpha_{D}\)（这里记 \(D = \operatorname{Dyn}(G)\)）。为验证这个态射是同构，可设 \(G\) 分裂；此时它无非是以 \(\alpha\)，\(\alpha \in \Delta\)，为分量的态射 \(T \to (G_{m, S})^{\Delta}\)，而后者是同构（\SGARefLink{sga3:xxii:rem:4.3.8}{第 XXII 讲 4.3.8}）。

\emph{第二个证明。} 由 \SGARefLink{sga3:xxiv:2.8}{2.8}、\SGARefLink{sga3:xxiv:3.5}{3.5} 和 \SGARefLink{sga3:xxiv:3.11}{3.11}，可设 \(G = Q-\acute{E}p_{S'/S}(R)\)，\(T\) 是典范极大环面。在 \(S'\) 上，由\SGARefLink{sga3:xxii:rem:4.3.8}{第 XXII 讲 4.3.8}，有由单根（分别为单余根）定义的同构 \(T_{S'} \xrightarrow{\sim} (G_{m, S'})^{\Delta}\)。群 \(E\) 通过置换 \(\Delta\) 作用于右端。但 \(S'/S\) 经 \(E \to \operatorname{Aut}(\Delta)\) 相伴的丛为 \(\operatorname{Dyn}(G)\)（\SGARefLink{sga3:xxiv:3.7}{3.7}），故立即得到结论。

利用附录的命题 \SGARefLink{sga3:xxiv:proposition:8-4}{8.4}，推出：
\textbf{\SGAInlineRefTarget{sga3:xxiv:corollary:3-14}推论 3.14.} \emph{在上述条件下，有}
\[
  H^1(S,T)\xrightarrow{\sim}
  H^1\bigl(\operatorname{Dyn}(G),\mathbb G_m\bigr)
  \xrightarrow{\sim}\operatorname{Pic}\bigl(\operatorname{Dyn}(G)\bigr).
\]
\emph{特别地，当 \(S\) 半局部时，\(H^{1}(S, T) = 0\)。}

\textbf{\SGAInlineRefTarget{sga3:xxiv:remark:3-15}注 3.15.} 设 \(S\) 是概形，\(G\)（分别为 \(G'\)）是约化 \(S\)-群，\(B \supset T\)（分别为 \(B' \supset T'\)）是 \(G\)（分别为 \(G'\)）的 Killing 偶，\(u \in Isomext(G, G')(S)\)。置（参见 \SGARefLink{sga3:xxiv:2.10}{2.10}）
\[
  P=\underline{\operatorname{Isomint}}_u(G,B,T;G',B',T').
\]
它是 \(T_{ad}\) 下的主齐性丛（作用为 \((f, t) \mapsto f \circ int(t)\)）。

另一方面，设 \(D = \operatorname{Dyn}(G) = \operatorname{Dyn}(G')\)（经由 \(u\) 等同，\SGARefLink{sga3:xxiv:3.5}{3.5}），并设 \(L = Isom_{D}(g_{D}, g'_{D})\) 是由可逆 \(\mathcal O_D\)-模
\[
  \operatorname{Hom}_{\mathcal O_D}(\mathfrak g_D,\mathfrak g'_D)
  =\mathfrak g'_D\otimes(\mathfrak g_D)^\vee .
\]
定义的 \(G_{m,D}\) 下的主齐性丛。

每个 \(f \in P(S')\) 都定义从 \(g_{D} \otimes_{O_{S}} O_{S'}\) 到 \(g'_{D} \otimes_{O_{S}} O_{S'}\) 的同构，因而有典范态射
\[ P \to \prod_{D/S} L. \]
这个态射是同构，并且与上述作用群的同构
\[
  T_{\mathrm{ad}}\xrightarrow{\sim}
  \prod_{D/S}\mathbb G_{m,D}.
\]
相容。事实上，只须在两个群都已钉扎的情形下验证，此时很容易。

特别地，在 \SGARefLink{sga3:xxiv:corollary:3-14}{3.14} 的同构
\[ H^{1}(S, T_{ad}) \xrightarrow{\sim} \operatorname{Pic}(\operatorname{Dyn}(G)) \]
下，丛 \(P\) 的类变为 \(c\ell(g'^{D}) - c\ell(g_{D})\)。因此丛 \(P\) 平凡，当且仅当 \(g'^{D}\) 与 \(g_{D}\) 同构。

若 \((G, B, T)\) 可拟分裂，例如取 \(G\) 为赋予典范 Killing 偶的群 \(G'_{q-\acute{e}p.}\)，则 \(P\) 的类的像无非是 \SGARefLink{sga3:xxiv:3.9}{3.9} 中定义的 \(G'\) 的拟分裂障碍。

\subsection*{\SGAInlineRefTarget{sga3:xxiv:subsection:3-16}3.16. 对称}
\SGAInlineRefTarget{sga3:xxiv:subsubsection:3-16-1}\subsubsection*{3.16.1.}
设 \(S\) 是概形，\(G\) 是约化 \(S\)-群，\(B \supset T\) 是 \(G\) 的 Killing 偶。记 \(D = \operatorname{Dyn}(G)\)。回忆（\SGARefLink{sga3:xxii:prop:5.9.2}{第 XXII 讲 5.9.2}），存在唯一的 \(G\) 的 Borel 子群 \(B^{-}\)，使得 \(B \cap B^{-} = T\)。若 \(X \in \Gamma(D,g_{D})^{\times}\) 定义了 \(G\) 关于 \((B, T)\) 的拟钉扎（\SGARefLink{sga3:xxiv:3.9}{3.9}），则 \(Y = -X^{-1} \in \Gamma(D,(g_{D})^{\vee})^{\times}\) 定义了 \(G\) 关于 \((B^{-}, T)\) 的拟钉扎；称它为\emph{相反拟钉扎}。

若 \(R\) 是钉扎约化根数据，\(E\) 是约化 \(S\)-群 \(G\) 的一个 \(R\)-钉扎，则按如下方式定义一个称为\emph{与 \(E\) 相反}的 \(R\)-钉扎 \(E^{-}\)：保留同一个极大环面 \(T\)，将同构 \(D_{S}(M) \xrightarrow{\sim} T\) 换成其相反，并用 \(Y_{\alpha} = -X^{-1}_{\alpha} \in\Gamma(S, g^{-\alpha})^{\times}\)，\(\alpha \in \Delta\)，来“钉扎”。\(E^{-}\) 的底层拟钉扎，正是与 \(E\) 的底层拟钉扎相反的拟钉扎。

\textbf{注。} 沿用\SGARefLink{sga3:xx:3.1}{第 XX 讲 3.1} 的记号，若置
\[
  w_\alpha(X_\alpha)
  =\exp(X_\alpha)\exp(-X_\alpha^{-1})\exp(X_\alpha).
\]
则有 \(w_{\alpha}(X_{\alpha}) = w_{-\alpha}(Y_{\alpha})\)（同上，3.1 (vi)）。

\textbf{\SGAInlineRefTarget{sga3:xxiv:proposition:3-16-2}命题 3.16.2.} \emph{设 \(S\) 是概形，\(G\) 是约化 \(S\)-群。}
\emph{\SGAInlineRefTarget{sga3:xxiv:proposition:3-16-2:item:i}(i) 设 \(T\) 是 \(G\) 的极大环面；存在唯一的}
\[
  i_T\in
  \left(
    \underline{\operatorname{Aut}}_{S\text{-}\mathrm{gr.}}(G,T)/T_{\mathrm{ad}}
  \right)(S)
  \subset
  \underline{\operatorname{Aut}}_{S\text{-}\mathrm{gr.}}(T).
\]
\emph{使得对 \(T\) 的每个截面 \(t\)，都有 \(i_{T}(t) = t^{-1}\)。}

\emph{\SGAInlineRefTarget{sga3:xxiv:proposition:3-16-2:item:ii}(ii) 设 \((B, T)\) 是 \(G\) 的 Killing 偶；存在唯一的截面}
\[
  w_{B,T}\in
  \bigl(\operatorname{Norm}_G(T)/T\bigr)(S)=W_G(T)(S).
\]
\emph{使得 \(int(w_{B, T})(B) = B^{-}\)（这里有显然的术语滥用）。}

\emph{\SGAInlineRefTarget{sga3:xxiv:proposition:3-16-2:item:iii}(iii) 设 \(Q = (B, T, X)\) 是 \(G\) 的拟钉扎，\(Q^{-} = (B^{-}, T, Y)\) 是相反拟钉扎；存在唯一的 \(G\) 的内自同构}
\[ n_{Q} \in Norm_{ad(G)}(T_{ad})(S) \subset \underline{\operatorname{Aut}}_{S\text{-}\mathrm{gr.}}(G) \]
\emph{使得 \(n_{Q}(Q) = Q^{-}\)，即 \(n_{Q}(T) = T\)、\(n_{Q}(B) = B^{-}\)、\(n_{Q}(X) = Y\)。}

\emph{\SGAInlineRefTarget{sga3:xxiv:proposition:3-16-2:item:iv}(iv) 设 \((R, E)\) 是 \(G\) 的钉扎，\((R, E^{-})\) 是相反钉扎；存在唯一的 \(G\) 的自同构}
\[
  u_E\in
  \underline{\operatorname{Aut}}_{S\text{-}\mathrm{gr.}}(G,T)
  \subset
  \underline{\operatorname{Aut}}_{S\text{-}\mathrm{gr.}}(G).
\]
\emph{使得 \(u_{E}(E) = E^{-}\)，即对 \(T\) 的每个截面 \(t\) 有 \(u_{E}(t) = t^{-1}\)，而对每个 \(\alpha \in \Delta\) 有 \(Ad(u_{E}) X_{\alpha} =Y_{\alpha}\)。}

\emph{证明。} \SGARefLink{sga3:xxiv:proposition:3-16-2:item:ii}{(ii)} 由\SGARefLink{sga3:xxii:corollary:5-5-5:item:ii}{第 XXII 讲 5.5.5(ii)} 推出；再由 \SGARefLink{sga3:xxiv:3.10}{3.10} 得 \SGARefLink{sga3:xxiv:proposition:3-16-2:item:iii}{(iii)}，由\SGARefLink{sga3:xxiii:theorem:4-1}{第 XXIII 讲 4.1} 得 \SGARefLink{sga3:xxiv:proposition:3-16-2:item:iv}{(iv)}。最后，为证明 \SGARefLink{sga3:xxiv:proposition:3-16-2:item:i}{(i)}，可设 \(G\) 已钉扎。存在性例如由 \SGARefLink{sga3:xxiv:proposition:3-16-2:item:iv}{(iv)} 推出；唯一性则来自如下事实：\(G\) 的一个在 \(T\) 上诱导恒等的自同构，由 \(T_{ad}\) 的一个截面给出（\SGARefLink{sga3:xxiv:2.11}{2.11}）。

\textbf{\SGAInlineRefTarget{sga3:xxiv:corollary:3-16-3}推论 3.16.3.} \emph{有}
\[
  i_T^2=w_{B,T}^2=n_Q^2=u_E^2=e.
\]
\emph{此外，若 \(G \neq e\)，则 \(i_{T}\)（分别为 \(u_{E}\)）\(\neq e\)；若 \(G\) 不是环面，则 \(w_{B, T}\)（分别为 \(n_{Q}\)）\(\neq e\)。}

\textbf{\SGAInlineRefTarget{sga3:xxiv:corollary:3-16-4}推论 3.16.4.} \emph{在 \SGARefLink{sga3:xxiv:proposition:3-16-2:item:iii}{(iii)}（分别为 \SGARefLink{sga3:xxiv:proposition:3-16-2:item:iv}{(iv)}）的情形下，\(n_{Q}\)（分别为 \(u_{E}\)）经由典范态射}
\[ Norm_{ad(G)}(T_{ad}) \to W_{ad(G)}(T_{ad}) \simeq W_{G}(T), \]
\emph{分别地}
\[
  \underline{\operatorname{Aut}}_{S\text{-}\mathrm{gr.}}(G,T)
  \longrightarrow
  \underline{\operatorname{Aut}}_{S\text{-}\mathrm{gr.}}(G,T)/T_{\mathrm{ad}}.
\]
\emph{投影到 \(w_{B, T}\)（分别为 \(i_{T}\)）。}

\textbf{\SGAInlineRefTarget{sga3:xxiv:corollary:3-16-5}推论 3.16.5.} \emph{上述定义与基变换相容，并且对于同构具有函子性（其意义是显然的）。}

\textbf{\SGAInlineRefTarget{sga3:xxiv:proposition:3-16-6}命题 3.16.6.} \emph{\SGAInlineRefTarget{sga3:xxiv:proposition:3-16-6:item:i}(i) 对概形 \(S\) 上的每个约化群 \(G\)，可以唯一地定义一个元素}
\[ s_{G} \in Autext(G)(S) \]
\emph{使得这一构造对于 \(G\) 的同构具有函子性，与基变换相容，并且只要 \(T\) 是约化 \(S\)-群 \(G\) 的极大环面，\(s_{G}\) 就是 \(i_{T}\) 在典范态射}
\[
  \underline{\operatorname{Aut}}_{S\text{-}\mathrm{gr.}}(G,T)/T_{\mathrm{ad}}
  \longrightarrow\underline{\underline{\operatorname{Autext}}}(G).
\]
\emph{下的像。}

\emph{\SGAInlineRefTarget{sga3:xxiv:proposition:3-16-6:item:ii}(ii) 有 \(s^{2}_{G} = e\)，并且 \(s_{G}\) 是 \(Autext(G)\) 的中心元素。}

\emph{\SGAInlineRefTarget{sga3:xxiv:proposition:3-16-6:item:iii}(iii) 在 \SGARefLink{sga3:xxiv:proposition:3-16-2:item:ii}{3.16.2(ii)} 的条件下，若将 \(\underline{\operatorname{Aut}}_{S\text{-}\mathrm{gr.}}(G, B, T) / T_{ad}\) 等同于 \(Autext(G)\)（\SGARefLink{sga3:xxiv:2.2}{2.2}），则有}
\[
  w_{B,T}i_T=i_Tw_{B,T}=s_G.
\]

\emph{\SGAInlineRefTarget{sga3:xxiv:proposition:3-16-6:item:iv}(iv) 在 \SGARefLink{sga3:xxiv:proposition:3-16-2:item:iv}{3.16.2(iv)} 的条件下，若将 \(\operatorname{Aut}_{S-gr.-\acute{e}p}(G, R, E)\) 等同于 \(Autext(G)\)（\SGARefLink{sga3:xxiv:1.3:item:iii}{1.3(iii)}），则有}
\[
  n_Eu_E=u_En_E=s_G.
\]

\emph{证明。} \SGARefLink{sga3:xxiv:proposition:3-16-6:item:i}{(i)} 用下降很容易证明。另一方面，由于 \(i_{T}\) 显然是 \(\underline{\operatorname{Aut}}_{S\text{-}\mathrm{gr.}}(T)\) 中平方为 \(e\) 的中心截面，\SGARefLink{sga3:xxiv:proposition:3-16-6:item:ii}{(ii)} 立即推出；\SGARefLink{sga3:xxiv:proposition:3-16-6:item:iii}{(iii)} 经由下降是 \SGARefLink{sga3:xxiv:proposition:3-16-6:item:iv}{(iv)} 的推论。最后，在 \SGARefLink{sga3:xxiv:proposition:3-16-6:item:iv}{(iv)} 的条件下，显然 \(n_{E} u_{E} = u_{E} n_{E}\)，且 \(G\) 的这个自同构保持钉扎；在所作的等同下，它因此等于自己在 \(Autext(G)\) 中的像；但 \(n_{E}\) 是内自同构，而 \(u_{E}\) 投影到 \(s_{G}\)。

\textbf{\SGAInlineRefTarget{sga3:xxiv:remark:3-16-7}注 3.16.7.} \SGARefLink{sga3:xxiv:proposition:3-16-6:item:i}{(i)} 由 \SGARefLink{sga3:xxiv:proposition:3-16-6:item:iii}{(iii)}，可在分类中的每种情形下显式确定 \(s_{G}\)：对每个不可约钉扎根数据，只须将关于原点的对称与 Weyl 群中的对称复合（即 Weyl 群中使 \(w_{0}(\Delta) = -\Delta\) 的元素 \(w_{0}\)）。得到如下结果：除 \(A_{n}\)（\(n \geqslant 2\)）、\(D_{n}\)（\(n\) 为奇数）和 \(E_6\) 外，均有 \(s_{G} = e\)；在这三种情形中，\(s_{G}\) 是唯一非平凡的“外自同构”。
\begin{enumerate}
\def\labelenumi{(\roman{enumi})}
\setcounter{enumi}{1}
\item 自同构 \(u_{E}\) 就是半单 Lie 代数理论中用来构造“紧实形式”的自同构。
\end{enumerate}

\textbf{\SGAInlineRefTarget{sga3:xxiv:remark:3-16-8}注 3.16.8.} 在 \SGARefLink{sga3:xxiv:subsubsection:3-16-1}{3.16.1} 中，我们在 \(G\) 的拟钉扎所成的 \(S\)-概形
\[
Q = Isomint(G_{q-\acute{e}p.}, G)
\]
上定义了一个对合（参见 \SGARefLink{sga3:xxiv:3.12.1}{3.12.1}）；通过从 \(G_{q-\acute{e}p.}\) 到 \(G\) 的结构转移，立即可见这一对合由 \(ad(G_{q-\acute{e}p.})(S)\) 中一个元素的作用给出：这个元素就是由 \(G_{q-\acute{e}p.}\) 的典范拟钉扎所定义（\SGARefLink{sga3:xxiv:proposition:3-16-2:item:iii}{3.16.2(iii)}）的元素 \(n_{0}\)。

同样，在 \(G\) 的 \(R\)-钉扎所成的 \(S\)-概形
\[ P = Isom_{S-gr.}(\acute{E}p_{S}(R), G) \]
上定义了一个对合（参见 \SGARefLink{sga3:xxiv:1.20}{1.20}）。仿照前述论证，可见这一对合由 \(\acute{E}p_{S}(R)\) 的自同构 \(u_{0}\) 的作用给出；它由 \(\acute{E}p_{S}(R)\) 的典范钉扎定义（\SGARefLink{sga3:xxiv:proposition:3-16-2:item:iv}{3.16.2(iv)}）。
% 已译源行 2167--2418。

\section*{\SGAInlineRefTarget{sga3:xxiv:section:4}4. 约化群及约化群下主丛的等平凡性}
\subsection*{\SGAInlineRefTarget{sga3:xxiv:subsection:4-1}4.1. 定义. 等平凡性定理}

\textbf{\SGAInlineRefTarget{sga3:xxiv:definition:4-1-1}定义 4.1.1.} \emph{设 \(S\) 是概形，\(G\) 是 \(S\)-群概形，\(P\) 是 \(G\) 下的主齐性丛。称 \(P\)}局部等平凡\emph{（分别为}半局部等平凡\emph{）}，如果对每个点 \(s \in S\)（分别为包含于仿射开集中的 \(S\) 的每个有限点集 \(F\)），都存在包含 \(s\)（分别为 \(F\)）的 \(S\) 的开集 \(U\)，以及有限满平展态射 \(S' \to U\)，使得 \(P_{S'}\) 平凡。

\textbf{\SGAInlineRefTarget{sga3:xxiv:definition:4-1-2}定义 4.1.2.} \emph{设 \(S\) 是概形，\(G\) 是约化 \(S\)-群。称 \(G\)}局部等平凡\emph{（分别为}半局部等平凡\emph{）}，如果对每个点 \(s \in S\)（分别为包含于仿射开集中的 \(S\) 的每个有限点集 \(F\)），都存在包含 \(s\)（分别为 \(F\)）的 \(S\) 的开集 \(U\)，以及有限满平展态射 \(S' \to U\)，使得 \(G_{S'}\) 可分裂。*

\textbf{\SGAInlineRefTarget{sga3:xxiv:remark:4-1-3}注 4.1.3.} (i) 按定义，\SGARefLink{sga3:xxiv:1.17}{1.17} 的范畴等价 \SGARefLink{sga3:xxiv:definition:4-1-1}{4.1.1} 保持局部（分别为半局部）等平凡性。
\begin{enumerate}
\def\labelenumi{(\roman{enumi})}
\setcounter{enumi}{1}
\item 在 \SGARefLink{sga3:xxiv:definition:4-1-1}{4.1.1} 的条件中补入：\(G\) 在 \(S\) 上局部有限呈示。则主丛 \(P\)（或约化群 \(G\)）局部等平凡（分别为半局部等平凡），当且仅当对每个 \(S' \to S\)，若 \(S'\) 局部（分别为半局部），则\(P_{S'}\) 等平凡（或 \(G_{S'}\) 等平凡），即存在有限满平展态射 \(S'' \to S'\)，使得 \(P_{S''}\) 平凡（或 \(G_{S''}\) 分裂）。
\item 对环面，\SGARefLink{sga3:xxiv:definition:4-1-2}{定义 4.1.2} 与\SGARefLink{sga3:ix:definition:1:1}{第 IX 讲 1.1} 的定义一致。
\end{enumerate}

\subsubsection*{4.1.4.}
回忆（\SGARefLink{sga3:xxii:section:4.3}{第 XXII 讲 4.3} 和 \SGARefLink{sga3:xxiv:lemma:6-2}{6.2}），对每个约化群 \(G\)，已引入群 \(rad(G)\)、\(corad(G)\) 及 \(d\acute{e}r(G)\)。群 \(rad(G)\)、\(corad(G)\) 都是环面，当 \(G\) 分裂时它们也分裂；此外，存在同源 \(rad(G) \to corad(G)\)。\(S\)-群 \(d\acute{e}r(G)\) 半单，有 \(G / d\acute{e}r(G) =corad(G)\)；因此对 \(G\) 下的每个主齐性丛 \(P\)，\(P / d\acute{e}r(G)\) 是 \(corad(G)\) 下的主齐性丛。于是有：

\textbf{\SGAInlineRefTarget{sga3:xxiv:theorem:4-1-5}定理 4.1.5.} \emph{设 \(S\) 是概形，\(G\) 是类型恒定的约化 \(S\)-群。}
\emph{\SGAInlineRefTarget{sga3:xxiv:theorem:4-1-5:item:i}(i) 下列条件等价：}
\emph{\SGAInlineRefTarget{sga3:xxiv:theorem:4-1-5:item:a}(a) \(G\) 局部（分别为半局部）等平凡。}
\emph{\SGAInlineRefTarget{sga3:xxiv:theorem:4-1-5:item:b}(b) 环面 \(rad(G)\) 如此。}
\emph{(b') 环面 \(corad(G)\) 如此。}
\emph{\SGAInlineRefTarget{sga3:xxiv:theorem:4-1-5:item:c}(c) 与 \(G\) 相伴的 \(S\) 的 Galois 覆盖（\SGARefLink{sga3:xxiv:3.11}{3.11}）如此。}
\emph{若 \(T\) 是 \(G\) 的极大环面，则这些条件还等价于}
\emph{\SGAInlineRefTarget{sga3:xxiv:theorem:4-1-5:item:d}(d) 环面 \(T\) 局部（分别为半局部）等平凡。}
\emph{\SGAInlineRefTarget{sga3:xxiv:theorem:4-1-5:item:ii}(ii) 设 \(P\) 是 \(G\) 下的主齐性丛；\(P\) 局部（分别为半局部）等平凡的充要条件，是 \(corad(G)\)-主丛 \(P / d\acute{e}r(G)\) 如此。}

\textbf{\SGAInlineRefTarget{sga3:xxiv:corollary:4-1-6}推论 4.1.6.} \emph{当 \(G\) 半单，或 \(S\) 局部 Noether 且正规（或更一般地，几何单枝）时，\SGARefLink{sga3:xxiv:theorem:4-1-5:item:i}{(i)} 的条件成立。当 \(G\) 局部（分别为半局部）等平凡时，\SGARefLink{sga3:xxiv:theorem:4-1-5:item:ii}{(ii)} 的条件成立。}

对 \SGARefLink{sga3:xxiv:theorem:4-1-5:item:i}{(i)}，第一个断言由 \SGARefLink{sga3:xxiv:theorem:4-1-5:item:b}{(b)} 显然成立，第二个断言由 \SGARefLink{sga3:xxiv:theorem:4-1-5:item:c}{(c)} 及\SGARefLink{sga3:x:corollary:5-14}{第 X 讲 5.14} 和 5.15 推出。对 \SGARefLink{sga3:xxiv:theorem:4-1-5:item:ii}{(ii)}，只须注意，由\SGARefLink{sga3:viii:corollary:4:5}{定理 90}，分裂环面下的主丛半局部等平凡。

\subsection*{\SGAInlineRefTarget{sga3:xxiv:subsection:4-2}4.2. 证明：半单情形}
为便于以后引用，先证明：
\textbf{\SGAInlineRefTarget{sga3:xxiv:proposition:4-2-1}命题 4.2.1.} \emph{设 \(S\) 是概形，\(G\) 是约化 \(S\)-群，\(T\) 是 \(G\) 的极大环面（分别地，\(B\) 是 Borel 子群，\(B \supset T\) 是 \(G\) 的 Killing 偶），\(P\) 是 \(G\) 下的主齐性丛，\(G'\) 是与 \(P\) 相伴的 \(G\) 的扭形式（经由态射 \(int : G \to \operatorname{Aut}_{S-gr.}(G)\)）。有典范同构}
\[
  P/\operatorname{Norm}_G(T)\xrightarrow{\sim}\operatorname{Tor}(G'),
  \qquad
  P/B\xrightarrow{\sim}\operatorname{Bor}(G'),
  \qquad
  P/T\xrightarrow{\sim}\operatorname{Kil}(G').
\]

按构造，\(G'\) 是 \(P \times_{S} G\) 对 \(G\) 的某个作用（\(((p, g')) g = (pg, g^{-1} g'g)\)）的商；因此有态射 \(P \times_{S} G \to G'\)，即一个态射
\[ P \to \operatorname{Hom}_{S}(G, G'), \]
立即可见它经由态射
\[
  f:P\longrightarrow\operatorname{Isom}_{S\text{-}\mathrm{gr.}}(G,G').
\]
分解（为验证此断言，可设 \(G = P\)，此时有 \(G' = G\)、\(f = int\)）。集合映射 \(p\mapsto f(p)(T)\) 定义一个态射
\[ P \to Tor(G'). \]
为验证这个态射确实诱导所述同构 \(P / Norm_{G}(T) \xrightarrow{\sim} Tor(G')\)，可再设 \(P = G\)，此时归结为\SGARefLink{sga3:xxii:cor:5.8.3:item:iii}{第 XXII 讲 5.8.3(iii)}。（实际上，应将该处陈述换成上述陈述。）对 \(\operatorname{Bor}\)、\(\operatorname{Kil}\) 作类似论证。

\textbf{\SGAInlineRefTarget{sga3:xxiv:proposition:4-2-2}命题 4.2.2.} \emph{设 \(S\) 是半局部概形，\(G\) 是类型恒定的半单 \(S\)-群。}
\emph{\SGAInlineRefTarget{sga3:xxiv:proposition:4-2-2:item:i}(i) \(G\) 等平凡。}
\emph{(ii) \(G\) 下的每个主丛都等平凡。}

先证 \SGARefLink{sga3:xxiv:proposition:4-2-2:item:i}{(i)}。作有限满平展基扩张后，由 \SGARefLink{sga3:xxiv:3.9.2}{3.9.2}，可设 \(G\) 拟分裂。但此时 \(G\) 按构造等平凡（\SGARefLink{sga3:xxiv:3.11}{3.11}，群 \(E\) 有限）。为证明 \SGARefLink{sga3:xxiv:3.11:item:ii}{(ii)}，由 \SGARefLink{sga3:xxiv:3.11:item:i}{(i)} 可设 \(G\) 分裂；于是归结为：

\textbf{\SGAInlineRefTarget{sga3:xxiv:lemma:4-2-3}引理 4.2.3.} \emph{设 \(S\) 是半局部概形。分裂约化群下的每个主丛都等平凡。}

事实上，沿用 \SGARefLink{sga3:xxiv:proposition:4-2-1}{4.2.1} 的记号，其中 \(B \supset T\) 表示分裂群 \(G\) 的典范 Killing 偶，由 \SGARefLink{sga3:xxiv:3.9.2}{3.9.2}，在作有限满平展基扩张后，\(S\)-概形 \(Kil(G')\) 有截面。因此可将 \(G\) 的结构群约化到 \(T\)；但 \(T\) 分裂，故 \(H^{1}(S, T) = 0\)（定理 90）。

\textbf{\SGAInlineRefTarget{sga3:xxiv:corollary:4-2-4}推论 4.2.4.} \emph{设 \(S\) 是概形，}
\[ 1 \to G \to G' \to G'' \to 1 \]
\emph{是 \(S\)-群概形的正合序列（对于 (fpqc)），其中 \(G\) 半单且类型恒定。设 \(P\) 是 \(G'\) 下的主齐性丛，假设相伴的层\footnote{编者注：我们将“丛”换成“层”，并将 \(G'\) 换成 \(G''\)。}\(P / G\) 可表（例如 \(G''\) 在 \(S\) 上仿射）。\(P\) 局部等平凡（分别为半局部等平凡）的充要条件，是 \(P / G\) 如此（作为 \(G''\) 下的丛）。}

若 \(P\) 平凡，则 \(P / G\) 也平凡，说明此条件必要。反之，设 \(S\) 局部（分别为半局部），\(P / G\) 等平凡，从而由 \(S\) 的一个有限满平展扩张 \(S'\) 平凡化。将基扩张到 \(S'\) 后，可把 \(P_{S'}\) 的结构群约化到 \(G_{S'}\)。但 \(S'\) 仍半局部，\(G_{S'}\) 半单且类型恒定，所以所得的丛等平凡（\SGARefLink{sga3:xxiv:proposition:4-2-2}{4.2.2}）。

\subsection*{\SGAInlineRefTarget{sga3:xxiv:subsection:4-3}4.3. 证明：一般情形}
首先注意，将 \SGARefLink{sga3:xxiv:corollary:4-2-4}{4.2.4} 用于正合序列
\[ 1 \to d\acute{e}r(G) \to G \to corad(G) \to 1. \]
立即得到 \SGARefLink{sga3:xxiv:theorem:4-1-5:item:ii}{4.1.5(ii)}。

因此来证明 \SGARefLink{sga3:xxiv:proposition:4-2-2:item:i}{(i)}。若 \(G\) 分裂，则 \(rad(G)\)、\(corad(G)\) 分裂，\(rev(G)\) 也是如此；所以 \SGARefLink{sga3:xxiv:theorem:4-1-5:item:a}{(a)} 蕴含 \SGARefLink{sga3:xxiv:theorem:4-1-5:item:b}{(b)}、(b') 和 \SGARefLink{sga3:xxiv:theorem:4-1-5:item:c}{(c)}。

\subsubsection*{4.3.1.}
有 \SGARefLink{sga3:xxiv:theorem:4-1-5:item:c}{(c)} \(\Rightarrow\) \SGARefLink{sga3:xxiv:theorem:4-1-5:item:a}{(a)}。设 \(R\) 是 \(G\) 的类型；有正合序列
\[ 1 \to ad(\acute{E}p_{S}(R)) \to A_{S}(R) \to E_{S} \to 1. \]
将 \SGARefLink{sga3:xxiv:corollary:4-2-4}{4.2.4} 用于丛 \(P = Isom_{S-gr.}(\acute{E}p_{S}(R), G)\)\linebreak 及其相伴丛 \(rev(G) = P /ad(\acute{E}p_{S}(R))\)，即得 \SGARefLink{sga3:xxiv:theorem:4-1-5:item:c}{(c)} \(\Rightarrow\) \SGARefLink{sga3:xxiv:theorem:4-1-5:item:a}{(a)}。

\subsubsection*{4.3.2.}
有 \SGARefLink{sga3:xxiv:theorem:4-1-5:item:b}{(b)} \(\Rightarrow\) \SGARefLink{sga3:xxiv:theorem:4-1-5:item:a}{(a)}。只须证明，若 \(rad(G)\) 分裂，则 \(G\) 半局部等平凡。现设 \(G_{0} =\acute{E}p_{S}(R)\)；考虑有序对 \((G',f)\) 的范畴，其中 \(G'\) 是 \(G_0\) 的形式，\(f\) 是从 \(rad(G_{0})\) 到 \(rad(G)\) 的同构。

已知（\SGARefLink{sga3:xxiv:2.16}{2.16}）这一范畴等价于某个 \(S\)-群 \(H\) 下主齐性丛所成范畴，该群是有限扭常群被半单群的扩张。只须证明 \(H\) 下的每个主丛都半局部等平凡；这立即由 \SGARefLink{sga3:xxiv:corollary:4-2-4}{4.2.4} 推出。

\subsubsection*{4.3.3.}
有 (b') \(\Rightarrow\) \SGARefLink{sga3:xxiv:theorem:4-1-5:item:a}{(a)}。可以像上面那样论证（而且出现的还是同一个群 \(H\)）。也可看到 \SGARefLink{sga3:xxiv:theorem:4-1-5:item:b}{(b)} 与 (b') 等价：与一个局部分裂环面同源 的环面也局部分裂（参见\SGARefLink{sga3:ix:proposition:2:11:item:iii:01}{第 IX 讲 2.11(iii)}）。

\subsubsection*{4.3.4.}
有 \SGARefLink{sga3:xxiv:theorem:4-1-5:item:d}{(d)} \(\Rightarrow\) \SGARefLink{sga3:xxiv:theorem:4-1-5:item:a}{(a)}。只须证明，一个有分裂极大环面的类型恒定的群半局部等平凡；这立即由 \SGARefLink{sga3:xxiv:2.14}{2.14} 推出。

\subsubsection*{4.3.5.}
有 \SGARefLink{sga3:xxiv:theorem:4-1-5:item:a}{(a)} \(\Rightarrow\) \SGARefLink{sga3:xxiv:theorem:4-1-5:item:d}{(d)}。只须证明分裂群的极大环面半局部等平凡。更精确地，有：

\textbf{\SGAInlineRefTarget{sga3:xxiv:lemma:4-3-6}引理 4.3.6.} \emph{设 \(S\) 是半局部概形，\(G\) 是约化 \(S\)-群，\(T_0\) 是 \(G\) 的分裂极大环面，\(W_{0}= Norm_{G}(T_{0}) / T_{0}\)（这是局部常 \(S\)-群；若 \(G\) 类型恒定，则由 \SGARefLink{sga3:xxiv:2.14}{2.14} 它是常群），\(T\) 是 \(G\) 的极大环面。}
\emph{存在态射 \(S' \to S\)，它是 \(W_0\) 下的主齐性丛（因而平展、有限且满；若 \(G\) 类型恒定，甚至是主 Galois 覆盖），使得 \(T_{S'}\) 通过 \(G(S')\) 的某个元素与 \((T_{0})_{S'}\) 共轭（因而特别地是分裂的）。}

事实上，已知 \(Transp_{G}(T_{0}, T)\) 是 \(Norm_{G}(T_{0})\) 下的主齐性丛（例如参见\SGARefLink{sga3:xi:corollary:5:4:bis}{第 XI 讲 5.4 bis}）。置 \(S' = Transp_{G}(T_{0}, T) / T_{0}\)；它是 \(W_0\) 下的主齐性丛。将基从 \(S\) 扩张到 \(S'\)，便可把 \(Transp_{G}(T_{0}, T)\) 的结构群约化到 \(T_0\)。现在 \(S'\) 半局部，\(T_0\) 分裂，所以 \(Transp_{G}(T_{0}, T)\) 在 \(S'\) 上有截面。\footnote{编者注：由“定理 90”。}\qed

\subsection*{\SGAInlineRefTarget{sga3:xxiv:subsection:4-4}4.4. 极大环面存在性的应用}
利用 Grothendieck 的极大环面存在性定理（\SGARefLink{sga3:xiv:3.20}{第 XIV 讲 3.20}），可以大大加强上述结果。先作如下陈述：

\textbf{\SGAInlineRefTarget{sga3:xxiv:theorem:4-4-1}定理 4.4.1.} \emph{设 \(S\) 是半局部概形，\(R\) 是钉扎约化根数据，\(W\) 是它的 Weyl 群，\(E\) 是它的自同构群。（回忆，\(E\) 自然作用于 \(W\)，而半直积 \(A = W \cdot E\) 等同于不带钉扎的 \(R\) 的自同构群，参见\SGARefLink{sga3:XXI:6.7.2}{第 XXI 讲 6.7.2}。）}
\emph{(i) \(\acute{E}p_{S}(R)\) 下的每个主齐性丛都由一个群为 \(W\) 的主 Galois 覆盖 \(S' \to S\) 平凡化。}
\emph{(ii) 设 \(G\) 是类型为 \(R\) 的约化 \(S\)-群；设 \(rev(G) = Isomext(\acute{E}p_{S}(R), G)\) 是相伴的群为 \(E\) 的 \(S\) 的 Galois 覆盖。设 \(W_0\) 是与 \(rev(G)\) 相伴的 \(W_S\) 的形式。存在态射 \(S' \to S\)，它是 \(W_0\) 下的主齐性丛，使得 \(G_{S'}\) 可拟分裂（即有 Borel 子群）。}
\emph{\SGAInlineRefTarget{sga3:xxiv:theorem:4-4-1:item:iii}(iii) 每个类型为 \(R\) 的约化 \(S\)-群都由一个群为 \(A\) 的主 Galois 覆盖 \(\tilde{S} \to S\) 分裂。}

先陈述：
\textbf{\SGAInlineRefTarget{sga3:xxiv:proposition:4-4-2}命题 4.4.2.} \emph{设 \(S\) 是概形，\(G\) 是约化 \(S\)-群，\(T\) 是 \(G\) 的极大环面，\(P\) 是 \(G\) 下的主齐性丛，\(G'\) 是与 \(P\) 相伴的 \(G\) 的扭形式（于是，参见 \SGARefLink{sga3:xxiv:proposition:4-2-1}{4.2.1}，有典范同构 \(P / Norm_{G}(T) \xrightarrow{\sim} Tor(G')\)）。设 \(T'\) 是 \(G'\) 的极大环面，定义复合态射}
\[
  S\longrightarrow\operatorname{Tor}(G')
  \xrightarrow{\sim}P/\operatorname{Norm}_G(T).
\]
\emph{考虑典范态射}
\[
  P\longrightarrow P/T\longrightarrow P/\operatorname{Norm}_G(T).
\]
\emph{并取它们沿上述态射的逆像：}
\[
\begin{tikzcd}[column sep=large,row sep=large]
P \arrow[r]
&
P/T \arrow[r]
&
P/\operatorname{Norm}_G(T)
\\
H \arrow[u] \arrow[r]
&
S' \arrow[u] \arrow[r]
&
S \arrow[u]
\end{tikzcd}
\]
\emph{则 \(S'\)（分别为 \(H\)）是 \(S\)（分别为 \(S'\)）上 \(W_{G}(T)\)（分别为 \(T_{S'}\)）下的主齐性丛。此外，若让 \(W_{G}(T)\) 以显然的方式作用于 \(T\)，则与 \(S'\) 相伴的丛同构于 \(T'\)。}

前两个断言显然，最后一个断言的证明与 \SGARefLink{sga3:xxiv:2.6}{2.6} 中相应断言的证明相同。

\textbf{\SGAInlineRefTarget{sga3:xxiv:corollary:4-4-3}推论 4.4.3.} \emph{设 \(S\) 是半局部概形，\(G\) 是约化 \(S\)-群，\(T\) 是 \(G\) 的极大环面。假设下列两个条件之一满足：}
\emph{\SGAInlineRefTarget{sga3:xxiv:corollary:4-4-3:item:i}(i) \(T\) 分裂。}
\emph{\SGAInlineRefTarget{sga3:xxiv:corollary:4-4-3:item:ii}(ii) \(T\) 包含在 \(G\) 的一个 Borel 子群中，并且 \(G\) 为伴随或单连通。}
\emph{再设 \(P\) 是 \(G\) 下的主齐性丛。存在一个 \(S\)-概形 \(S'\)，它是 \(W_{G}(T)\) 下的主齐性丛，使得 \(P_{S'}\) 平凡。}

事实上，若 \(G'\) 是与 \(P\) 相伴的 \(G\) 的形式，则 \(G'\) 有极大环面 \(T'\)（\SGARefLink{sga3:xiv:3.20}{第 XIV 讲 3.20}）。沿用前一命题的记号，可见 \(H^{1}(S', T_{S'}) = 0\)（对 \SGARefLink{sga3:xxiv:corollary:4-4-3:item:i}{(i)} 用定理 90，对 \SGARefLink{sga3:xxiv:corollary:4-4-3:item:ii}{(ii)} 用 \SGARefLink{sga3:xxiv:corollary:3-14}{3.14}）。态射 \(H \to S'\) 有截面，故 \(P_{S'}\) 在 \(S'\) 上也有截面。\qed

现在证明定理。断言 \SGARefLink{sga3:xxiv:corollary:4-4-3:item:i}{(i)} 是前一推论的特例（取 \(G =\acute{E}p_{S}(R)\)，赋予其典范分裂环面）。来证 \SGARefLink{sga3:xxiv:corollary:4-4-3:item:ii}{(ii)}。已知（\SGARefLink{sga3:xxiv:3.12}{3.12}）\(G\) 是
\[ G_{0} = Q-\acute{E}p_{rev(G)/S}(R). \]
的内扭形式。

若 \(T_0\) 是 \(G_0\) 的典范极大环面，则 \(W_{G_{0}}(T_{0})\) 就是陈述中的群 \(W_0\)。\(G_0\) 的形式 \(G\) 对应于 \(ad(G_{0})\) 下的主齐性丛 \(P\)（\(P = Isomint(G_{0}, G)\)）。群 \(W_{ad(G_{0})}(T^{ad}_{0})\) 典范同构于 \(W_0\)，而将 \SGARefLink{sga3:xxiv:corollary:4-4-3}{4.4.3} 用于 \((ad(G_{0}), T^{ad}_{0}, P)\)，即得所求结果；\SGARefLink{sga3:xxiv:corollary:4-4-3}{4.4.3} 的假设 \SGARefLink{sga3:xxiv:corollary:4-4-3:item:ii}{(ii)} 确实满足。最后来证 \SGARefLink{sga3:xxiv:theorem:4-4-1:item:iii}{(iii)}。沿用 \SGARefLink{sga3:xxiv:corollary:4-4-3:item:ii}{(ii)} 的记号，有图
\[
\begin{tikzcd}[column sep=large,row sep=large]
&
\operatorname{rev}(G) \arrow[d,"E_S"']
\\
S'
  \arrow[r,"W_0"']
&
S
\end{tikzcd}
\]
已知 \(G_{S'}\) 同构于 \((G_{0})_{S'}\)，而 \((G_{0})_{rev(G)}\) 可分裂。若置 \(\tilde{S} = S' \times_{S} rev(G)\)，则 \(G_{\tilde{S}}\) 可分裂，只须验证 \(\tilde{S}\) 确为 \(S\) 的群为 \(A\) 的主 Galois 覆盖；这由下述更一般的引理推出（它自然适用于任意景）：
% 已译源行 2419--2723。

\textbf{\SGAInlineRefTarget{sga3:xxiv:lemma:4-4-4}引理 4.4.4.} \emph{设 \(S\) 是概形，\(G\)、\(H\) 是两个 \(S\)-群概形，\(G \to \operatorname{Aut}_{S-gr.}(H)\) 是 \(G\) 在 \(H\) 上的作用，\(E\) 是 \(G\)-主齐性丛，\(F\) 是 \(H_0\)-主齐性丛，其中 \(H_0\) 是与 \(E\) 相伴的 \(H\) 的形式。则 \(E \times_{S} F\) 自然带有半直积 \(H \cdot G\) 下的主齐性丛结构。}

以 \((e, g) \mapsto eg\)（分别为 \((f, u) \mapsto f u\)）表示 \(G\) 在 \(E\) 上（分别为 \(H_0\) 在 \(F\) 上）的（右）作用。以
\[
  p:E\times_S H\longrightarrow H_0.
\]
表示典范投影（按定义，\(H_0\) 是 \(E \times_{S} H\) 对 \(G\) 的作用 \((e, h) g = (eg, g^{-1} h g)\) 的商）。考虑态射
\[
  r:E\times_S F\times_S H\times_S G\longrightarrow E\times_S F.
\]
其集合层面的定义为
\[
  r(e,f,h,g)=\bigl(eg,f\cdot p(e,h)\bigr).
\]
态射 \(r\) 确实定义半直积 \(H \cdot G\) 在 \(E \times_{S} F\) 上的作用。事实上，在集合层面有
\[
  r\bigl(r(e,f,h,g),h',g'\bigr)
  =\bigl(egg',f\,p(e,h)\,p(eg,h')\bigr).
\]
但
\[
  p(e,h)p(eg,h')
  =p(e,h)p(e,gh'g^{-1})
  =p(e,hgh'g^{-1}).
\]
因此
\[
  r\bigl(r(e,f,h,g),h',g'\bigr)
  =r(e,f,hgh'g^{-1},gg').
\]
这正是所要证明的。

现在为证明这个作用律确实给出主齐性丛结构，可设 \(E\)、\(F\) 平凡，此时立即可见 \(E \times_{S} F\) 也是平凡丛。

\subsection*{\SGAInlineRefTarget{sga3:xxiv:subsection:4-5}4.5. 半单群的极大环面的等平凡性}
\footnote{\textit{编者注。} 本小节 \SGARefLink{sga3:xxiv:subsection:4-5}{4.5} 为增补。}

这里指出以下结果，它已隐含在 \SGARefLink{sga3:xxiv:2.6}{2.6} 中（参见编者注 (13)），并且在\SGARefLink{sga3:ix:comments:1:2}{第 IX 讲 1.2} 末尾提及。回忆（\SGARefLink{sga3:XXIII:5.11}{XXIII 5.11}），对每个概形 \(S\) 以及每个约化根数据 \(R\)，以 \(\acute{E}p_{S}(R)\) 表示类型为 \(R\) 的唯一钉扎 \(S\)-群（由第 XXV 讲存在），以 \(T_{S}(R)\) 表示它的典范极大环面。

\textbf{\SGAInlineRefTarget{sga3:xxiv:proposition:4-5-1}命题 4.5.1.} \emph{设 \(G\) 是类型恒为 \(R\) 的半单 \(S\)-群；假设 \(G\) 有极大环面 \(T\)（由 \SGARefLink{sga3:xiv:3.20}{XIV 3.20}，\(S\) 半局部时即为如此）。}
\emph{\SGAInlineRefTarget{sga3:xxiv:proposition:4-5-1:item:i}(i) 则 \(P = Isom_{S-gr.}(\acute{E}p_{S}(R), T_{S}(R); G, T) / T_{S}(R)_{ad}\) 是有限群 \(\operatorname{Aut}(R)\) 下的主 \(S\)-丛，且 \(T_P\) 是 \(G_P\) 的分裂极大环面。因此 \(T\) 等平凡。}
\emph{\SGAInlineRefTarget{sga3:xxiv:proposition:4-5-1:item:ii}(ii) 此外，对包含于一个仿射开集中的 \(P\) 的每个有限点集 \(F\)，存在包含 \(F\) 的 \(P\) 的开集 \(U\)，使得 \(G_U\) 分裂。}

\emph{证明。} 正合序列
\[
  1\longrightarrow T_S(\mathscr R)_{\mathrm{ad}}
  \longrightarrow
  \underline{\operatorname{Aut}}_{S\text{-}\mathrm{gr.}}
  \bigl(\operatorname{\acute{E}p}_S(\mathscr R),T_S(\mathscr R)\bigr)
  \longrightarrow
  \underline{\operatorname{Aut}}_{S\text{-}\mathrm{gr.}}\bigl(T_S(\mathscr R)\bigr).
\]
对每个 \(S' \to S\) 诱导丛的态射 \(P_{S'} \to Isom_{S'-gr.}(T_{S'}(R), T_{S'})\)，所以若 \(P_{S'}\) 有截面，则 \(T_{S'}(R)\) 与 \(T_{S'}\) 同构。特别地，对 \(S' = P\) 即为如此。

再注意，\(\operatorname{Aut}_{S-gr.}(\acute{E}p_{S}(R), T_{S}(R)) / T_{S}(R)_{ad}\) 无非是常 \(S\)-群 \(\operatorname{Aut}(R)_{S}\)，而 \(P\) 等同于 \(\operatorname{Aut}(R)_{S}\)-主齐性丛 \(Isom(R, \tilde{R}(G, T))\)，其中 \(\tilde{R}(G, T)\) 是与 \((G, T)\) 相伴的扭根数据（参见 \SGARefLink{sga3:xxii:prop:1.10}{XXII 1.10}）。

于是 \SGARefLink{sga3:xxiv:proposition:4-5-1:item:ii}{(ii)} 由 \SGARefLink{sga3:xxiv:proposition:4-5-1:item:i}{(i)} 与 \SGARefLink{sga3:xxiv:2.14}{2.14} 推出，只须注意，在包含 \(F\) 的充分小的仿射开集 \(U\) 上，各 \(g_{\alpha}\)（\(\alpha\) 遍历 \(R\) 的根）都是秩为 1 的自由 \(\mathcal O_U\)-模。

\section*{\SGAInlineRefTarget{sga3:xxiv:section:5}5. 伴随或单连通群的典范分解}
本节利用第 1 节的结果，将伴随（分别为单连通）群的经典分解推广到概形情形。为避免叙述过于繁重，只略述证明，将细节留给读者；事实上，所用的总是主丛理论中完全标准的论证：结构群约化、扭转等。

\subsection*{\SGAInlineRefTarget{sga3:xxiv:subsection:5-1}5.1.}
回忆（\SGARefLink{sga3:XXI:7.4}{第 XXI 讲 7.4}），Dynkin 图是它的各连通分支的不交并，而各连通分支也是 Dynkin 图。此外，与根数据对应的每个非空连通 Dynkin 图都同构于\SGARefLink{sga3:XXI:7.4.6}{第 XXI 讲 7.4.6} 中列出的一个标准图（\(A_{n}\)、\(B_{n}\)、\ldots、\(G_2\)）。下文只考虑连通分支为上述类型之一的 Dynkin 图。设 \(T\) 为这些标准图组成的集合。对每个 Dynkin 图 \(D\)，设 \(n(t) = n_{D}(t)\) 是 \(D\) 中同构于 \(t\) 的连通分支数，其中 \(t \in T\)。\(D\) 的\emph{类型}按定义为 \(\sum_{t \in T} n_{D}(t) \cdot t\)。

类型为 \(t\) 的 Dynkin 图称为\emph{类型 \(t\) 的单图}，类型为 \(n \cdot t\) 的 Dynkin 图称为\emph{类型 \(t\) 的等型图}。设 \(D_0\) 是 \(D\) 的连通分支集，
\[ a : D_{0} \to T \]
是显然的映射。\(D\) 的类型无非是 \(\sum_{x \in D_{0}} a(x)\)。

\subsection*{\SGAInlineRefTarget{sga3:xxiv:subsection:5-2}5.2.}
设 \(S\) 是概形，\(D\) 是一个 Dynkin \(S\)-概形（满足上述限制条件）。将态射对 \(L \Rightarrow D\) 的余核记为 \(D_0\)，其中 \(L =\)\footnote{编者注：仿照 \SGARefLink{sga3:xxiv:3.2}{3.2}，将 \(R\) 换成了 \(L\)。}\(D\) 的键概形。它是 \(D\) 的“连通分支概形”（当 \(D\) 为常概形时，其存在性显然；一般情形由下降得到）；它是有限扭常 \(S\)-概形。有典范态射
\[ a : D_{0} \to T_{S}. \]
对 \(t \in T\)，置 \(a^{-1}(t) = D_{0, t}\)；这是 \(D_0\) 的子概形，它在 \(D\) 中的逆像记为 \(D_{t}\)，称为 Dynkin 概形 \(D\) 的\emph{类型 \(t\) 的等型分量}。每个 \(D_{t}\) 都是 \(D\) 的子概形，且有
\[
  D=\coprod_{t\in T}D_t.
\]
注意，若 \(D\) 在 \(s\) 处的类型为 \(\sum_{t} n(s, t) \cdot t\)，则 \(D_{0, t}\) 在 \(s \in S\) 处的次数为 \(n(s, t)\)。

\subsection*{\SGAInlineRefTarget{sga3:xxiv:subsection:5-3}5.3.}
下文只考虑半单伴随（分别为单连通）群。为简化用语，我们对伴随群陈述结果；若将各处“伴随”换成“单连通”，所有陈述仍然成立。

回忆，伴随约化根数据在相差同构的意义下由其 Dynkin 图的类型确定。因此，若伴随根数据 \(R\) 的 Dynkin 图具有该类型，则称其\emph{具有类型 \(\sum n(t) \cdot t\)}；对半单伴随群 \(G\)，若其类型由类型为 \(\sum n(t) \cdot t\) 的伴随根数据给出，也作同样定义。若 \(R\) 或 \(G\) 的类型为 \(t\)（分别为 \(n \cdot t\)，\(n > 0\)），则称它\emph{为类型 \(t\) 的单的}（分别为\emph{类型 \(t\) 的等型的}）。

若 \(G\) 是半单伴随群，则按 \SGARefLink{sga3:xxiv:subsection:5-2}{5.2} 中定义的意义使用记号 \(Dyn_{0}(G)\) 及 \(Dyn_{0, t}(G)\)。

\subsection*{\SGAInlineRefTarget{sga3:xxiv:subsection:5-4}5.4.}
设 \(t_i\)，\(i=1,2,\ldots,n\)，是 \(T\) 的互不相同的元素，\(R_i\) 是类型 \(t_{i}\) 的等型钉扎伴随根数据。考虑乘积 \(R = R_{1} \times \cdot\cdot\cdot \times R_{n}\)（\SGARefLink{sga3:XXI:6.4.1}{第 XXI 讲 6.4.1}）。设 \(S\) 是使各 \(\acute{E}p_{S}(R_{i})\) 都存在的概形（参见第 XXV 讲）。立即可验证，存在典范同构
\[
  \tag{*}\label{sga3:xxiv:subsection:5-4:equation:star}
  \operatorname{\acute{E}p}_S(\mathscr R)
  =
  \operatorname{\acute{E}p}_S(\mathscr R_1)
  \times_S\cdots\times_S
  \operatorname{\acute{E}p}_S(\mathscr R_n).
\]
% 已译源行 2724--2897。

\par\medskip\noindent\SGAInlineRefTarget{sga3:xxiv:proposition:5-5}\textbf{命题 5.5.---}\enspace
\emph{设 \(S\) 是概形，\(G\) 是伴随半单 \(S\)-群。存在唯一分解}
\[
  G\simeq\prod_{t\in\mathrm T}G_t,
\]
\emph{其中 \(G_t\) 是类型 \(t\) 的等型伴随半单 \(S\)-群。此外，这一分解诱导同构}
\[
  \operatorname{Dyn}_t(G)\simeq\operatorname{Dyn}(G_t),
  \qquad
  \underline{\operatorname{Aut}}_{S\text{-}\mathrm{gr.}}(G)
  \simeq
  \prod_{t\in\mathrm T}
  \underline{\operatorname{Aut}}_{S\text{-}\mathrm{gr.}}(G_t).
\]

当 \(G\) 分裂时，已在上面证明。一般情形可设 \(G\) 的类型恒为 \(\mathscr R\)。利用 \(A_S(\mathscr R)\) 的上述分解及 \SGARefLink{sga3:xxiv:1.17}{1.17}，便得到所求的 \(G\) 的分解。其余断言由下降得到。

\par\medskip\noindent\SGAInlineRefTarget{sga3:xxiv:remark:5-6}\textbf{注 5.6.---}\enspace
更一般地，若 \(G\)、\(H\) 是两个伴随半单群，则以下交换图中有典范同构：
\[
\begin{tikzcd}[column sep=huge,row sep=large]
\underline{\operatorname{Isom}}_{S\text{-}\mathrm{gr.}}(G,H)
  \arrow[r,"\sim"] \arrow[d]
&
\displaystyle\prod_{t\in\mathrm T}
\underline{\operatorname{Isom}}_{S\text{-}\mathrm{gr.}}(G_t,H_t)
  \arrow[d]
\\
\underline{\operatorname{Isomext}}(G,H)
  \arrow[r,"\sim"] \arrow[d,"\wr"']
&
\displaystyle\prod_{t\in\mathrm T}
\underline{\operatorname{Isomext}}(G_t,H_t)
  \arrow[d,"\wr"]
\\
\underline{\operatorname{Isom}}_{\mathrm{Dyn}}
  \bigl(\operatorname{Dyn}(G),\operatorname{Dyn}(H)\bigr)
  \arrow[r,"\sim"]
&
\displaystyle\prod_{t\in\mathrm T}
\underline{\operatorname{Isom}}_{\mathrm{Dyn}}
  \bigl(\operatorname{Dyn}(G_t),\operatorname{Dyn}(H_t)\bigr).
\end{tikzcd}
\]

\par\medskip\noindent\SGAInlineRefTarget{sga3:xxiv:remark:5-7}\textbf{注 5.7.---}\enspace
\(G_t\) 有如下内在刻画，这里只陈述而不证明：它是 \(G\) 的最大的、正规（实际上为特征）且类型 \(t\) 等型的约化子群。

上述命题将伴随半单群的研究归结到等型伴随半单群的研究。现在进行这一研究。

\subsection*{\SGAInlineRefTarget{sga3:xxiv:subsection:5-8}5.8.}
设 \(\mathscr R\) 是类型 \(t\) 的钉扎约化伴随单根数据，\(I\) 是非空有限集，\(\mathscr R^I\) 是 \(\mathscr R\) 的各副本 \(\mathscr R_i\)，\(i\in I\)，的乘积根数据。记 \(\mathscr R\) 的自同构群为 \(E\)，将它等同于其 Dynkin 图 \(D\) 的自同构群。\(\mathscr R^I\) 的 Dynkin 图等同于 \(D^I\)，因而有正合序列
\[
  1\longrightarrow\operatorname{Aut}(D)^I
  \longrightarrow\operatorname{Aut}(D^I)
  \longrightarrow\operatorname{Aut}(I)
  \longrightarrow1,
\]
其中 \(\operatorname{Aut}(I)\) 是 \(I\) 的置换群。因此典范同构
\[
  \operatorname{\acute Ep}_S(\mathscr R)^I
  \simeq
  \operatorname{\acute Ep}_S(\mathscr R^I)
\]
诱导正合序列
\[
  1\longrightarrow A_S(\mathscr R)^I
  \longrightarrow A_S(\mathscr R^I)
  \longrightarrow\underline{\operatorname{Aut}}(I_S)
  \longrightarrow1.
\]
最后的群是与 \(\operatorname{Aut}(I)\) 相伴的常 \(S\)-群。还应注意，\(I\) 典范地等同于 \(D^I\) 的连通分支集。

设 \(G\) 是类型为 \(\mathscr R^I\) 的半单 \(S\)-群，从而定义一个 \(A_S(\mathscr R^I)\) 下的主齐性丛 \(P\)（参见 \SGARefLink{sga3:xxiv:1.17}{1.17}）。\(\operatorname{Dyn}(G)\)（\SGARefLink{sga3:xxiv:3.7}{3.7}）与 \(\operatorname{Dyn}_0(G)\)（\SGARefLink{sga3:xxiv:subsection:5-2}{5.2}）经由下降的定义表明，\(P\) 经
\[
  A_S(\mathscr R^I)\longrightarrow
  \underline{\operatorname{Aut}}(I_S)
\]
相伴的丛恰为 \(\underline{\operatorname{Isom}}_S(I_S,\operatorname{Dyn}_0(G))\)，对应于 \(I_S\) 的形式 \(\operatorname{Dyn}_0(G)\)。再次使用范畴等价 \SGARefLink{sga3:xxiv:1.17}{1.17} 及上述正合序列，形式地推出：存在类型为 \(\mathscr R\) 的约化 \(\operatorname{Dyn}_0(G)\)-群 \(G_0\)，以及 \(S\)-同构
\[
  G\simeq\prod_{\operatorname{Dyn}_0(G)/S}G_0.
\]

\par\medskip\noindent\SGAInlineRefTarget{sga3:xxiv:proposition:5-9}\textbf{命题 5.9.---}\enspace
\emph{设 \(S\) 是概形，\(G\) 是类型 \(t\) 的等型伴随半单 \(S\)-群。存在类型 \(t\) 的单的伴随半单 \(\operatorname{Dyn}_0(G)\)-群 \(G_0\)，以及唯一的 \(S\)-同构}
\[
  G\simeq\prod_{\operatorname{Dyn}_0(G)/S}G_0.
\]
\emph{此外，这一同构诱导正合序列}
\[
  1\longrightarrow
  \prod_{\operatorname{Dyn}_0(G)/S}
    \underline{\operatorname{Aut}}_{\mathrm{gr.}}(G_0)
  \longrightarrow
  \underline{\operatorname{Aut}}_{S\text{-}\mathrm{gr.}}(G)
  \longrightarrow
  \underline{\operatorname{Aut}}_S\bigl(\operatorname{Dyn}_0(G)\bigr)
  \longrightarrow1.
\]

显然可设 \(G\) 的类型恒为 \(n\,t\)。第一个断言已经证明（唯一性也显然）。第二个断言经由下降从分裂情形得到。

结合 \SGARefLink{sga3:xxiv:proposition:5-5}{5.5} 与 \SGARefLink{sga3:xxiv:proposition:5-9}{5.9}，得到：
\par\medskip\noindent\SGAInlineRefTarget{sga3:xxiv:proposition:5-10}\textbf{命题 5.10.---}\enspace
\emph{设 \(S\) 是概形，\(G\) 是伴随半单 \(S\)-群，并置 \(\mathfrak D=\operatorname{Dyn}(G)\)。}
\emph{(i) 有典范分解}
\[
  G\simeq
  \prod_{t\in\mathrm T}
  \prod_{\mathfrak D_{0,t}/S}G_{0,t}
  \simeq
  \prod_{\operatorname{Dyn}_0(G)/S}G_0,
\]
\emph{其中每个 \(G_{0,t}\) 都是类型 \(t\) 的单的伴随 \(\operatorname{Dyn}_{0,t}(G)\)-群；等价地，\(G_0\) 是伴随半单 \(\operatorname{Dyn}_0(G)\)-群，其在 \(x\in\operatorname{Dyn}_0(G)\) 处的类型为 \(a(x)\in\mathrm T\)，这里 \(a:\operatorname{Dyn}_0(G)\to\mathrm T_S\) 定义于 \SGARefLink{sga3:xxiv:subsection:5-2}{5.2}。}

\emph{(ii) 上述分解诱导下列竖写的彼此同构的正合序列。这里 \(\underline{\operatorname{Aut}}_{S,a}(\mathfrak D_0)\) 表示 \(\operatorname{Dyn}_0(G)\) 中与 \(a:\operatorname{Dyn}_0(G)\to\mathrm T_S\) 交换的自同构所成的概形：}
\[
\begin{tikzcd}[column sep=large,row sep=large]
1 \arrow[d] & 1 \arrow[d] \\
\displaystyle
\prod_{t\in\mathrm T}\prod_{\mathfrak D_{0,t}/S}
  \underline{\operatorname{Aut}}_{\mathrm{gr.}}(G_{0,t})
  \arrow[r,"\sim"] \arrow[d]
&
\displaystyle
\prod_{\operatorname{Dyn}_0(G)/S}
  \underline{\operatorname{Aut}}_{\mathrm{gr.}}(G_0)
  \arrow[d]
\\
\underline{\operatorname{Aut}}_{S\text{-}\mathrm{gr.}}(G)
  \arrow[r,equal] \arrow[d]
&
\underline{\operatorname{Aut}}_{S\text{-}\mathrm{gr.}}(G)
  \arrow[d]
\\
\displaystyle\prod_{t\in\mathrm T}
  \underline{\operatorname{Aut}}_S
    \bigl(\operatorname{Dyn}_{0,t}(G)\bigr)
  \arrow[r,"\sim"] \arrow[d]
&
\underline{\operatorname{Aut}}_{S,a}
  \bigl(\operatorname{Dyn}_0(G)\bigr)
  \arrow[d]
\\
1 & 1
\end{tikzcd}
\]

\par\medskip\noindent\SGAInlineRefTarget{sga3:xxiv:corollary:5-11}\textbf{推论 5.11.---}\enspace
\emph{在上述假设下，以下三个范畴等价：}
\begin{enumerate}
\item[(i)] \(G\) 下主齐性丛所成范畴；
\item[(ii)] \(G_0\) 下主齐性丛所成范畴；
\item[(iii)] 各 \(G_{0,t}\)，\(t\in\mathrm T\)，下主齐性丛所成范畴的乘积。
\end{enumerate}
这形式地由上述分解与 \SGARefLink{sga3:xxiv:proposition:8-4}{8.4} 推出。

\par\medskip\noindent\SGAInlineRefTarget{sga3:xxiv:corollary:5-12}\textbf{推论 5.12.---}\enspace
\emph{有典范同构}
\[
  \underline{\operatorname{Tor}}(G)
  \simeq
  \prod_{t\in\mathrm T}
  \prod_{\mathfrak D_{0,t}/S}
    \underline{\operatorname{Tor}}(G_{0,t})
  \simeq
  \prod_{\operatorname{Dyn}_0(G)/S}
    \underline{\operatorname{Tor}}(G_0),
\]
\emph{将 \(\underline{\operatorname{Tor}}\) 换成 \(\underline{\operatorname{Bor}}\)，分别换成 \(\underline{\operatorname{Kil}}\)，也成立。}

这立即由分裂情形得到。

\par\medskip\noindent\SGAInlineRefTarget{sga3:xxiv:remark:5-13}\textbf{注 5.13.---}\enspace
典范态射
\[
  \operatorname{Dyn}(G)\longrightarrow\operatorname{Dyn}_0(G)
\]
使我们可将 \(\operatorname{Dyn}(G)\) 视为 \(\operatorname{Dyn}_0(G)\) 上的 Dynkin 概形；事实上，它就是 \(G_0\) 的 Dynkin 概形 \(\operatorname{Dyn}(G_0)\)。

同样，若 \(T\subset B\) 是 \(G\) 的 Killing 偶，典范地对应于 \(G_0\) 的 Killing 偶 \(T_0\subset B_0\)，则由 \SGARefLink{sga3:xxiv:3.9}{3.9} 属于
\[
  \operatorname{Pic}\bigl(\operatorname{Dyn}(G)\bigr)
  =
  \operatorname{Pic}\bigl(\operatorname{Dyn}(G_0)\bigr),
\]
的 \(G\) 与 \(G_0\) 的拟分裂障碍相同。因此：

\par\medskip\noindent\SGAInlineRefTarget{sga3:xxiv:corollary:5-14}\textbf{推论 5.14.---}\enspace
\emph{下列条件等价：}
\begin{enumerate}
\item[(i)] \(G\) 可拟分裂；
\item[(ii)] \(G_0\) 可拟分裂；
\item[(iii)] 每个 \(G_{0,t}\)，\(t\in\mathrm T\)，都可拟分裂。
\end{enumerate}

\section*{\SGAInlineRefTarget{sga3:xxiv:section:6}6. 约化群的 Borel 子群的自同构}

\par\medskip\noindent\SGAInlineRefTarget{sga3:xxiv:lemma:6-1}\textbf{引理 6.1.---}\enspace
\emph{设 \(S\) 是概形，\((G,T,M,R)\) 是分裂 \(S\)-群，\(\Delta\) 是 \(R\) 的单根系，\(B\) 是相应的 Borel 子群。则 \(B^u\) 作为 \textup{(fppf)} 群层，由各 \(U_\alpha\)，\(\alpha\in\Delta\)，生成。}

设 \(H\) 是由各 \(U_\alpha\)，\(\alpha\in\Delta\)，生成的 \(B^u\) 的子群层。对 \(\operatorname{ord}(\beta)=\operatorname{ord}_\Delta(\beta)>0\)（参见\SGARefLink{sga3:XXI:3.2.15}{第 XXI 讲 3.2.15}）归纳，证明对 \(\beta\in R_+\) 有 \(H\supset U_\beta\)。当 \(\operatorname{ord}(\beta)=1\) 时，断言由假设成立。设 \(\operatorname{ord}(\beta)>1\)，并假设只要 \(\operatorname{ord}(\gamma)<\operatorname{ord}(\beta)\)，就有 \(U_\gamma\subset H\)。存在 \(\alpha\in\Delta\)，使得 \(\beta-\alpha\in R_+\)（\SGARefLink{sga3:XXI:3.1.1}{第 XXI 讲 3.1.1}）。设 \(p\) 是使 \(\beta-p\alpha=\beta'\in R_+\) 的最大整数。则 \(U_\alpha,U_{\beta'}\subset H\)，且 \(\beta'-\alpha\notin R\)。因此只须证明：

\par\medskip\noindent\SGAInlineRefTarget{sga3:xxiv:lemma:6-2}\textbf{引理 6.2.---}\enspace
\emph{沿用\SGARefLink{sga3:XXIII:6.4}{第 XXIII 讲 6.4} 的记号。设 \(p=1\)，即 \(\beta-\alpha\) 不是根。若 \(H\) 是 \(B^u\) 的子群层，且 \(U_\alpha,U_\beta\subset H\)，则只要 \(i\alpha+j\beta\in R\)、\(i>0\)、\(j>0\)，便有 \(U_{i\alpha+j\beta}\subset H\)。}

按照 \(q=0,1,2,3\) 区分四种情形。下文中，\(x\)、\(y\) 表示 \(\mathbb G_{a,S'}\) 的任意截面，其中 \(S'\to S\)。

当 \(q=0\) 时，无须证明。

当 \(q=1\) 时，有
\[
 p_{\alpha+\beta}(\pm x)
 =p_\beta(-y)p_\alpha(-x)p_\beta(y)p_\alpha(x)\in H(S'),
\]
因而 \(U_{\alpha+\beta}\subset H\)。

当 \(q=2\) 时，同样有
\[
 p_{\alpha+\beta}(\pm xy)
 p_{2\alpha+\beta}(\pm x^2y)
 =p_\beta(-y)p_\alpha(-x)p_\beta(y)p_\alpha(x)\in H(S').
\]
因此，改变某些符号后，有
\[
 F(x,y)=p_{\alpha+\beta}(xy)p_{2\alpha+\beta}(x^2y)\in H(S').
\]
由于 \(p_{\alpha+\beta}\) 与 \(p_{2\alpha+\beta}\) 交换，
\[
 F(x,1)F(1,-x)=p_{2\alpha+\beta}(x^2-x).
\]
\setcounter{footnote}{27}
\footnote{\textit{编者注。} 将 \(F(-1,x)\) 改为 \(F(1,-x)\)。}
若 \(a\in\mathbb G_a(S)\)，则方程 \(X^2-X=a\) 定义了 \(S\) 的一个自由满扩张，即 \(\operatorname{Spec}\mathcal O_S[X]/(X^2-X-a)\)。因此 \(p_{2\alpha+\beta}(a)\in H(S')\)，故 \(U_{2\alpha+\beta}\subset H\)；再由 \(F(1,a)\in H(S')\)，得 \(U_{\alpha+\beta}\subset H\)。

当 \(q=3\) 时，同样有
\[
 F(x,y)=
 p_{\alpha+\beta}(xy)
 p_{2\alpha+\beta}(x^2y)
 p_{3\alpha+\beta}(x^3y)
 p_{3\alpha+2\beta}(x^3y^2)\in H(S'),
\]
并且由于
\[
 p_{3\alpha+2\beta}(\pm x)
 =F(1,1)^{-1}p_\beta(-x)F(1,1)p_\beta(x)\in H(S'),
\]
得到 \(U_{3\alpha+2\beta}\subset H\)，因此
\[
 K(x,y)=
 p_{\alpha+\beta}(xy)
 p_{2\alpha+\beta}(x^2y)
 p_{3\alpha+\beta}(x^3y)\in H(S').
\]
模 \(U_{3\alpha+2\beta}\) 计算
\[
 K(x+y,1)K(-x,1)K(1,y)^{-1}
\]
得到
\[
 p_{2\alpha+\beta}(2x^2+y^2+2xy-y)
 p_{3\alpha+\beta}(y^3+3xy^2+3x^2y-y)\in H(S').
\]
若 \(a\in\mathbb G_a(S)\)，则“方程组”
\[
  x^2=-xy-y+1-a,\qquad
  y^2=3y-2+3a
\]
定义了 \(S\) 的自由满扩张。此时 \(y^3+3xy^2+3x^2y-y=0\)，上述表达式给出 \(p_{2\alpha+\beta}(a)\in H(S')\)。\footnote{\textit{编者注。} 将 \(p_{3\alpha+\beta}(a)\) 改为 \(p_{2\alpha+\beta}(a)\)，并补上论证的末尾。}
所以 \(H\) 包含 \(U_{2\alpha+\beta}\) 与 \(U_{3\alpha+2\beta}\)，而且
\[
 E(x,y)=p_{\alpha+\beta}(xy)p_{3\alpha+\beta}(x^3y)\in H(S').
\]
由于 \(p_{\alpha+\beta}(xy)\) 与 \(p_{3\alpha+\beta}\) 交换，上述计算适用：
\[
 E(1,x)E(1,-x)=p_{3\alpha+\beta}(x^3-x).
\]
故 \(U_{3\alpha+\beta}\subset H\)，进而 \(U_{\alpha+\beta}\subset H\)。

\par\medskip\noindent\SGAInlineRefTarget{sga3:xxiv:remark:6-2-1}\textbf{注 6.2.1.---}\enspace
上述证明表明，假设“\(H\) 包含 \(U_\alpha\) 和 \(U_\beta\)”可换成：“\(H\) 包含 \(U_\alpha\) 或 \(U_\beta\)，并且 \(H\) 在 \(\operatorname{int}(U_\alpha)\) 和 \(\operatorname{int}(U_\beta)\) 下不变。”

\par\medskip\noindent\SGAInlineRefTarget{sga3:xxiv:theorem:6-3}\textbf{定理 6.3.---}\enspace
\emph{设 \(S\) 是概形，\(G\)、\(G'\) 是两个半单 \(S\)-群，\(B\)（分别为 \(B'\)）是 \(G\)（分别为 \(G'\)）的 Borel 子群。每个同构 \(B\xrightarrow{\sim}B'\) 都由唯一的同构 \(G\xrightarrow{\sim}G'\) 诱导。}

断言对于平展拓扑是局部的，因此可设 \(G\) 分裂：\(G=(G,T,M,R)\)，其中 \(B\) 由 \(R\) 的正根系 \(R_+\) 定义。给定的同构 \(u:B\xrightarrow{\sim}B'\) 诱导从 \(T\) 到 \(B'\) 的一个极大环面 \(T'\)（因而也是 \(G'\) 的极大环面）的同构。给定的同构 \(T\simeq D_S(M)\) 给出同构 \(T'\simeq D_S(M)\)，在这个同构下，\(R_+\) 的元素成为 \(B'\) 关于 \(T'\) 的常根，因而 \(R\) 的元素成为 \(G'\) 关于 \(T'\) 的常根。由于 \(G\)、\(G'\) 半单，余根由对偶性确定（\SGARefLink{sga3:XXI:1.2.5}{第 XXI 讲 1.2.5}）；这说明 \((T',M,R)\) 是 \(G'\) 的一个分裂结构，使得 \(\mathscr R(G)=\mathscr R(G')\)。

应用\SGARefLink{sga3:XXIII:5.1}{第 XXIII 讲 5.1}（唯一性定理），得到唯一的同构 \(G\xrightarrow{\sim}G'\)，它在 \(T\) 以及各 \(U_\alpha\)，\(\alpha\in\Delta\)，上与 \(u\) 一致。由 \SGARefLink{sga3:xxiv:lemma:6-1}{6.1}，它在 \(B\) 上的限制就是 \(u\)。

\par\medskip\noindent\SGAInlineRefTarget{sga3:xxiv:remark:6-3-1}\textbf{注 6.3.1.---}\enspace
(i) 使用\SGARefLink{sga3:xxii:prop:4.1.9}{第 XXII 讲 4.1.9}，并仿照该处 4.2.12 的论证，可以将定理中的“同构”换成“同源”（分别为“中心同源”）。

(ii) 对约化群，该定理不成立。例如取 \(G=G'=\mathrm{SL}_2\times\mathbb G_m\)，将它等同于群
\[
 \left\{
 \begin{pmatrix}
   a&b&0\\
   c&d&0\\
   0&0&h
 \end{pmatrix}
 \ \middle|\ ad-bc=1,\ h\text{ 可逆}
 \right\}.
\]
% 已译源行 2898--3340。

\setcounter{footnote}{29}
取 \(B=B'\) 为由 \(c=0\) 定义的 Borel 子群，并令 \(u\) 为 \(B\) 的如下自同构：
\[
  u(a,b,d,h)=(a,b,d,ah).
\]

\par\medskip\noindent\SGAInlineRefTarget{sga3:xxiv:corollary:6-4}\textbf{推论 6.4.---}\enspace
\emph{函子
\[
  (G,B)\longmapsto B
\]
从有序对 \((G,B)\) 所成范畴到 \(S\)-群概形范畴完全忠实；这里 \(G\) 是半单 \(S\)-群，\(B\) 是 \(G\) 的 Borel 子群，两个范畴均以同构为态射。}

\par\medskip\noindent\SGAInlineRefTarget{sga3:xxiv:corollary:6-5}\textbf{推论 6.5.---}\enspace
\emph{设 \(S\) 是概形，\(G\) 是半单 \(S\)-群，\(B\) 是 \(G\) 的 Borel 子群，\(B'\) 是在 \textup{(fpqc)} 拓扑下局部同构于 \(B\) 的 \(S\)-群。则 \(B'\) 在局部有限平展拓扑下局部同构于 \(B\)，\footnote{\textit{编者注。} 在\SGARefLink{sga3:iv:subsection:6:3}{第 IV 讲 6.3} 中，这一拓扑记为 \textup{(étf)}。换言之，对每个 \(s\in S\)，存在 \(s\) 的开邻域 \(U\) 及有限满平展态射 \(V\to U\)，使得 \(B'\times_S V\simeq B\times_S V\)。}而且存在半单 \(S\)-群 \(G'\)，使得 \(B'\) 是它的 Borel 子群。此外，\(G'\) 在相差一个在 \(B'\) 上诱导恒等的唯一同构的意义下唯一。}

\par\medskip\noindent\SGAInlineRefTarget{sga3:xxiv:corollary:6-6}\textbf{推论 6.6.---}\enspace
\emph{设 \(S\) 是概形，\(G\) 是半单 \(S\)-群，\(B\) 是 \(G\) 的 Borel 子群。则 \(\underline{\operatorname{Aut}}_{S\text{-}\mathrm{gr.}}(B)\) 可由一个仿射光滑 \(S\)-概形表示，\(\underline{\operatorname{Autext}}(B)\) 可由一个平展有限 \(S\)-概形表示，并且显然的态射诱导正合序列之间的同构。}
\[
\begin{tikzcd}[column sep=small,row sep=large]
1 \arrow[r]
&
\underline{B}_{\mathrm{ad}}(G)
  \arrow[r]
  \arrow[d,"\wr" description]
&
\underline{\Aut}_{S\text{-}\mathrm{gr.}}(G,B)
  \arrow[r]
  \arrow[d,"\wr" description]
&
\underline{\Autext}(G)
  \arrow[r]
  \arrow[d,"\wr" description]
&
1
\\
1 \arrow[r]
&
\underline{B}_{\mathrm{ad}}
  \arrow[r]
&
\underline{\Aut}_{S\text{-}\mathrm{gr.}}(G)
  \arrow[r]
&
\underline{\Autext}(B)
  \arrow[r]
&
1 .
\end{tikzcd}
\]

这立即由 \SGARefLink{sga3:xxiv:2.1}{2.1} 及上述结果推出。请读者在这一框架下建立第 1、2、3、4 节结果的类似形式。

\par\medskip\noindent\SGAInlineRefTarget{sga3:xxiv:remark:6-7}\textbf{注 6.7.---}\enspace
若 \(S\) 是概形，\(B\) 是 \(S\)-群，则 \(B\) 至多能成为一个半单群的 Borel 子群，该半单群由 \(B\) 唯一确定。尚须刻画哪些 \(S\)-群 \(B\) 是半单群的 Borel 子群。\footnote{\textit{编者注。} 可以问，若一个光滑仿射 \(S\)-群的几何纤维都是半单群的 Borel 子群，则它本身是否也是某个半单 \(S\)-群的 Borel 子群。我们删去了原文中以域 \(k\) 上的对偶数概形为 \(S\) 的反例，因为正如 M.~Demazure 所指出的，这个反例是错误的。}

\section*{\SGAInlineRefTarget{sga3:xxiv:section:7}7. 当 \(G\) 约化时函子 \(\underline{\operatorname{Hom}}_{S\text{-}\mathrm{gr.}}(G,H)\) 的可表性}
\subsection*{\SGAInlineRefTarget{sga3:xxiv:subsection:7-1}7.1. 分裂情形}
\subsubsection*{7.1.1.}
设 \(S\) 是概形，\(G\) 是钉扎 \(S\)-群，\(T\) 是它的极大环面，\(\Delta\) 是单根集。对 \(\alpha\in\Delta\)，设
\[
  u_\alpha\in U_\alpha^\times(S),
  \qquad
  w_\alpha\in\operatorname{Norm}_G(T)(S)
\]
是由钉扎定义的元素。

设 \(H\) 是 \(S\)-群概形。我们考察函子 \(\underline{\operatorname{Hom}}_{S\text{-}\mathrm{gr.}}(G,H)\)，更精确地，考察态射
\[
  q:\underline{\operatorname{Hom}}_{S\text{-}\mathrm{gr.}}(G,H)
  \longrightarrow
  \underline{\operatorname{Hom}}_{S\text{-}\mathrm{gr.}}(T,H).
\]
设
\[
  f_T:T\longrightarrow H
\]
是 \(S\)-群同态，并置 \(q^{-1}(f_T)=\mathscr F\)。于是
\[
  \mathscr F(S')
  =
  \left\{
    f\in\operatorname{Hom}_{S'\text{-}\mathrm{gr.}}
      (G_{S'},H_{S'})
    \,\middle|\,
    f|_{T_{S'}}=(f_T)_{S'}
  \right\}.
\]
有 \(S\)-函子的态射
\[
  i:\mathscr F\longrightarrow H^{2n},
\]
其中 \(n=\operatorname{Card}(\Delta)\)，它在点上的定义为
\[
  i(f)=\bigl(f(u_\alpha),f(w_\alpha)\bigr)_{\alpha\in\Delta}.
\]
由\SGARefLink{sga3:XXIII:1.9}{第 XXIII 讲 1.9}，\(i\) 是单态射。

\par\medskip\noindent\SGAInlineRefTarget{sga3:xxiv:proposition:7-1-2}\textbf{命题 7.1.2.---}\enspace
\emph{若 \(H\) 在 \(S\) 上分离，则 \(\mathscr F\) 可表，且 \(i\) 是闭浸入。}

相对可表性的通常技巧\footnote{\textit{编者注。} 例如参见\SGARefLink{sga3:xxii:thm:5.8.1}{第 XXII 讲 5.8.1} 的证明。}表明，只须证明以下断言。给定截面
\[
  v_\alpha,h_\alpha\in H(S),
  \qquad \alpha\in\Delta,
\]
则存在一个延拓 \((f_T)_{S'}\) 且满足
\[
  f(u_\alpha)=v_\alpha,
  \qquad
  f(w_\alpha)=h_\alpha
\]
的 \(S'\)-同态
\[
  f:G_{S'}\longrightarrow H_{S'}
\]
的那些 \(S\)-概形 \(S'\)，恰是经 \(S\) 的某个闭子概形分解的那些概形。显然可设 \(S\) 仿射。

为方便起见，称仿射概形的态射 \(Y\to X\)\footnote{\textit{编者注。} 将 \(X\to Y\) 改为 \(Y\to X\)。}满足条件 \textup{(L)}，如果 \(Y\) 是某个 \(\mathcal O(X)\)-代数的谱，而该代数作为 \(\mathcal O(X)\)-模自由。在仿射概形范畴中，满足 \textup{(L)} 的态射的乘积及复合仍满足 \textup{(L)}，且 \textup{(L)} 在基变换下稳定。

\par\medskip\noindent\SGAInlineRefTarget{sga3:xxiv:lemma:7-1-3}\textbf{引理 7.1.3.---}\enspace
\emph{假设 \(S\) 仿射，并设 \(\alpha\in\Delta\)。考虑态射
\[
  a:T\times_S T\longrightarrow G_{a,S},
  \qquad
  a(t,t')=\alpha(t)+\alpha(t').
\]
则：
\begin{enumerate}
\item[(i)] \(a\) 忠实平坦且有限呈示；
\item[(ii)] 若 \(R\) 是 \(a\) 的纤维平方，则结构态射 \(R\to S\) 满足 \textup{(L)}。
\end{enumerate}}

态射 \(a\) 是满的，因为 \(\alpha:T\to G_{m,S}\) 是满的，而且它显然有限呈示。由于 \(\alpha\) 满足 \textup{(L)}，对这两个断言，都只须证明
\[
  u:G_{m,S}^2\longrightarrow G_{a,S},
  \qquad
  u(x,y)=x+y
\]
满足 \textup{(L)}。等价地，对每个环 \(A\)，环
\[
  A[X,Y,X^{-1},Y^{-1}]
\]
必须在其子环 \(A[X+Y]\) 上自由。\footnote{\textit{编者注。} 简化了原文中的证明。}现在 \(A[X,Y]\) 在
\[
  B=A[X+Y,XY],
\]
上以 \(\{1,X\}\) 为基自由，因为 \(X^2-(X+Y)X+XY=0\)。所以
\[
  A[X^{\pm1},Y^{\pm1}]
  =A[X,Y]_{XY}
\]
在
\[
  B_{XY}=A[X+Y,(XY)^{\pm1}]
\]
上以 \(\{1,X\}\) 为基自由，并且在 \(A[X+Y]\) 上以
\[
  (XY)^p,\quad X(XY)^p,
  \qquad p\in\mathbb Z.
\]
为基自由。

\par\medskip\noindent\SGAInlineRefTarget{sga3:xxiv:lemma:7-1-4}\textbf{引理 7.1.4.---}\enspace
\emph{设 \(\alpha\in\Delta\)，并设
\[
  b:T\times_S T\longrightarrow H
\]
是在点上由
\[
  b(t,t')
  =
  \operatorname{int}(f_T(t))(v_\alpha)\,
  \operatorname{int}(f_T(t'))(v_\alpha).
\]
定义的态射。设 \(f_\alpha:U_\alpha\to H\) 是 \(S\)-态射。下列条件等价：\footnote{\textit{编者注。} 下文将 \(U_\alpha\) 的群律按加法记。若 \(p_\alpha:G_{a,S}\xrightarrow{\sim}U_\alpha\) 是满足 \(p_\alpha(1)=u_\alpha\) 的同构，\(x=p_\alpha(z)\)，则 \(\alpha(t)x\) 表示 \(p_\alpha(\alpha(t)z)\)，而 \(a(t,t')u_\alpha\) 表示 \(p_\alpha(a(t,t'))=\alpha(t)u_\alpha+\alpha(t')u_\alpha\)。}
\begin{enumerate}
\item[\SGAInlineRefTarget{sga3:xxiv:lemma:7-1-4:item:i}(i)] \(f_\alpha\) 是群同态，\(f_\alpha(u_\alpha)=v_\alpha\)，并且对所有 \(t\in T(S')\)、\(x\in U_\alpha(S')\)、\(S'\to S\)，有
  \[
    f_T(t)f_\alpha(x)f_T(t)^{-1}
    =
    f_\alpha(\alpha(t)x)
  \]
\item[\SGAInlineRefTarget{sga3:xxiv:lemma:7-1-4:item:ii}(ii)] \(f_\alpha(u_\alpha)=v_\alpha\)，并且对所有 \(t,t'\in T(S')\)、\(S'\to S\)，有
  \[
    f_\alpha(a(t,t')u_\alpha)=b(t,t')
  \]
\end{enumerate}}

若 \(f_\alpha\) 满足 \SGARefLink{sga3:xxiv:lemma:7-1-4:item:i}{(i)}，则
\[
  f_\alpha(\alpha(t)u_\alpha)
  =
  \operatorname{int}(f_T(t))(v_\alpha),
\]
立即得到 \SGARefLink{sga3:xxiv:lemma:7-1-4:item:ii}{(ii)}。反之，假设 \SGARefLink{sga3:xxiv:lemma:7-1-4:item:ii}{(ii)}。由于 \(a\) 是 \textup{(fpqc)} 覆盖，为证明 \SGARefLink{sga3:xxiv:lemma:7-1-4:item:i}{(i)} 的最后一个条件，只须证明对 \(t,t',t''\in T(S')\) 有
\[
  f_T(t)\,
  f_\alpha(a(t',t'')u_\alpha)\,
  f_T(t)^{-1}
  =
  f_\alpha(\alpha(t)a(t',t'')u_\alpha).
\]
这等同于
\[
  f_T(t)b(t',t'')f_T(t)^{-1}=b(tt',tt''),
\]
而此式直接由 \(b\) 的定义得到。

尚须证明 \(f_\alpha\) 是群同态。前一性质给出
\[
\begin{aligned}
 f_\alpha(\alpha(t)u_\alpha)
 f_\alpha(\alpha(t')u_\alpha)
 &=
 \bigl(f_T(t)f_\alpha(u_\alpha)f_T(t)^{-1}\bigr)
 \bigl(f_T(t')f_\alpha(u_\alpha)f_T(t')^{-1}\bigr)\\
 &=b(t,t')\\
  &=f_\alpha\bigl(\alpha(t)u_\alpha+
                   \alpha(t')u_\alpha\bigr).
\end{aligned}
\]
所以只要 \(x,x'\) 是概形稠密的开子集 \(U_\alpha^\times\subset U_\alpha\) 的截面，就有 \(f_\alpha(x+x')=f_\alpha(x)f_\alpha(x')\)。结果遂由\SGARefLink{sga3:xviii:prop:1.4}{第 XVIII 讲 1.4} 得到。

\subsubsection*{7.1.5.}
暂时固定 \(\alpha\in\Delta\)。由于 \(a\) 忠实平坦且拟紧，\(G_{a,S}\) 等同于 \(T\times_S T\) 对 \(a\) 所定义的等价关系 \(R\) 的商。以 \(i_1,i_2:R\to T\times_S T\) 表示它的两个投影。
\[
\begin{tikzcd}[column sep=huge]
R
  \arrow[r,shift left=.7ex,"i_1"]
  \arrow[r,shift right=.7ex,"i_2"']
&
T\times_S T
\end{tikzcd}
\]

\(b\) 经 \(a\) 分解\footnote{\textit{编者注。} 即条件 \SGARefLink{sga3:xxiv:lemma:7-1-4:item:ii}{7.1.4(ii)} 成立。}的充要条件是
\[
  b\circ i_1=b\circ i_2.
\]
设 \(K\) 是 \(b\circ i_1\) 与 \(b\circ i_2\) 的等化子。
\[
\begin{tikzcd}[column sep=huge]
R
  \arrow[r,shift left=.7ex,"b\circ i_1"]
  \arrow[r,shift right=.7ex,"b\circ i_2"']
&
H
\end{tikzcd}
\]
于是所需条件为 \(K=R\)。由于 \(H\) 在 \(S\) 上分离，\(K\) 是 \(R\) 的闭子概形。另一方面，由 \SGARefLink{sga3:xxiv:lemma:7-1-3}{7.1.3}，\(R\) 在 \(S\) 上本质自由。
% 已译源行 3341--3698。

\subsubsection*{7.1.6.}
由此可知，对 \(S\)-概形 \(S'\)，在 \(S'\) 上存在满足 \SGARefLink{sga3:xxiv:lemma:7-1-4:item:i}{7.1.4(i)} 的 \(f_\alpha\)（且存在时唯一），当且仅当
\[
  K_{S'}=R_{S'}
\]
并且所得态射满足
\[
  f_\alpha(u_\alpha)=v_\alpha.
\]
第一个条件等价于要求 \(S'\to S\) 经 \(S\) 的某个闭子概形分解（\SGARefLink{sga3:viii:theorem:6:4}{第 VIII 讲 6.4}\footnote{\textit{编者注。} 另见第 VI\textsubscript{B} 讲的增补 \SGARefLink{thm:sga3-VIB-6.2.3}{6.2.3}。}）。将 \(S\) 换成这个闭子概形后，第二个条件又定义了 \(S\) 的一个闭子概形，因为它表示分离 \(S\)-概形 \(H\) 的两个截面相等。

所以，将 \(S\) 换成所得的闭子概形后，可设存在满足 \SGARefLink{sga3:xxiv:lemma:7-1-4:item:i}{7.1.4(i)} 的态射 \(f_\alpha:U_\alpha\to H\)。将对所有 \(\alpha\in\Delta\) 得到的闭子概形相交，可设此条件对每个 \(\alpha\in\Delta\) 成立。

\subsubsection*{7.1.7.}
同样，对每个 \(\alpha\in\Delta\)，考虑两个 \(S\)-群同态
\[
  f_T\circ\operatorname{int}(w_\alpha),
  \qquad
  \operatorname{int}(h_\alpha)\circ f_T
  :T\longrightarrow H.
\]
由于 \(H\) 在 \(S\) 上分离，\(T\) 在 \(S\) 上本质自由，同样的论证表明，将 \(S\) 换成一个闭子概形后，可设对每个 \(\alpha\in\Delta\) 都有
\[
  f_T\circ\operatorname{int}(w_\alpha)
  =
  \operatorname{int}(h_\alpha)\circ f_T
\]

\subsubsection*{7.1.8.}
现在，生成元与关系定理（\SGARefLink{sga3:xxiii:subsection:3-5}{第 XXIII 讲 3.5}）表明，具有所需性质的群同态 \(f:G\to H\) 存在，当且仅当截面
\[
  h_\alpha,\qquad
  v_\alpha,\qquad
  f_T\bigl(\alpha^*(-1)\bigr),
  \qquad \alpha\in\Delta,
\]
之间的某个有限代数关系集成立：
\[
  R_i\!\left(
    (h_\alpha)_\alpha,
    (v_\alpha)_\alpha,
    \bigl(f_T(\alpha^*(-1))\bigr)_\alpha
  \right)=e,
  \qquad i=1,\ldots,n.
\]
由于 \(H\) 在 \(S\) 上分离，这些关系又定义了 \(S\) 的一个闭子概形，从而证明了\SGARefLink{sga3:xxiv:proposition:7-1-2}{命题 7.1.2}。

\par\medskip\noindent\SGAInlineRefTarget{sga3:xxiv:corollary:7-1-9}\textbf{推论 7.1.9.---}\enspace
\emph{设 \(S\) 是概形，\(G\) 是分裂约化 \(S\)-群，\(T\) 是它的极大环面，\(H\) 是分离 \(S\)-群概形。设 \(\mathcal P\) 是态射的一个性质，满足：
\begin{enumerate}
\item[(i)] 每个闭浸入都有 \(\mathcal P\)；
\item[(ii)] 具有 \(\mathcal P\) 的态射的复合也有 \(\mathcal P\)；
\item[(iii)] \(\mathcal P\) 在基变换下稳定；
\item[(iv)] 结构态射 \(H\to S\) 有 \(\mathcal P\)。
\end{enumerate}
则典范态射
\[
  \underline{\operatorname{Hom}}_{S\text{-}\mathrm{gr.}}(G,H)
  \longrightarrow
  \underline{\operatorname{Hom}}_{S\text{-}\mathrm{gr.}}(T,H)
\]
相对可由具有 \(\mathcal P\) 的分离态射表示。\footnote{\textit{编者注。} 更明确地，对每个可表的 \(S'\to\underline{\operatorname{Hom}}_{S\text{-}\mathrm{gr.}}(T,H)\)，纤维积 \(\underline{\operatorname{Hom}}_{S\text{-}\mathrm{gr.}}(G,H)\times S'\) 可表，且它到 \(S'\) 的态射分离并具有 \(\mathcal P\)。}}

事实上，可设 \(G\) 已钉扎。结构态射 \(H^{2n}\to S\) 有 \(\mathcal P\)，断言遂由 \SGARefLink{sga3:xxiv:proposition:7-1-2}{7.1.2} 得到。

\par\medskip\noindent\SGAInlineRefTarget{sga3:xxiv:corollary:7-1-10}\textbf{推论 7.1.10.---}\enspace
\emph{设 \(S\) 是概形，\(G\) 是分裂约化 \(S\)-群，\(H\) 是具有仿射纤维的光滑 \(S\)-群概形。\footnote{\textit{编者注。} 删去了 \(H\) 拟分离这一多余假设。}则 \(\underline{\operatorname{Hom}}_{S\text{-}\mathrm{gr.}}(G,H)\) 可由一个在 \(S\) 上局部有限呈示且分离的 \(S\)-概形表示。}

由于 \(H\) 在 \(S\) 上光滑，其单位连通分支 \(H^0\) 有定义；它具有仿射纤维，并在 \(S\) 上光滑、分离且有限呈示（第 VI\textsubscript{B} 讲，\SGARefLink{thm:sga3-VIB-3.10}{3.10} 与 \SGARefLink{sga3:xxiv:proposition:5-5}{5.5}）。由于 \(G\) 的纤维连通，可将 \(H\) 换成 \(H^0\)。此时 \(G\)、\(H\) 都有限呈示，因而可设 \(S\) 是 Noether 概形，并将 \SGARefLink{sga3:xxiv:corollary:7-1-9}{7.1.9} 用于性质 \(\mathcal P\)：“有限型”。由\SGARefLink{sga3:xv:theorem:8:8}{第 XV 讲 8.9}，\(\underline{\operatorname{Hom}}_{S\text{-}\mathrm{gr.}}(T,H)\) 可由一个分离的局部有限型 \(S\)-概形表示。

\par\medskip\noindent\SGAInlineRefTarget{sga3:xxiv:remark:7-1-11}\textbf{注 7.1.11.---}\enspace
若 \(H\) 在 \(S\) 上仿射，则可将对第 XV 讲的引用换成\SGARefLink{sga3:xi:corollary:4:2}{第 XI 讲 4.2}。

\subsection*{\SGAInlineRefTarget{sga3:xxiv:subsection:7-2}7.2. 一般情形}
\par\medskip\noindent\SGAInlineRefTarget{sga3:xxiv:proposition:7-2-1}\textbf{命题 7.2.1.---}\enspace
\emph{设 \(S\) 是概形，\(G\) 是约化 \(S\)-群，\(T\) 是 \(G\) 的极大环面。设 \(H\) 是 \(S\)-群概形，其结构态射 \(H\to S\) 满足：
\begin{quote}
\textup{(+)} 每个点 \(s\in S\) 都有一个开邻域 \(U\)，使得 \(H_U\to U\) 拟射影。
\end{quote}
则典范态射
\[
  \underline{\operatorname{Hom}}_{S\text{-}\mathrm{gr.}}(G,H)
  \longrightarrow
  \underline{\operatorname{Hom}}_{S\text{-}\mathrm{gr.}}(T,H)
\]
相对可由满足 \textup{(+)} 的态射表示。}

当 \(G\) 关于 \(T\) 分裂时，取 \(\mathcal P=\textup{(+)}\) 应用 7.1.9。当 \(G\) 局部等平凡时，例如 \(G\) 半单时（4.2.2），注意断言对于局部有限平展拓扑是局部的。事实上，拟射影性对于全局有限平展拓扑是局部的，并保证此拓扑下下降的有效性（SGA~1，VIII，7.7）。一般情形由下一引理推出。

\par\medskip\noindent\SGAInlineRefTarget{sga3:xxiv:lemma:7-2-2}\textbf{引理 7.2.2.---}\enspace
\emph{设 \(S\) 是概形，\(G\) 是约化 \(S\)-群，\(T\) 是 \(G\) 的极大环面，\(G'\) 是 \(G\) 的导群（第 XXII 讲 6.2），\(T'=T\cap G'\) 是与 \(T\) 相应的 \(G'\) 的极大环面（第 XXII 讲 6.2.8）。由 \(G,G',T,T'\) 组成的自然图在 \(S\)-群层范畴中是余笛卡尔图。等价地，对每个 \(S\)-群层 \(H\)，相应的 Hom 函子方图是笛卡尔图。}
\[
\begin{tikzcd}[column sep=large,row sep=large]
G & G' \arrow[l] \\
T \arrow[u] & T' \arrow[l] \arrow[u]
\end{tikzcd}
\]
\[
\begin{tikzcd}[column sep=large,row sep=large]
\underline{\operatorname{Hom}}_{S\text{-}\mathrm{gr.}}(G,H)
  \arrow[r]
  \arrow[d]
&
\underline{\operatorname{Hom}}_{S\text{-}\mathrm{gr.}}(G',H)
  \arrow[d]
\\
\underline{\operatorname{Hom}}_{S\text{-}\mathrm{gr.}}(T,H)
  \arrow[r]
&
\underline{\operatorname{Hom}}_{S\text{-}\mathrm{gr.}}(T',H)
\end{tikzcd}
\]

若 \(\operatorname{rad}(G)\) 是 \(G\) 的根基，则 \(\operatorname{rad}(G)\subset T\)，故
\[
  \operatorname{rad}(G)\cap G'
  =
  \operatorname{rad}(G)\cap T'
  =
  K.
\]
\(G\) 中的乘法诱导同源（\SGARefLink{sga3:xxii:subsection:6-2}{第 XXII 讲 6.2}）
\[
  i:G'\times_S\operatorname{rad}(G)\longrightarrow G,
  \qquad
  j:T'\times_S\operatorname{rad}(G)\longrightarrow T.
\]
设
\[
  f_{G'}:G'\to H,\qquad
  f_T:T\to H,\qquad
  f_{T'}:T'\to H
\]
是满足
\[
  f_{G'}|_{T'}=f_T|_{T'}=f_{T'}.
\]
的 \(S\)-群同态。我们证明，存在唯一的 \(S\)-群同态 \(f:G\to H\)，诱导 \(f_{G'}\) 与 \(f_T\)。置
\[
  f_{\mathrm{rad}}
  =
  f_T|_{\operatorname{rad}(G)}
\]
\footnote{\textit{编者注。} 将 \(f_r\) 换成了 \(f_{\mathrm{rad}}\)。}
并定义
\[
  f_1=f_{G'}\cdot f_{\mathrm{rad}}:
  G'\times_S\operatorname{rad}(G)\longrightarrow H.
\]
\(f\) 存在（而且此时显然唯一）的充要条件，是 \(f_1\) 在 \(i\) 的核上平凡，即 \(f_{G'}\) 与 \(f_{\mathrm{rad}}\) 在 \(K\) 上一致。同样，由 \(f_T\) 的存在可知，\(f_{T'}\) 与 \(f_{\mathrm{rad}}\) 在 \(K\) 上一致。最后，\(f_{G'}\) 与 \(f_{T'}\) 在 \(K\subset T'\) 上一致，引理得证。

仿照 \SGARefLink{sga3:xxiv:corollary:7-1-10}{7.1.10} 的论证，由\SGARefLink{sga3:xxiv:proposition:7-2-1}{命题 7.2.1} 得：
\par\medskip\noindent\SGAInlineRefTarget{sga3:xxiv:corollary:7-2-3}\textbf{推论 7.2.3.---}\enspace
\emph{设 \(S\) 是概形，\(G\) 是约化 \(S\)-群，\(H\) 是具有仿射纤维的光滑拟射影 \(S\)-群概形。则 \(\underline{\operatorname{Hom}}_{S\text{-}\mathrm{gr.}}(G,H)\) 可由一个在 \(S\) 上局部有限呈示且分离的 \(S\)-概形表示。}

\subsection*{\SGAInlineRefTarget{sga3:xxiv:subsection:7-3}7.3. 特征零所特有的现象}
若 \(G\)、\(H\) 是光滑 \(S\)-群，Lie 代数为 \(\mathfrak g\)、\(\mathfrak h\)，则有典范态射
\[
  \Lie:
  \underline{\operatorname{Hom}}_{S\text{-}\mathrm{gr.}}(G,H)
  \longrightarrow
  \underline{\operatorname{Hom}}_{\mathcal O_S\text{-}\mathrm{Lie}}
    (\mathfrak g,\mathfrak h),
\]
其中右端函子的意义是显然的。

\par\medskip\noindent\SGAInlineRefTarget{sga3:xxiv:proposition:7-3-1}\textbf{命题 7.3.1.---}\enspace
\emph{设 \(S\) 是特征零的概形，即 \(\mathbb Q\)-概形；设 \(G\) 是约化 \(S\)-群，\(H\) 是具有仿射纤维的光滑拟射影 \(S\)-群概形。
\begin{enumerate}
\item[\SGAInlineRefTarget{sga3:xxiv:proposition:7-3-1:item:i}(i)] \(\underline{\operatorname{Hom}}_{S\text{-}\mathrm{gr.}}(G,H)\) 可由一个光滑分离 \(S\)-概形表示。
\item[(ii)] 若 \(G\) 半单，则这个概形在 \(S\) 上仿射且有限呈示。
\item[\SGAInlineRefTarget{sga3:xxiv:proposition:7-3-1:item:iii}(iii)] 若 \(G\) 单连通（\SGARefLink{sga3:xxii:def:4.3.3}{第 XXII 讲 4.3.3}），则典范态射
\[
  \Lie:
  \underline{\operatorname{Hom}}_{S\text{-}\mathrm{gr.}}(G,H)
  \longrightarrow
  \underline{\operatorname{Hom}}_{\mathcal O_S\text{-}\mathrm{Lie}}
    (\mathfrak g,\mathfrak h)
\]
是同构。
\end{enumerate}}
% 已译源行 3699--3975。

\par\medskip\noindent\SGAInlineRefTarget{sga3:xxiv:lemma:7-3-2}\textbf{引理 7.3.2.---}\enspace
\emph{设 \(k\) 是特征零的域，\(G\) 是约化 \(k\)-群，\(V\) 是有限维 \(k\)-向量空间，\(G\to\operatorname{GL}(V)\) 是线性表示。则
\[
  H^0(G,V)=H^0(\mathfrak g,V),
  \qquad
  H^i(G,V)=0\quad(i>0).
\]}

第一个等式一般对光滑连通群成立。\footnote{\textit{编者注。} 包含关系 \(V^G\subset V^{\mathfrak g}\) 显然。反之，\(C=\operatorname{Centr}_G(V^{\mathfrak g})\) 是 \(G\) 的闭子群，由于 \(\operatorname{char}(k)=0\)，它是光滑的。由\SGARefLink{sga3:ii:theorem:5-3-1}{第 II 讲 5.3.1}，\(\Lie(C)=\operatorname{Centr}_{\mathfrak g}(V^{\mathfrak g})=\mathfrak g\)；由于 \(C\) 光滑、\(G\) 连通，\(C=G\)。所以 \(V^{\mathfrak g}\subset V^G\)。另见 \([{\rm DG70}]\)，\S II.6，命题 2.1(c)。}对约化群，也可如下证明。可设 \(k\) 代数闭，因而 \(G\) 分裂，且由同构于 \(G_{m,k}\) 的子群生成；\footnote{\textit{编者注。} \(G\) 的极大环面的并在 \(G\) 中稠密；参见 \emph{Bible}，\S6.5，定理 5，或 \([{\rm Ch05}]\)，\S6.6，定理 6。}而对这个群，断言是直接的。

第一个等式表明，当 \(G\) 半单时，函子 \(H^0(G,V)\) 对 \(V\) 正合：此时 \(\mathfrak g\) 是半单 Lie 代数，应用 \([{\rm BLie}]\)，\S I.6，习题 1(b) 即可。当 \(G\) 约化时，断言仍然成立。事实上，引入 \(G\) 的根基 \(R\)\footnote{\textit{编者注。} 注意 \(R\) 是环面。}及半单商群 \(G/R\)，有
\[
  H^0(G,V)=H^0\bigl(G/R,H^0(R,V)\bigr),
\]
所以 \(H^0(G,-)\) 是两个正合函子的复合。\SGARefLink{sga3:i:theorem:5-3-1}{第 I 讲 5.3.1} 于是给出：对 \(i>0\)，有 \(H^i(G,V)=0\)。

\par\medskip\noindent\SGAInlineRefTarget{sga3:xxiv:remark:7-3-3}\textbf{注 7.3.3.---}\enspace
在上述假设下，若 \(G\) 半单，则
\[
  H^1(\mathfrak g,V)=H^2(\mathfrak g,V)=0;
\]
参见 \([{\rm BLie}]\)，同上，(b) 和 (d)。

\subsubsection*{7.3.4.}
现在证明命题。由 \SGARefLink{sga3:xxiv:corollary:7-2-3}{7.2.3}，\(\underline{\operatorname{Hom}}_{S\text{-}\mathrm{gr.}}(G,H)\) 已可由一个在 \(S\) 上局部有限呈示且分离的 \(S\)-概形表示。为证明它光滑，只须证明它无穷小光滑（\SGARefLink{sga3:xi:remark:1:8}{第 XI 讲 1.8}），\footnote{\textit{编者注。} 这意味着，对每个闭点为 \(s'\) 的局部 Artin \(S\)-概形 \(S'\)，每个 \(\kappa(s')\)-群同态 \(G_{s'}\to H_{s'}\) 都可提升为 \(S'\)-群同态 \(G_{S'}\to H_{S'}\)。由\SGARefLink{sga3:iii:corollary:2-8}{第 III 讲 2.8}，这来自 \(H^2(G_{s'},V)\) 的消失，其中 \(V=\Lie(H_{s'}/\kappa(s'))\)。}而这由\SGARefLink{sga3:iii:corollary:2-8:item:i}{第 III 讲 2.8(i)} 及 \SGARefLink{sga3:xxiv:lemma:7-3-2}{7.3.2} 推出。这就证明了 \SGARefLink{sga3:xxiv:proposition:7-3-1:item:i}{(i)}。

其次，说明 \SGARefLink{sga3:iii:corollary:2-8:item:ii}{(ii)} 如何由 \SGARefLink{sga3:xxiv:proposition:7-3-1:item:iii}{(iii)} 推出。首先，只须证明 \(\underline{\operatorname{Hom}}_{S\text{-}\mathrm{gr.}}(G,H)\) 在 \(S\) 上仿射；因为由 \SGARefLink{sga3:xxiv:proposition:7-3-1:item:i}{(i)} 它光滑，此时便有限呈示。无论如何，
\[
  \underline{\operatorname{Hom}}_{\mathcal O_S\text{-}\mathrm{Lie}}
    (\mathfrak g,\mathfrak h)
\]
由
\[
  \underline{\operatorname{Hom}}_{\mathcal O_S\text{-}\mathrm{mod}}
    (\mathfrak g,\mathfrak h)
  \simeq
  \mathrm W\!\left(
    \mathfrak g^\vee\otimes_{\mathcal O_S}\mathfrak h
  \right),
\]
的一个闭子概形表示，而后者是仿射 \(S\)-概形。因此，当 \(G\) 单连通时，所需结论成立。

一般情形可设 \(G\) 分裂，从而 \(G\simeq\operatorname{\acute Ep}_S(\mathscr R)\)。引入单连通根数据 \(\operatorname{sc}(\mathscr R)\)（\SGARefLink{sga3:XXI:6.5.5}{第 XXI 讲 6.5.5}）。存在性定理（\SGARefLink{sga3:xxv:theorem:1:1}{第 XXV 讲 1.1}）给出一个单连通 \(S\)-群 \(\widetilde G\) 及中心同源
\[
  \widetilde G\longrightarrow G.
\]
它的核 \(K\) 是有限可对角化 \(S\)-群，因 \(S\) 特征为零，故它是扭常 \(S\)-群。函子 \(\underline{\operatorname{Hom}}_{S\text{-}\mathrm{gr.}}(K,H)\) 显然由一个分离 \(S\)-概形表示。例如，若 \(K\simeq(\mathbb Z/n\mathbb Z)_S\)，则它同构于
\[
  \ker\!\left(H\xrightarrow{\,n\,}H\right).
\]
有带点 \(S\)-概形的正合序列
\[
  1\longrightarrow
  \underline{\operatorname{Hom}}_{S\text{-}\mathrm{gr.}}(G,H)
  \xrightarrow{u}
  \underline{\operatorname{Hom}}_{S\text{-}\mathrm{gr.}}
    (\widetilde G,H)
  \longrightarrow
  \underline{\operatorname{Hom}}_{S\text{-}\mathrm{gr.}}(K,H).
\]
所以 \(u\) 是闭浸入，\(\underline{\operatorname{Hom}}_{S\text{-}\mathrm{gr.}}(G,H)\) 在 \(S\) 上仿射。

\SGAInlineRefTarget{sga3:xxiv:subsubsection:7-3-5}\subsubsection*{7.3.5.}
尚须证明 \SGARefLink{sga3:xxiv:proposition:7-3-1:item:iii}{(iii)}。仿照 \SGARefLink{sga3:xxiv:proposition:7-3-1:item:i}{(i)} 的证明，利用 \SGARefLink{sga3:xxiv:remark:7-3-3}{7.3.3}，可证
\[
  \underline{\operatorname{Hom}}_{\mathcal O_S\text{-}\mathrm{Lie}}
    (\mathfrak g,\mathfrak h)
\]
在 \(S\) 上光滑。因此可设 \(S=\operatorname{Spec}(k)\)，其中 \(k\) 代数闭且特征为零。只须证明 \(\Lie\) 在 \(k\)-点上为双射，并在对应点处的切空间之间诱导双射。先证明
\[
  \operatorname{Hom}_{k\text{-}\mathrm{gr.}}(G,H)
  \longrightarrow
  \operatorname{Hom}_{k\text{-}\mathrm{Lie}}
    (\mathfrak g,\mathfrak h)
\]
为双射。

若 \(u:G\to H\) 是 \(k\)-群同态，则它的图是 \(G\times_k H\) 的连通子群，由其 Lie 代数确定；该 Lie 代数就是 \(\Lie(u)\) 的图。因此映射为单射。反之，若 \(v:\mathfrak g\to\mathfrak h\) 是 Lie 代数同态，则它的图 \(\mathfrak k\) 是 \(\mathfrak g\times\mathfrak h\) 的 Lie 子代数，同构于 \(\mathfrak g\)。由于 \(\mathfrak g\) 等于自身的导代数，\(\mathfrak k\) 也如此。因而 Chevalley 的定理 \([{\rm Ch51}]\)，\S14，定理 15，表明 \(\mathfrak k\) 是 \(G\times_k H\) 的某个连通子群 \(K\) 的 Lie 代数。\footnote{\textit{编者注。} 另见 \([{\rm Bo91}]\)，II 7.9。下面一句也是增补的。}

由于 \(\mathfrak k\simeq\mathfrak g\) 半单，\(K\) 是半单 \(k\)-群（法文此处印为 \(G\)；根据本段论证及下一句，应为 \(K\)）。此外，
\[
  \dim K=\dim\mathfrak k=\dim\mathfrak g=\dim G,
\]
且 \(K\cap H\) 有限，因为它的 Lie 代数为零。所以典范态射
\[
  \operatorname{pr}_1:K\longrightarrow G
\]
有限且占优；由于 \(G\) 连通，它是满的，因而是同源。又由于 \(G\) 单连通、\(K\) 半单，\SGARefLink{sga3:XXI:6.2.7}{第 XXI 讲 6.2.7} 表明 \(\operatorname{pr}_1\) 是同构。因此，\(K\) 是某个满足 \(\Lie(u)=v\) 的 \(k\)-群同态 \(u:G\to H\) 的图。

最后，设 \(u:G\to H\) 是 \(k\)-群同态。
\[
  \underline{\operatorname{Hom}}_{k\text{-}\mathrm{gr.}}(G,H)
\]
在 \(u\) 处的切空间等同于 \(Z^1(G,\mathfrak h)\)（\SGARefLink{sga3:ii:subsection:4-2}{第 II 讲 4.2}），其中 \(G\) 经 \(\operatorname{Ad}_H\circ u\) 作用于 \(\mathfrak h\)。同样，
\[
  \underline{\operatorname{Hom}}_{k\text{-}\mathrm{Lie}}
    (\mathfrak g,\mathfrak h)
\]
在 \(\Lie(u)\) 处的切空间等同于 \(Z^1(\mathfrak g,\mathfrak h)\)。须证明
\[
  Z^1(G,\mathfrak h)\longrightarrow
  Z^1(\mathfrak g,\mathfrak h)
\]
为双射。
\[
\begin{tikzcd}[column sep=small,row sep=large]
0 \arrow[r]
&
\mathfrak h^G
  \arrow[r]
  \arrow[d]
&
\mathfrak h
  \arrow[r]
  \arrow[d,equal]
&
Z^1(G,\mathfrak h)
  \arrow[r]
  \arrow[d]
&
H^1(G,\mathfrak h)
  \arrow[d]
\\
0 \arrow[r]
&
\mathfrak h^{\mathfrak g}
  \arrow[r]
&
\mathfrak h
  \arrow[r]
&
Z^1(\mathfrak g,\mathfrak h)
  \arrow[r]
&
H^1(\mathfrak g,\mathfrak h)
\end{tikzcd}
\]
由 \SGARefLink{sga3:xxiv:lemma:7-3-2}{7.3.2} 与 \SGARefLink{sga3:xxiv:remark:7-3-3}{7.3.3}，
\[
  \mathfrak h^G=\mathfrak h^{\mathfrak g},
  \qquad
  H^1(G,\mathfrak h)=0=H^1(\mathfrak g,\mathfrak h),
\]
即得所需结论。

\par\medskip\noindent\SGAInlineRefTarget{sga3:xxiv:remark:7-3-6}\textbf{注 7.3.6.---}\enspace
以下断言大概成立：若 \(S\) 局部 Noether，\(G\) 是单连通半单 \(S\)-群，\(H\) 是光滑 \(S\)-群概形，\(\widehat G,\widehat H\) 是它们沿单位截面的形式完备化，则每个同态 \(\widehat G\to\widehat H\) 都来自唯一的同态 \(G\to H\)。这将推广 \SGARefLink{sga3:xxiv:proposition:7-3-1:item:iii}{7.3.1(iii)}。当 \(S=\operatorname{Spec}(k)\)，而 \(H\) 仿射且有限型时，这来自 Dieudonné 的一个定理 \([{\rm Di57}]\)，\S15，定理 4。\footnote{\textit{编者注。} 形式群 \(\widehat G\)、\(\widehat H\) 的定义参见第 VII\textsubscript{B} 讲。设 \(S=\operatorname{Spec}(k)\)，其中 \(k\) 是域。由该处 2.2.1，给定形式 \(k\)-群态射 \(v:\widehat G\to\widehat H\)，等价于给定 Hopf \(k\)-代数态射 \(\phi:\operatorname{Dist}(G)\to\operatorname{Dist}(H)\)，其中 \(\operatorname{Dist}(G)\) 表示 \(G\) 的（原点处）分布代数；参见第 VII\textsubscript{A} 讲 \SGARefLink{sga3:VIIA:2.1}{2.1}，\([{\rm DG70}]\)，\S II.4，\SGARefLink{sga3:xxiv:lemma:6-1}{6.1}，或 \([{\rm Ja87}]\)，I 7.7。在 \([{\rm Di57}]\) 和 \([{\rm Ta75}]\) 中，它称为 \(G\) 的超代数。\([{\rm Di57}]\) 的定理 4（另见 \([{\rm Ta75}]\)，0.3.4 (f)、(g)）在这一框架下推广了 \SGARefLink{sga3:xxiv:subsubsection:7-3-5}{7.3.5} 中使用的 Chevalley 定理：得到一个连通闭 \(k\)-子群 \(K\subset G\times H\)，使得 \(\widehat K\) 等于 \(v\) 的图。由于 \(\operatorname{Dist}(K)\to\operatorname{Dist}(G)\) 是同构，\(K\to G\) 有限平展。因此 \(K\) 半单，再由\SGARefLink{sga3:XXI:6.2.7}{第 XXI 讲 6.2.7}，\(K\to G\) 是同构（因为 \(G\) 单连通；另见 \([{\rm Ta75}]\)，\SGARefLink{sga3:xxiv:1.8}{1.8} 与 \SGARefLink{sga3:xxiv:2.2}{2.2}）。更一般地，该处研究了具有性质 \textup{(SC)} 的 \(k\)-群 \(G\)：每个有限平展 \(k\)-群态射 \(G'\to G\)，只要 \(G'\) 连通，就是同构。最后，上述论证表明，每个有限维 \(\operatorname{Dist}(G)\)-模都唯一地成为 \(G\)-模；关于推广到分裂约化 \(k\)-群 \(G\)（甚至推广到 \(G\) 的 Borel 子群），参见 \([{\rm Ja87}]\)，II 1.20（及该处参考文献）。}
% 已译源行 3976--4245。

\subsection*{\SGAInlineRefTarget{sga3:xxiv:subsection:7-4}7.4. 一个例子}
作为例子，我们来确定
\[
  \underline{\operatorname{Hom}}_{S\text{-}\mathrm{gr.}}
    (\mathrm{SL}_{2,S},\mathrm{SL}_{2,S}).
\]

\subsubsection*{7.4.1.}
回忆（第 XX 讲第 5 节），\(\mathrm{SL}_{2,S}\) 是 \(S\) 上满足 \(ad-bc=1\) 的矩阵
\[
  \begin{pmatrix}a&b\\c&d\end{pmatrix}
\]
所成的 \(S\)-群概形。一个极大环面 \(T\) 由单态射
\[
  \alpha^*:\mathbb G_{m,S}\longrightarrow\mathrm{SL}_{2,S},
  \qquad
  \alpha^*(z)=
  \begin{pmatrix}z&0\\0&z^{-1}\end{pmatrix}.
\]
定义。\(G\) 关于 \(T\) 的一个根由
\[
  \alpha\bigl(\alpha^*(z)\bigr)=z^2,
\]
定义，相应的单态射为
\[
  p:\mathbb G_{a,S}\longrightarrow\mathrm{SL}_{2,S},
  \qquad
  p(x)=
  \begin{pmatrix}1&x\\0&1\end{pmatrix}.
\]
最后，置 \(u=p(1)\)，相应的 Weyl 群代表元为
\[
  w=
  \begin{pmatrix}0&1\\-1&0\end{pmatrix}.
\]

又回忆（\SGARefLink{sga3:xx:6.2}{第 XX 讲 6.2}），\(\mathrm{SL}_{2,S}\) 由 \(T\)、\(p(\mathbb G_{a,S})\) 及 \(w\) 生成，满足关系
\[
  \alpha^*(z)p(x)\alpha^*(z^{-1})=p(z^2x),
\]
\[
  w\alpha^*(z)w^{-1}=\alpha^*(z^{-1}),
  \qquad
  w^2=\alpha^*(-1),
  \qquad
  (wu)^3=1.
\]

\subsubsection*{7.4.2.}
设 \(f\) 是 \(\mathrm{SL}_{2,S}\) 的自同态。它首先定义同态
\[
  f\circ\alpha^*:\mathbb G_{m,S}\longrightarrow\mathrm{SL}_{2,S}.
\]
闭子群 \(\ker(f\circ\alpha^*)\subset\mathbb G_{m,S}\) 在 \(S\) 上局部地等于 \(\mu_{n,S}\)，其中 \(n\geq1\)，或等于 \(\mathbb G_{m,S}\)。因此，将 \(S\) 限制到适当部分后，可设存在整数 \(n\in\mathbb N\) 以及单态射
\[
  f':\mathbb G_{m,S}\longrightarrow\mathrm{SL}_{2,S}
\]
使得
\[
  f\circ\alpha^*(z)=f'(z^n).
\]
由单态射 \(\mathbb G_{m,S}\to\mathrm{SL}_{2,S}\) 的共轭性，在作平展覆盖基扩张后，存在 \(G\) 的截面 \(g\)，使得
\[
  f\circ\alpha^*(z)
  =
  \operatorname{int}(g)\circ\alpha^*(z^n).
\]
将 \(f\) 换成其一个内共轭后，便归结到对某个 \(n\in\mathbb N\) 有
\[
  f\circ\alpha^*(z)=\alpha^*(z^n)
\]
的情形。

\subsubsection*{7.4.3.}
现在考虑
\[
  f\circ p:\mathbb G_{a,S}\longrightarrow\mathrm{SL}_{2,S}.
\]
它满足
\[
  \alpha^*(z^n)(f\circ p)(x)\alpha^*(z^n)^{-1}
  =
  (f\circ p)(z^2x).
\]
若 \(n=0\)，则立即得到 \(f\circ p\) 在 \(\mathbb G_{a,S}\) 的位似下不变，因而为常值。若 \(n\neq0\)，则可应用\SGARefLink{sga3:xxii:prop:4.1.9}{第 XXII 讲 4.1.9}：映射 \(x\mapsto x^n\) 是 \(\mathbb G_{a,S}\) 的自同态，而且存在 \(\lambda\in\mathbb G_{m,S}\)，使得
\[
  f\circ p(x)=p(\lambda x^n).
\]
当 \(n=0\) 时，取 \(\lambda=1\)，此关系也成立。再次扩张 \(S\) 后，选取满足 \(z^2=\lambda\) 的 \(z\in\mathbb G_{m,S}\)。将 \(f\) 换成 \(\operatorname{int}(\alpha^*(z))\circ f\)，便归结到
\[
  f\circ\alpha^*(z)=\alpha^*(z^n),
  \qquad
  f\circ p(x)=p(x^n).
\]

\subsubsection*{7.4.4.}
最后，必须有
\[
  f(w)\alpha^*(z^n)f(w)^{-1}=\alpha^*(z^n)^{-1},
  \qquad
  (f(w)u)^3=1.
\]
由\SGARefLink{sga3:xx:3.8}{第 XX 讲 3.8}，这蕴含 \(f(w)=w^n\)。由于 \(G\) 由 \(T\)、\(p(\mathbb G_{a,S})\) 及 \(w\) 生成，得到对每个 \(S'\to S\) 及每个 \(\begin{pmatrix}a&b\\c&d\end{pmatrix}\in G(S')\)，有
\[
  f\!\begin{pmatrix}a&b\\c&d\end{pmatrix}
  =
  \begin{pmatrix}a^n&b^n\\c^n&d^n\end{pmatrix}.
\]

\subsubsection*{7.4.5.}
综上，在 \(S\) 上对平展拓扑局部地，\(\mathrm{SL}_{2,S}\) 的每个自同态 \(f\) 都可写成
\[
  f=\phi\circ F_n,
\]
其中 \(\phi\) 是内自同构，\(n\geq0\)，并且
\[
  F_n\!\begin{pmatrix}a&b\\c&d\end{pmatrix}
  =
  \begin{pmatrix}a^n&b^n\\c^n&d^n\end{pmatrix}.
\]
若 \(f=\phi\circ F_n\)，则整数 \(n\) 由任一纤维 \(f_s\) 确定，例如由 \(\ker(f_s)\) 确定。因此 \(n\) 是 \(S\) 上的局部常值函数。

\subsubsection*{7.4.6.}
由此可知，对 \(\mathrm{SL}_{2,S}\) 的每个自同态 \(f\)，概形 \(S\) 都有一个典范的开闭子概形分解
\[
  S_0,\qquad S_1,\qquad S_{p^n},
\]
其中 \(p^n\) 遍历素数的正整数次幂，使得：
\begin{enumerate}
\item[(i)] \(f_{S_0}\) 是零态射；
\item[(ii)] \(f_{S_1}\) 是同构，且必为内自同构；
\item[(iii)] \(S_{p^n}\) 的特征为 \(p\)，而且 \(f_{S_{p^n}}\) 有唯一分解
  \[
    f_{S_{p^n}}=\phi\circ F_p^n,
  \]
  其中 \(\phi\) 是内自同构，\(F_p\) 是 \(\mathrm{SL}_{2,\mathbb F_p}\) 的 Frobenius 自同态。
\end{enumerate}

\subsubsection*{7.4.7.}
等价地，
\[
  \underline{\operatorname{Hom}}_{\mathbb Z\text{-}\mathrm{gr.}}
    (\mathrm{SL}_{2,\mathbb Z},\mathrm{SL}_{2,\mathbb Z})
\]
是由以下各概形组成的不交并概形：
\begin{enumerate}
\item[(i)] 一个同构于 \(\operatorname{Spec}(\mathbb Z)\) 的概形；
\item[(ii)] 一个同构于
  \[
    \underline{\operatorname{Aut}}_{\mathbb Z\text{-}\mathrm{gr.}}
      (\mathrm{SL}_{2,\mathbb Z})
    \simeq
    \operatorname{ad}(\mathrm{SL}_{2,\mathbb Z});
  \]
  的概形；
\item[(iii)] 对每个素数 \(p\) 及每个整数 \(n>0\)，一个同构于
  \[
    \underline{\operatorname{Aut}}_{\mathbb F_p\text{-}\mathrm{gr.}}
      (\mathrm{SL}_{2,\mathbb F_p})
    \simeq
    \operatorname{ad}(\mathrm{SL}_{2,\mathbb F_p}).
  \]
  的概形。
\end{enumerate}

\subsubsection*{7.4.8.}
特别地：
\begin{enumerate}
\item[(i)] \(\underline{\operatorname{Hom}}_{\mathbb F_p\text{-}\mathrm{gr.}}(\mathrm{SL}_{2,\mathbb F_p},\mathrm{SL}_{2,\mathbb F_p})\) 有无穷多个连通分支，因而不拟紧。
\item[(ii)] 若 \(S\) 是不等特征的概形，则 \(\underline{\operatorname{Hom}}_{S\text{-}\mathrm{gr.}}(\mathrm{SL}_{2,S},\mathrm{SL}_{2,S})\) 在 \(S\) 上不平坦。
\end{enumerate}
% 已译源行 4246--4461。

\section*{\SGAInlineRefTarget{sga3:xxiv:section:8}8. 附录：Hensel 环上光滑群的上同调\\上同调与函子 \(\prod\)}

\par\medskip\noindent\SGAInlineRefTarget{sga3:xxiv:proposition:8-1}\textbf{命题 8.1.---}\enspace
\emph{设 \(S\) 是 Hensel 局部概形，\(s\) 是它的闭点，\(G\) 是光滑 \(S\)-群概形，且 \(G\) 的每个有限子集都包含于一个仿射开子集。%
\begingroup
\renewcommand{\thefootnote}{*}%
\footnote{最后这个条件实际上不必要；参见 A.~Grothendieck，\textit{Le groupe de Brauer III}，载于 \textit{Dix Exposés sur la cohomologie des schémas}，North-Holland，1968，定理 11.7 及注 11.8.3。}%
\endgroup
\addtocounter{footnote}{-1}
\begin{enumerate}
\item[(i)] 若 \(P\) 是 \(G\) 下的主齐性丛，则存在平凡化 \(P\) 的有限满平展态射 \(S'\to S\)。此外，
  \[
    \operatorname{Fib}(S,G)\simeq H^1_{\acute et}(S,G).
  \]
\item[\SGAInlineRefTarget{sga3:xxiv:proposition:8-1:item:ii}(ii)] 对每个有限满平展态射 \(S'\to S\)，典范映射
  \[
    H^1(S'/S,G)
    \longrightarrow
    H^1(S'\otimes_S\kappa(s)/\kappa(s),G_s)
  \]
  为双射。
\item[(iii)] 典范映射
  \[
    \operatorname{Fib}(S,G)
    \longrightarrow
    \operatorname{Fib}(\kappa(s),G_s)
  \]
  为双射。
\end{enumerate}}

\subsubsection*{8.1.2.}
若 \(K\) 是 \(\kappa(s)\) 的有限可分扩张，则存在有限满平展态射 \(S'\to S\)，使得
\[
  K\simeq S'\otimes_S\kappa(s).
\]
\footnote{\textit{编者注。} 概形 \(S'\) 仍为 Hensel 局部概形。我们也保留了原文的编号：没有第 8.1.1 小节。}
若 \(P\) 是 \(G\) 下的主齐性丛，则 \(P\) 在 \(S\) 上光滑，故 \(P_s\) 在 \(\kappa(s)\) 上光滑。因此存在有限可分扩张 \(K/\kappa(s)\)，使得 \(P_K\) 有截面（EGA IV\textsubscript{4}，17.15.10）。按上述方式表示 \(K\)，Hensel 引理（EGA IV\textsubscript{4}，18.5.17）表明 \(P_{S'}\) 有截面。这证明了 (i) 的第一个断言。

反之，若 \(P\) 是平展拓扑下 \(G\) 的主齐性层，则有有限满平展态射 \(S'\to S\) 平凡化 \(P\)：Hensel 局部概形的每个平展覆盖族都由一个只包含一个这种形式的态射的族所支配。由对 \(G\) 所作的下降假设，\(P\) 可表（SGA~1，VIII，7.6）。这完成了 (i) 的证明，并表明 (ii) 蕴含 (iii)。尚须证明 (ii)。

\subsubsection*{8.1.3.}
(ii) 中的映射为单射。设 \(P\)、\(Q\) 是由 \(S'\) 平凡化的 \(G\) 下两个主齐性丛，考虑 \(S\)-群层
\[
  \mathcal H=
  \underline{\operatorname{Isom}}_{G\text{-}\mathrm{bundles}}(P,Q).
\]
由于 \(\mathcal H_{S'}\) 同构于 \(G_{S'}\)，仿照上述论证，由 \(G\) 的第二个假设可知 \(\mathcal H\) 可表。若 \(\mathcal H(\kappa(s))\neq\varnothing\)，则由 Hensel 引理有 \(\mathcal H(S)\neq\varnothing\)，故 \(P\simeq Q\)。

\subsubsection*{8.1.4.}
为证明 \SGARefLink{sga3:xxiv:proposition:8-1:item:ii}{(ii)} 中的满射性，等价地只须证明
\[
  Z^1(S'/S,G)
  \longrightarrow
  Z^1(S'\otimes_S\kappa(s)/\kappa(s),G_s)
\]
为满射。设 \(Z^1\) 是由
\[
  Z^1(T)=Z^1(S'\times_ST/T,G_T).
\]
定义的 \(S\)-函子。于是上述映射就是
\[
  Z^1(S)\longrightarrow Z^1(\kappa(s)).
\]
再用 Hensel 引理，将断言归结为证明 \(Z^1\) 可由光滑 \(S\)-概形表示。

\subsubsection*{8.1.5.}
先证明 \(Z^1\) 可由局部有限呈示的 \(S\)-概形表示。对 \(i=0,1,\ldots\)，由
\[
  C_i(T)=C_i(S'\times_ST/T,G_T).
\]
定义 \(S\)-函子 \(C_i\)。等价地，
\[
  C_i(T)
  =
  G\bigl((S'/S)^{i+1}\times_ST\bigr),
  \qquad
  C_i
  =
  \underline{\operatorname{Hom}}_S((S'/S)^{i+1},G).
\]
由于 \(Z^1\) 可由 \(C_1\)、\(C_2\) 通过纤维积得到，只须证明对 \(i=1,2\)，\(C_i\) 可由局部有限呈示的 \(S\)-概形表示。

\subsubsection*{8.1.6.}
设 \(T\to S\) 是使 \(S'\) 分裂的有限满平展态射。则
\[
  C_i\times_ST
  =
  \underline{\operatorname{Hom}}_T
    ((S'\times_ST/T)^{i+1},G_T)
\]
由 \(G_T\) 的 \(n\) 个副本的乘积表示，其中 \(n\) 为 \((S'/S)^{i+1}\) 的次数。再次应用对 \(G\) 所作的假设，可知 \(C_i\) 由局部有限呈示的 \(S\)-概形表示（SGA~1，VIII，同上）。

\subsubsection*{8.1.7.}
尚须证明 \(Z^1\) 光滑。按定义，若 \(T\) 是仿射 \(S\)-概形，\(T_0\) 是由平方为零的理想 \(J\) 定义的闭子概形，则须证明
\[
  Z^1(T)\longrightarrow Z^1(T_0)
\]
为满射。由于 \(G\) 光滑，典范映射 \(G(T)\to G(T_0)\) 为满射，所以只须证明
\[
  H^1(S'\times_ST/T,G_T)
  \longrightarrow
  H^1(S'\times_ST_0/T_0,G_{T_0})
\]
为双射。稍改记号并无害地推广假设后，便归结为下一引理。

\par\medskip\noindent\SGAInlineRefTarget{sga3:xxiv:lemma:8-1-8}\textbf{引理 8.1.8.---}\enspace
\emph{设 \(S\)、\(S'\) 是仿射概形，\(S'\to S\) 忠实平坦，\(J\) 是 \(S\) 上平方为零的理想，\(S_0\) 是由它定义的闭子概形。设 \(G\) 是光滑 \(S\)-群。置
\[
  S'_0=S'\times_SS_0,
  \qquad
  G_0=G\times_SS_0.
\]
则典范映射
\[
  H^1(S'/S,G)\longrightarrow H^1(S'_0/S_0,G_0)
\]
为双射。}

\par\medskip\noindent\SGAInlineRefTarget{sga3:xxiv:remark:8-1-9}\textbf{注 8.1.9.---}\enspace
若 \(G\) 交换，则对每个 \(H^i\)，\(i>0\)，相同断言也成立，证明类似。

\smallskip\noindent\emph{证明。}
设 \(M_0\) 是拟凝聚 \(\mathcal O_{S_0}\)-模
\[
  M_0
  =
  \Lie(G_0/S_0)\otimes_{\mathcal O_{S_0}}J.
\]
对每个 \(S_0\)-概形 \(T_0\)，置
\[
  M_0(T_0)
  =
  H^0\!\left(
    T_0,
    M_0\otimes_{\mathcal O_{S_0}}\mathcal O_{T_0}
  \right),
\]
并以
\[
  M(T)=M_0(T_0),
  \qquad
  \overline G(T)=G_0(T_0),
  \qquad
  T_0=T\times_SS_0.
\]
定义 \(S\)-群函子
\[
  M=\prod_{S_0/S}M_0,
  \qquad
  \overline G=\prod_{S_0/S}G_0
\]
由\SGARefLink{sga3:iii:corollary:0-9}{第 III 讲 0.9} 及 (\SGARefLink{sga3:iii:remark:0-6:equation:0-6-2}{0.6.2})，对每个仿射 \(S\)-概形 \(T\)，有函子性的正合序列
\[
  1\longrightarrow M(T)\longrightarrow G(T)
  \longrightarrow\overline G(T)\longrightarrow1.
\]
须研究映射
\[
  H^1(S'/S,G)
  \longrightarrow
  H^1(S'_0/S_0,G_0)
  =
  H^1(S'/S,\overline G).
\]

先设 \(G\) 交换。此时有上同调正合序列
\[
  H^1(S'/S,M)
  \longrightarrow H^1(S'/S,G)
  \longrightarrow H^1(S'/S,\overline G)
  \longrightarrow H^2(S'/S,M).
\]
此外，
\[
  H^i(S'/S,M)
  =
  H^i(S'_0/S_0,M_0)
  =
  H^i(S'_0/S_0,\mathcal M_0)
  =
  0
  \qquad(i\neq0),
\]
最后一个等式由 TDTE I，B，引理 1.1 得到。

现在设 \(G\) 不交换。使用非交换上同调正合序列。若 \(u\in Z^1(S'/S,\overline G)\)，则 \(H^1(S'/S,G)\) 中在 \(H^1(S'/S,\overline G)\) 的像等于 \(u\) 的类的元素，属于相应上边缘映射
\[
  H^1(S'/S,M_u)\longrightarrow H^1(S'/S,G),
\]
的像，其中 \(M_u\) 是 \(M\) 通过 \(u\) 的扭转。同样，对 \(v\in Z^1(S'/S,\overline G)\)，有障碍
\[
  \Delta(v)\in H^2(S'/S,M_v)
\]
它为零，当且仅当 \(v\) 的类属于 \(H^1(S'/S,G)\) 的像。因此只须证明
\[
  H^1(S'/S,M_u)=0,
  \qquad
  H^2(S'/S,M_v)=0.
\]

回忆（\SGARefLink{sga3:iii:proposition:0-8}{第 III 讲 0.8}），上述正合序列定义的 \(G\) 在 \(M\) 上的作用，是函子性地由伴随表示
\[
  \operatorname{ad}:G_0
  \longrightarrow
  \underline{\operatorname{Aut}}_{\mathcal O_{S_0}}
    (\Lie(G_0/S_0)).
\]
诱导的。因此上闭链 \(u\)、\(v\) 分别定义了 \(\Lie(G_0/S_0)\) 上的下降资料；所得拟凝聚 \(\mathcal O_{S_0}\)-模分别记为 \(L_u\)、\(L_v\)。对 \(T\to S\)，
\[
  M_u(T)
  =
  H^0\!\left(
    T_0,
    L_u\otimes_{\mathcal O_{S_0}}J
      \otimes_{\mathcal O_{S_0}}\mathcal O_{T_0}
  \right),
\]
将 \(u\) 换成 \(v\) 也类似。因此
\[
  H^1(S'/S,M_u)=H^1(S'_0/S_0,L_u\otimes J),
\]
\[
  H^2(S'/S,M_v)=H^2(S'_0/S_0,L_v\otimes J),
\]
由已用过的结果，两个群均为零。
\hfill\(\square\)

\par\medskip\noindent\SGAInlineRefTarget{sga3:xxiv:proposition:8-2}\textbf{命题 8.2.---}\enspace
\emph{设 \(\mathcal C\) 是带有纤维积的范畴，赋予比典范拓扑更粗的一个拓扑。设 \(S'\to S\) 是 \(\mathcal C\) 的态射，\(G'\) 是 \(S'\)-群层，
\[
  G=\prod_{S'/S}G'
\]
是相应的 \(S\)-群层。设
\[
  H^1_S(S',G')\subset H^1(S',G')
\]
是由满足如下条件的 \(G'\) 下主齐性层的类所成子集：该层由 \(S'\) 上的某个筛平凡化，而这个筛是由 \(S\) 上适当的覆盖筛作基变换得到的。典范映射
\[
  H^1(S,G)\longrightarrow H^1(S',G'),
  \qquad
  P\longmapsto P\times_SS',
\]
诱导双射
\[
  H^1(S,G)\xrightarrow{\sim}H^1_S(S',G').
\]
其逆映射由
\[
  P'\longmapsto\prod_{S'/S}P'.
\]
诱导。}

对 \(\mathcal C/S\) 的每个对象 \(X\)，按定义有关于 \(X\) 函子性的同构
\[
  \operatorname{Hom}_S(X,G)
  \xrightarrow{\sim}
  \operatorname{Hom}_{S'}(X\times_SS',G').
\]
所以对每个 \(S\)-对象 \(T\)，有关于 \(T\) 函子性的双射
\[
  H^1(T/S,G)
  \xrightarrow{\sim}
  H^1(T'/S',G'),
  \qquad
  T'=T\times_SS'.
\]
将单个态射 \(T\to S\) 换成 \(S\) 的任意覆盖族，再取归纳极限，即得命题的第一部分。第二部分随即得到。

\par\medskip\noindent\SGAInlineRefTarget{sga3:xxiv:lemma:8-3}\textbf{引理 8.3.---}\enspace
\emph{在 \SGARefLink{sga3:xxiv:proposition:8-2}{8.2} 的假设下，等式
\[
  H^1_S(S',G')=H^1(S',G')
\]
在 \(S\) 上是局部的。更精确地，设存在覆盖族 \(\{S_i\to S\}\)，使得对每个 \(i\)，
\[
  H^1_{S_i}(S'\times_SS_i,G')
  =
  H^1(S'\times_SS_i,G').
\]
则 \(H^1_S(S',G')=H^1(S',G')\)。}

事实上，设 \(P'\) 是 \(G'\) 下的主齐性层，并置
\[
  P'_i=P'\times_{S'}(S'\times_SS_i).
\]
由假设，存在覆盖族 \(\{S_{ij}\to S_i\}\)，使得对每个 \(j\)，
\[
  P'\times_{S'}(S'\times_SS_{ij})
\]
有截面。但 \(\{S_{ij}\to S\}\) 是 \(S\) 的覆盖族，故 \(P'\) 由这个族经基变换得到的 \(S'\) 的覆盖族平凡化。

\par\medskip\noindent\SGAInlineRefTarget{sga3:xxiv:proposition:8-4}\textbf{命题 8.4.---}\enspace
\emph{设 \(S'\to S\) 是概形的有限平展态射。设 \(G'\) 是 \(S'\)-群层，
\[
  G=\prod_{S'/S}G'
\]
是相应的 \(S\)-群层。对于平展拓扑、局部有限平展拓扑以及 \textup{(fpqc)} 拓扑，下列函子诱导互逆的双射。}
\[
\begin{tikzcd}[column sep=huge,row sep=large]
\mathcal P
  \arrow[r,mapsto]
&
\mathcal P\times_S S'
\\
\displaystyle\prod_{S'/S}\mathcal P'
&
\mathcal P'
  \arrow[l,mapsto]
\end{tikzcd}
\]
\emph{特别地，
\[
  H^1(S,G)\simeq H^1(S',G').
\]}

由 \SGARefLink{sga3:xxiv:proposition:8-2}{8.2}，只须证明
\[
  H^1_S(S',G')=H^1(S',G').
\]
由 \SGARefLink{sga3:xxiv:lemma:8-3}{8.3}，可对局部有限平展拓扑作局部验证。因此可设 \(S'\) 是有限多个 \(S\) 的副本的不交并，记为 \(I_S\)，其中 \(I\) 为有限集。于是 \(G'\) 由 \(S\) 上的一族层 \((G_i)_{i\in I}\) 给出，而且
\[
  H^1(S',G')
  \simeq
  \prod_{i\in I}H^1(S,G_i).
\]
另一方面，
\[
  H^1(S,G)
  \simeq
  \prod_{i\in I}H^1(S,G_i).
\]
由 \SGARefLink{sga3:xxiv:proposition:8-2}{8.2} 得 \(H^1_S(S',G')=H^1(S',G')\)。
\hfill\(\square\)

\par\medskip\noindent\SGAInlineRefTarget{sga3:xxiv:remark:8-5}\textbf{注 8.5.---}\enspace
\SGARefLink{sga3:xxiv:proposition:8-2}{命题 8.2} 与 \SGARefLink{sga3:xxiv:proposition:8-4}{8.4} 可通过以下正合序列解释，其中 \(f:S'\to S\) 是给定态射：
\[
  1\longrightarrow H^1(S,f_*G')
  \longrightarrow H^1(S',G')
  \longrightarrow H^0(S,R^1f_*G').
\]
在交换情形，这个正合序列来自 Leray 谱序列。它在非交换情形仍成立。\footnote{\textit{编者注。} 参见 \([{\rm Gi71}]\) 的 \S V.3.1.4。}

以这种形式表述后，可以看到：若使用平展拓扑，则只假设 \(f\) 有限，甚至只假设它整，结果仍然成立。事实上，对每个 \(G'\)，由 SGA~4，VIII，5.3，有
\[
  R^1f_*(G')=\text{终层}
\]
\footnote{\textit{编者注。} 这一引用也指向 \([{\rm Gi71}]\) 的 \S V.3.1.4。}

对 \textup{(fpqc)}、\textup{(fppf)} 这样的拓扑，结果不成立，即使
\[
  S=\operatorname{Spec}(k),
  \qquad
  S'=\operatorname{Spec}(k[t]/t^2),
\]
其中 \(k\) 是特征为 \(p\neq0\) 的代数闭域，而 \(G'=\mu_p\) 或 \(\alpha_p\)。

同样，若删去 \(f\) 的有限性假设，则即使对平展拓扑，8.2 也不再成立。例如，取 \(f\) 为开浸入，且
\[
  S=\operatorname{Spec}(V),
\]
其中 \(V\) 是剩余域代数闭的完备离散赋值环；取 \(S'\) 为由泛点组成的开子概形，并取 \(G'\) 为常群 \((\mathbb Z/n\mathbb Z)_{S'}\)，其中 \(n\) 与 \(V\) 的剩余特征互素。则
\[
  H^1(S,G)=0,
  \qquad
  H^1(S',G')\neq0.
\]
将 \(S'\) 换成 \(S\amalg S'\)，得到类似例子，其中 \(S'\to S\) 平展且满，因而是所考虑拓扑的一个覆盖。

\section*{参考文献}
\noindent[BLie]\enspace
N.~Bourbaki, \textit{Groupes et algèbres de Lie}, Chap.~I, Hermann,
1960.\par

\smallskip\noindent[Ch51]\enspace
C.~Chevalley, \textit{Théorie des groupes de Lie}, t.~II,
\textit{Groupes algébriques}, Hermann, 1951.\par

\smallskip\noindent[Di57]\enspace
J.~Dieudonné, Lie groups and Lie hyperalgebras over a field of
characteristic \(p>0\). VI, \textit{Amer. J. Math.} \textbf{79}
(1957), no.~2, 331--388.\footnote{\textit{编者注。} 以上三条参考文献见于原文，下面的文献由编者增补。}\par

\smallskip\noindent[Bo91]\enspace
A.~Borel, \textit{Linear Algebraic Groups}, 第 2 版,
Springer-Verlag, 1991.\par

\smallskip\noindent[Ch05]\enspace
C.~Chevalley, \textit{Classification des groupes algébriques
semi-simples}, P.~Cartier、A.~Grothendieck、M.~Lazard 协作,
\textit{Collected Works}, vol.~3,
Springer, 2005.\par

\smallskip\noindent[DG70]\enspace
M.~Demazure、P.~Gabriel, \textit{Groupes algébriques}, Masson、
North-Holland, 1970.\par

\smallskip\noindent[Gi71]\enspace
J.~Giraud, \textit{Cohomologie non abélienne}, Springer-Verlag,
1971.\par

\smallskip\noindent[Ja87]\enspace
J.~C.~Jantzen, \textit{Representations of Algebraic Groups}, Academic
Press, 1987; 第 2 版, American Mathematical Society, 2003.\par

\smallskip\noindent[Jou83]\enspace
J.-P.~Jouanolou, \textit{Théorèmes de Bertini et applications},
Birkhäuser, 1983.\par

\smallskip\noindent[Ta75]\enspace
M.~Takeuchi, On coverings and hyperalgebras of affine algebraic
groups, \textit{Trans. Amer. Math. Soc.} \textbf{211} (1975),
249--275.\par
% 已译源行 4462--4994；第 XXIV 讲正文、附录、参考文献均已覆盖。
